% Limites et continuité — Mathématiques Tle A — Cours signé OnBuch+
% Programme officiel MINESEC. Compiler avec : tectonic MA02-limites-et-continuite-a.tex
\newcommand{\DOCMATIERE}{Mathématiques}
\newcommand{\DOCNIVEAU}{Tle A}
\newcommand{\DOCMODULE}{Relations et opérations fondamentales dans l'ensemble des nombres réels}
\newcommand{\DOCLECON}{A02}
\newcommand{\DOCTITRE}{Limites et continuité}
% OnBuch+ — préambule « cours signés » ENRICHI (compilé avec tectonic / XeLaTeX)
% Système visuel « Pop » d'OnBuch+ : fond crème (jamais blanc), bordures encre
% épaisses, ombres « dures » décalées, titres Archivo Black en capitales.
% Version MATHÉMATIQUES : graphes (pgfplots/tikz), boîtes pédagogiques
% (définition, propriété, méthode, attention, le savais-tu, exercice résolu,
% exercice, TP, bilan), page de garde et signature OnBuch+ ancrée en bas.
%
% Macros attendues dans main.tex AVANT \input{preamble} :
%   \DOCMATIERE  \DOCNIVEAU  \DOCMODULE  \DOCLECON (numéro)  \DOCTITRE
\documentclass[11pt]{article}

\usepackage{fontspec}
\usepackage[french]{babel}
\usepackage[a4paper,margin=2cm,top=3.4cm,bottom=2.8cm,headheight=1.4cm,footskip=1.3cm]{geometry}
\usepackage{amsmath,amssymb,amsfonts}
\usepackage{mathtools} % \xmapsto, \coloneqq... souvent utilisés par les modèles
\usepackage{cancel} % simplifications barrées
% \abs{z} (module, valeur absolue) et \norm{u} (norme), utilisés par les modèles
\providecommand{\abs}{}\let\abs\relax\DeclarePairedDelimiter{\abs}{\lvert}{\rvert}
\providecommand{\norm}{}\let\norm\relax\DeclarePairedDelimiter{\norm}{\lVert}{\rVert}
\usepackage[table]{xcolor} % [table] : fournit aussi \rowcolors (alternance de lignes)
\usepackage{pagecolor}
\usepackage{tikz}
\usetikzlibrary{arrows.meta,positioning,calc,shapes.geometric,shapes.misc,shapes.arrows,decorations.pathmorphing,decorations.markings,patterns,shadows,mindmap,trees,fit,backgrounds,matrix,intersections}
\usepackage{pgfplots}
\usepackage{pgf-pie} % diagrammes circulaires \pie (statistiques)
\pgfplotsset{compat=1.18}
\usepgfplotslibrary{fillbetween,groupplots} % aires entre courbes (calcul intégral)
\usepackage{enumitem}
\usepackage{fancyhdr}
% [most] charge toutes les bibliothèques tcolorbox (listings, minted, raster,
% documentation...) et gonfle inutilement la mémoire TeX sur un document long
% (~300 boîtes) : on ne charge que ce qui est réellement utilisé.
\usepackage[skins,breakable]{tcolorbox}
\usepackage{graphicx}
\usepackage{colortbl}
\usepackage{booktabs}
\usepackage{tabularx}
\usepackage{diagbox} % case d'en-tête diagonale des tableaux à double entrée
% Colonnes extensibles centrée / gauche / droite (C, L, R), souvent utilisées par les modèles
\newcolumntype{C}{>{\centering\arraybackslash}X}
\newcolumntype{L}{>{\raggedright\arraybackslash}X}
\newcolumntype{R}{>{\raggedleft\arraybackslash}X}
\usepackage{array}
\usepackage{multicol}
\usepackage{siunitx}
% parse-numbers=false : accepte aussi \SI{2,5 \times 10^{-3}}{...}
\sisetup{locale=FR,output-decimal-marker={,},group-minimum-digits=4,parse-numbers=false}
\DeclareSIUnit\ohm{\ensuremath{\symup{\Omega}}} % U+2126 absent de la police
\usepackage{qrcode}
% \degree : commande fréquemment utilisée par les modèles pour un angle ou une
% température en dehors de \SI{}{} (non fournie par les paquets chargés).
\providecommand{\degree}{^{\circ}}
% \qed (fin de démonstration), utilisé par les modèles sans amsthm
\providecommand{\qed}{\hfill$\square$}
% \barre : double barre des valeurs interdites dans un tableau de variations (tkz-tab)
\providecommand{\barre}{\ensuremath{\|}}
% environnement proof (amsthm non chargé) utilisé par les modèles
\ifdefined\proof\else\newenvironment{proof}[1][Démonstration]{\par\noindent\textit{#1.}\ }{\qed\par}\fi
\usepackage{lastpage}
\usepackage{hyperref}
\hypersetup{hidelinks}

% ============================================================
% Palette « Pop » OnBuch+ — mêmes hex que la classe P (plus_ui.dart)
% ============================================================
\definecolor{poporange}{HTML}{F59321}
\definecolor{popdark}{HTML}{E07A0C}
\definecolor{popcream}{HTML}{FFF6E8}
\definecolor{popink}{HTML}{1C1714}
\definecolor{poppurple}{HTML}{7A5AE0}
\definecolor{popgreen}{HTML}{2E9E6A}
\definecolor{poppink}{HTML}{E0457B}
\definecolor{popblue}{HTML}{2F7DE1}
\definecolor{popgold}{HTML}{C98A12}
\definecolor{popmuted}{HTML}{8A7F76}
\definecolor{popline}{HTML}{EADFD2}
\definecolor{popgray}{HTML}{8A7F76}
\colorlet{popgrey}{popgray}
% Alias tolérants (le modèle nomme parfois les couleurs différemment)
\colorlet{popwhite}{white}
\colorlet{popblack}{popink}
\colorlet{popred}{poppink}
\colorlet{popyellow}{popgold}
% Teintes claires pour fonds de tableaux / zones
\colorlet{popblueL}{popblue!10!white}
\colorlet{poporangeL}{poporange!14!white}
\colorlet{popgreenL}{popgreen!12!white}
\colorlet{poppurpleL}{poppurple!12!white}
\colorlet{poppinkL}{poppink!12!white}
\colorlet{popgoldL}{popgold!14!white}

\pagecolor{popcream}

% ============================================================
% Typographie
% ============================================================
\defaultfontfeatures{Ligatures=TeX}
\setmainfont{PlusJakartaSans-Medium.ttf}[Path=../../../fonts/,BoldFont=PlusJakartaSans-Bold.ttf]
% Maths en sans-serif (Fira Math) avec chiffres et lettres droites de Plus
% Jakarta, pour que les formules se fondent dans le texte.
\usepackage[math-style=ISO,bold-style=ISO]{unicode-math}
\setmathfont{FiraMath-Regular.otf}[Path=../../../fonts/]
\setmathfont{PlusJakartaSans-Medium.ttf}[Path=../../../fonts/,range={up/num,up/{Latin,latin},\mathup}]
\usepackage{icomma} % virgule décimale sans espace en mode math
% Caractères mathématiques Unicode absents des polices de texte (Plus Jakarta,
% Archivo Black) : rendus par leur équivalent mathématique, en texte comme en
% formule — sinon ils s'affichent en carré vide (ex. « Arithmétique dans ℤ »).
\usepackage{newunicodechar}
\newunicodechar{ℝ}{\ensuremath{\mathbb{R}}}
\newunicodechar{ℤ}{\ensuremath{\mathbb{Z}}}
\newunicodechar{ℂ}{\ensuremath{\mathbb{C}}}
\newunicodechar{ℕ}{\ensuremath{\mathbb{N}}}
\newunicodechar{ℚ}{\ensuremath{\mathbb{Q}}}
\newunicodechar{𝔻}{\ensuremath{\mathbb{D}}}
\newunicodechar{α}{\ensuremath{\alpha}}
\newunicodechar{β}{\ensuremath{\beta}}
\newunicodechar{γ}{\ensuremath{\gamma}}
\newunicodechar{δ}{\ensuremath{\delta}}
\newunicodechar{ε}{\ensuremath{\varepsilon}}
\newunicodechar{θ}{\ensuremath{\theta}}
\newunicodechar{λ}{\ensuremath{\lambda}}
\newunicodechar{μ}{\ensuremath{\mu}}
\newunicodechar{σ}{\ensuremath{\sigma}}
\newunicodechar{τ}{\ensuremath{\tau}}
\newunicodechar{φ}{\ensuremath{\varphi}}
\newunicodechar{ψ}{\ensuremath{\psi}}
\newunicodechar{ω}{\ensuremath{\omega}}
\newunicodechar{Σ}{\ensuremath{\Sigma}}
% majuscules produites par \MakeUppercase dans les titres (α-aminés -> Α)
\newunicodechar{Α}{\ensuremath{\alpha}}
\newunicodechar{Β}{\ensuremath{\beta}}
\newunicodechar{Γ}{\ensuremath{\Gamma}}
\newunicodechar{Δ}{\ensuremath{\Delta}}
\newunicodechar{Ε}{\ensuremath{\varepsilon}}
\newunicodechar{Θ}{\ensuremath{\Theta}}
\newunicodechar{Λ}{\ensuremath{\Lambda}}
\newunicodechar{Μ}{\ensuremath{\mu}}
\newunicodechar{Τ}{\ensuremath{\tau}}
\newunicodechar{Φ}{\ensuremath{\Phi}}
\newunicodechar{Ψ}{\ensuremath{\Psi}}
\newunicodechar{Ω}{\ensuremath{\Omega}}
\newunicodechar{↦}{\ensuremath{\mapsto}}
\newunicodechar{∈}{\ensuremath{\in}}
\newunicodechar{∉}{\ensuremath{\notin}}
\newunicodechar{∧}{\ensuremath{\wedge}}
\newunicodechar{⇔}{\ensuremath{\Leftrightarrow}}
\newunicodechar{⟺}{\ensuremath{\Longleftrightarrow}}
\newunicodechar{⇒}{\ensuremath{\Rightarrow}}
\newunicodechar{⟹}{\ensuremath{\Longrightarrow}}
\newunicodechar{∘}{\ensuremath{\circ}}
\newunicodechar{∪}{\ensuremath{\cup}}
\newunicodechar{∩}{\ensuremath{\cap}}
\newunicodechar{≡}{\ensuremath{\equiv}}
\newunicodechar{⊂}{\ensuremath{\subset}}
\newunicodechar{⊥}{\ensuremath{\perp}}
\newunicodechar{∃}{\ensuremath{\exists}}
\newunicodechar{∀}{\ensuremath{\forall}}
\newunicodechar{∛}{\ensuremath{\sqrt[3]{\,}}}
\newunicodechar{∜}{\ensuremath{\sqrt[4]{\,}}}
\newunicodechar{⃗}{\ensuremath{{}^{\to}}}
\newunicodechar{ⁿ}{\ifmmode^{n}\else\textsuperscript{\ensuremath{n}}\fi}
\newunicodechar{ˣ}{\ifmmode^{x}\else\textsuperscript{\ensuremath{x}}\fi}
\newunicodechar{ᵇ}{\ifmmode^{b}\else\textsuperscript{\ensuremath{b}}\fi}
\newunicodechar{ʳ}{\ifmmode^{r}\else\textsuperscript{\ensuremath{r}}\fi}
\newunicodechar{ᵘ}{\ifmmode^{u}\else\textsuperscript{\ensuremath{u}}\fi}
\newunicodechar{ʸ}{\ifmmode^{y}\else\textsuperscript{\ensuremath{y}}\fi}
\newunicodechar{ᵃ}{\ifmmode^{a}\else\textsuperscript{\ensuremath{a}}\fi}
\newunicodechar{ᵉ}{\ifmmode^{e}\else\textsuperscript{\ensuremath{e}}\fi}
\newunicodechar{ᵏ}{\ifmmode^{k}\else\textsuperscript{\ensuremath{k}}\fi}
\newunicodechar{ᵖ}{\ifmmode^{p}\else\textsuperscript{\ensuremath{p}}\fi}
\newunicodechar{ᴺ}{\ifmmode^{N}\else\textsuperscript{\ensuremath{N}}\fi}
\newunicodechar{ᶜ}{\ifmmode^{c}\else\textsuperscript{\ensuremath{c}}\fi}
\newunicodechar{ʰ}{\ifmmode^{h}\else\textsuperscript{\ensuremath{h}}\fi}
\newunicodechar{ᵠ}{\ifmmode^{q}\else\textsuperscript{\ensuremath{q}}\fi}
\newunicodechar{ₙ}{\ifmmode_{n}\else\textsubscript{\ensuremath{n}}\fi}
\newunicodechar{ₐ}{\ifmmode_{a}\else\textsubscript{\ensuremath{a}}\fi}
\newunicodechar{ᵢ}{\ifmmode_{i}\else\textsubscript{\ensuremath{i}}\fi}
\newunicodechar{ᵣ}{\ifmmode_{r}\else\textsubscript{\ensuremath{r}}\fi}
\newunicodechar{ₖ}{\ifmmode_{k}\else\textsubscript{\ensuremath{k}}\fi}
\newunicodechar{ᵦ}{\ifmmode_{\beta}\else\textsubscript{\ensuremath{\beta}}\fi}
\newunicodechar{ₓ}{\ifmmode_{x}\else\textsubscript{\ensuremath{x}}\fi}
\newfontfamily\popdisplay{ArchivoBlack-Regular.ttf}[Path=../../../fonts/]
\newfontfamily\poplogo{SpaceGrotesk-Bold.ttf}[Path=../../../fonts/]
\newfontfamily\popsemi{PlusJakartaSans-SemiBold.ttf}[Path=../../../fonts/]
\newfontfamily\popxbold{PlusJakartaSans-ExtraBold.ttf}[Path=../../../fonts/]
\newfontfamily\popxboldls{PlusJakartaSans-ExtraBold.ttf}[Path=../../../fonts/,LetterSpace=3.0]

\setlist[enumerate]{leftmargin=1.7em,itemsep=5pt,topsep=4pt}
\setlist[itemize]{leftmargin=1.7em,itemsep=4pt,topsep=4pt,label={\color{poporange}\textbullet}}
\setlength{\parindent}{0pt}
\setlength{\parskip}{5pt}
\renewcommand{\arraystretch}{1.25}

% pgfplots : style Pop par défaut
\pgfplotsset{
  every axis/.append style={axis line style={popink,line width=1pt},tick style={popink},
    grid style={popline,line width=0.5pt},label style={font=\small},tick label style={font=\footnotesize}},
  popcurve/.style={poporange,line width=1.8pt,no markers,smooth},
}
% Même style disponible dans un \draw TikZ simple (hors pgfplots)
\tikzset{popcurve/.style={poporange,line width=1.8pt}}

% ============================================================
% Pill / squiggle
% ============================================================
\newcommand{\poppill}[3]{%
  \tikz[baseline=-0.55ex]\node[draw=popink,line width=1.2pt,rounded corners=2.6pt,%
    fill=#2,inner xsep=6.5pt,inner ysep=3pt]{%
    \popxboldls\fontsize{7}{7}\selectfont\color{#3}\MakeUppercase{#1}};%
}
\newcommand{\popsquiggle}[1]{%
  \tikz[baseline=-0.4ex]{
    \draw[#1,line width=2.6pt,line cap=round]
      plot [smooth] coordinates {(0,0.09) (0.34,0.01) (0.68,0.21) (1.02,0.01) (1.36,0.10)};
  }%
}
% Mot-clé surligné (marqueur orange sous le texte)
\newcommand{\cle}[1]{{\popxbold\color{popdark}#1}}

% ============================================================
% En-tête / pied
% ============================================================
\pagestyle{fancy}
\fancyhf{}
\renewcommand{\headrulewidth}{0pt}
\renewcommand{\footrulewidth}{0pt}
\lhead{\raisebox{-0.15ex}{\poplogo\fontsize{13}{13}\selectfont\color{popink}OnBuch\color{poporange}+}}
\rhead{\poppill{\DOCMATIERE\ \textbullet\ \DOCNIVEAU\ \textbullet\ Leçon \DOCLECON}{white}{popink}}
\lfoot{\popxbold\fontsize{7.6}{7.6}\selectfont\color{popmuted}\MakeUppercase{Cours signé OnBuch+}\ \textbullet\ {\popsemi\DOCTITRE}}
\rfoot{\tikz[baseline=-0.6ex]\node[rounded corners=8.5pt,fill=popink,minimum height=17pt,minimum width=17pt,inner xsep=5pt,inner sep=0]{\popxbold\fontsize{8}{8}\selectfont\color{popcream}\thepage\,/\,\pageref*{LastPage}};}

% ============================================================
% Titres
% ============================================================
\newcommand{\chaptitre}[2]{%
  \begin{tcolorbox}[enhanced,colback=poporange,colframe=popink,boxrule=2.6pt,arc=11pt,
      drop shadow=popink,left=15pt,right=15pt,top=12pt,bottom=11pt,before skip=3pt,after skip=18pt]
    \poppill{#2}{popink}{popcream}\par\vspace{9pt}
    {\raggedright\hyphenpenalty=10000\exhyphenpenalty=10000\popdisplay\fontsize{22}{24}\selectfont\color{popcream}\MakeUppercase{#1}\par}
  \end{tcolorbox}
}
% Section numérotée : pastille orange avec le numéro + titre Archivo
\newcounter{popsec}
\newcommand{\coursec}[1]{%
  \stepcounter{popsec}\setcounter{popsub}{0}%
  \par\vspace{16pt}\noindent
  \begin{minipage}{\linewidth}
  \tikz[baseline=-0.7ex]\node[draw=popink,line width=1.6pt,fill=poporange,rounded corners=4pt,minimum size=22pt,inner sep=0]{\popdisplay\fontsize{12}{12}\selectfont\color{popcream}\thepopsec};%
  \hspace{8pt}{\popdisplay\fontsize{15}{17}\selectfont\color{popink}\MakeUppercase{#1}}\par
  \vspace{4pt}\noindent{\color{poporange}\rule{36pt}{3.6pt}}
  \end{minipage}\par\nopagebreak\vspace{8pt}
}
\newcounter{popsub}
\newcommand{\courssub}[1]{%
  \stepcounter{popsub}%
  \par\vspace{9pt}\noindent{\popxbold\fontsize{11.5}{13}\selectfont\color{popink}{\color{poporange}\thepopsec.\thepopsub}\hspace{6pt}#1}\par\nopagebreak\vspace{4pt}
}

% ============================================================
% Boîtes pédagogiques — même recette : fond blanc, bordure épaisse colorée,
% ombre dure encre, badge plein. Titre optionnel [#1].
% ============================================================
\tcbset{popbase/.style={breakable,enhanced,colback=white,boxrule=2.3pt,arc=9pt,
  shadow={3pt}{-3pt}{0pt}{popink},left=14pt,right=14pt,top=11pt,bottom=11pt,before skip=11pt,after skip=15pt,
  fontupper=\popsemi\fontsize{10.6}{14.5}\selectfont\color{popink},
  fontlower=\popsemi\fontsize{10.6}{14.5}\selectfont\color{popink}}}
% Coche dessinée (le glyphe ✓ n'existe pas dans les polices du document)
\newcommand{\popcheck}[1]{\tikz[baseline=-0.5ex]\draw[#1,line width=1.6pt,line cap=round,line join=round](-0.13,0.0)--(-0.04,-0.09)--(0.13,0.1);}
\newcommand{\popboxhead}[3]{\poppill{#1}{#2}{white}\ifx\relax#3\relax\else\hspace{6pt}{\popxbold\fontsize{10.6}{12}\selectfont\color{popink}#3}\fi\par\vspace{7pt}}

\newtcolorbox{aretenir}[1][]{popbase,colframe=popgreen,before upper={\popboxhead{À retenir}{popgreen}{#1}}}
\newtcolorbox{exemplebox}[1][]{popbase,colframe=poporange,before upper={\popboxhead{Exemple}{poporange}{#1}}}
\newtcolorbox{definition}[1][]{popbase,colframe=popblue,before upper={\popboxhead{Définition}{popblue}{#1}}}
\newtcolorbox{propriete}[1][]{popbase,colframe=poppurple,before upper={\popboxhead{Propriété}{poppurple}{#1}}}
\newtcolorbox{methode}[1][]{popbase,colframe=popgold,before upper={\popboxhead{Méthode}{popgold}{#1}}}
\newtcolorbox{attention}[1][]{popbase,colframe=poppink,before upper={\popboxhead{Attention}{poppink}{#1}}}
\newtcolorbox{experience}[1][]{popbase,colframe=popblue,colback=popblueL,before upper={\popboxhead{Expérience}{popblue}{#1}}}
% « Le savais-tu ? » : carte violette pleine (contexte, culture, Cameroun)
\newtcolorbox{savaistu}[1][]{popbase,colback=poppurple,colframe=popink,
  fontupper=\popsemi\fontsize{10.4}{14.2}\selectfont\color{white},
  before upper={\poppill{Le savais-tu\,?}{white}{poppurple}\ifx\relax#1\relax\else\hspace{6pt}{\popxbold\color{white}#1}\fi\par\vspace{7pt}}}
% Exercice résolu : énoncé en haut, \tcblower puis la solution sur fond crème
\newcounter{popexo}\newcounter{popexr}
\newtcolorbox{exoresolu}[1][]{popbase,colframe=poporange,colbacklower=popcream,
  segmentation style={popink,line width=1.2pt,dashed},
  before upper={\stepcounter{popexr}\popboxhead{Exercice résolu \thepopexr}{poporange}{#1}},
  before lower={\poppill{Solution}{popink}{popcream}\par\vspace{6pt}}}
% Exercice (non résolu en place) — la correction est dans la partie Corrigés
\newtcolorbox{exercice}[1][]{popbase,colframe=popink,
  before upper={\stepcounter{popexo}\popboxhead{Exercice \thepopexo}{popink}{#1}}}
% Corrigé
\newtcolorbox{corrige}[1][]{popbase,colframe=popgreen,colback=popgreenL,
  before upper={\popboxhead{Corrigé}{popgreen}{#1}}}
% Objectifs / prérequis / bilan
\newtcolorbox{objectifs}{popbase,colframe=popink,colback=poporangeL,
  before upper={\popboxhead{Objectifs de la leçon}{popink}{}}}
\newtcolorbox{prerequis}{popbase,colframe=popink,colback=popblueL,
  before upper={\popboxhead{Prérequis}{popblue}{}}}
\newtcolorbox{bilan}{popbase,colframe=popink,colback=white,boxrule=2.8pt,
  before upper={\popboxhead{Fiche bilan}{poporange}{L'essentiel à connaître pour l'examen}}}
% Figure encadrée avec légende
\newtcolorbox{popfigure}[1][]{enhanced,breakable=false,colback=white,colframe=popline,boxrule=1.4pt,arc=8pt,
  left=10pt,right=10pt,top=10pt,bottom=8pt,before skip=10pt,after skip=12pt,halign=center,
  lower separated=false,halign lower=center,
  fontlower=\popsemi\fontsize{9}{11.5}\selectfont\color{popmuted},
  #1}
\newcommand{\legende}[1]{\par\vspace{6pt}{\popsemi\fontsize{9}{11.5}\selectfont\color{popmuted}#1}}

% ============================================================
% Signature OnBuch+
% ============================================================
\newcommand{\poplogoinline}[1]{{\poplogo\fontsize{#1}{#1}\selectfont\color{popink}OnBuch\color{poporange}+}}
\newcommand{\signatureonbuch}{%
  \begin{tcolorbox}[enhanced,colback=popink,colframe=popink,boxrule=2.2pt,arc=10pt,
      drop shadow=poporange,left=14pt,right=14pt,top=10pt,bottom=10pt,before skip=0pt,after skip=0pt]
    \tikz[baseline=-0.7ex]\node[circle,fill=popgreen,draw=popcream,line width=1.3pt,minimum size=22pt,inner sep=0]{\popcheck{white}};%
    \hspace{9pt}\begin{minipage}[c]{0.62\linewidth}
      {\popxbold\fontsize{12}{13}\selectfont\color{popcream}Signé\ \textbullet\ {\poplogo OnBuch\color{poporange}+}}\par\vspace{2pt}
      {\raggedright\popsemi\fontsize{8}{10.5}\selectfont\color{popline}\MakeUppercase{Équipe pédagogique OnBuch+}\\\MakeUppercase{Conforme au programme officiel MINESEC}\par}
    \end{minipage}\hfill
    {\poplogo\fontsize{22}{22}\selectfont\color{popcream}OnBuch\color{poporange}+}
  \end{tcolorbox}%
}

% ============================================================
% Page de garde
% #1 titre · #2 matière · #3 niveau · #4 module · #5 n° leçon · #6 durée ·
% #7 description / objectifs (texte ou itemize)
% ============================================================
\newcommand{\pagegarde}[7]{%
  \thispagestyle{empty}
  \newgeometry{a4paper,margin=1.7cm,top=1.6cm,bottom=1.4cm}%
  \begin{tikzpicture}[remember picture,overlay]
    \fill[poporange!18] ([xshift=-3.2cm,yshift=-3.6cm]current page.north east) circle (4.6cm);
    \fill[poppurple!14] ([xshift=2.4cm,yshift=7.6cm]current page.south west) circle (3.8cm);
  \end{tikzpicture}%
  {\poplogo\fontsize{30}{30}\selectfont\color{popink}OnBuch\color{poporange}+}\hfill
  \raisebox{0.6ex}{\poppill{Cours signé \textbullet\ Premium}{popink}{popcream}}\par
  \vspace{1pt}\popsquiggle{poppurple}\par
  \vspace{8pt}
  \begin{tcolorbox}[enhanced,colback=poporange,colframe=popink,boxrule=2.8pt,arc=13pt,
      drop shadow=popink,left=18pt,right=18pt,top=12pt,bottom=13pt,after skip=10pt]
    \poppill{#2 \textbullet\ #3}{popink}{popcream}\hspace{5pt}\poppill{#4}{white}{popink}\par\vspace{8pt}
    {\popdisplay\fontsize{14}{14}\selectfont\color{popink}LEÇON #5}\par\vspace{4pt}
    {\raggedright\hyphenpenalty=10000\exhyphenpenalty=10000\popdisplay\fontsize{25}{27}\selectfont\color{popcream}\MakeUppercase{#1}\par}
    \vspace{8pt}
    {\popxbold\fontsize{9.5}{9.5}\selectfont\color{popink}\MakeUppercase{Durée conseillée : #6}}
  \end{tcolorbox}
  \begin{tcolorbox}[enhanced,colback=white,colframe=popink,boxrule=2.2pt,arc=10pt,
      drop shadow=popink,left=16pt,right=16pt,top=10pt,bottom=10pt,after skip=8pt]
    \poppill{Dans cette leçon}{poppurple}{white}\par\vspace{5pt}
    \popsemi\fontsize{9.3}{12.6}\selectfont\color{popink}#7
  \end{tcolorbox}
  \noindent\begin{minipage}[t]{0.487\linewidth}
    \begin{tcolorbox}[enhanced,colback=popgreen,colframe=popink,boxrule=2.2pt,arc=10pt,
        drop shadow=popink,left=11pt,right=11pt,top=10pt,bottom=10pt,height=3.1cm,valign=center]
      \tcbox[colback=white,colframe=popink,boxrule=1.3pt,arc=5pt,boxsep=0pt,nobeforeafter,
        left=3pt,right=3pt,top=3pt,bottom=3pt,valign=center]{\qrcode[height=1.9cm]{https://play.google.com/store/apps/details?id=com.luvvix.onbuch&hl=fr}}%
      \hspace{8pt}\begin{minipage}[c]{0.5\linewidth}
        {\popdisplay\fontsize{11}{12.5}\selectfont\color{white}\MakeUppercase{Télécharge OnBuch}}\par\vspace{4pt}
        {\raggedright\popsemi\fontsize{8.4}{11}\selectfont\color{white}Gratuit \textbullet\ Google Play\\Tuteur IA, annales, concours, résultats\par}
      \end{minipage}
    \end{tcolorbox}
  \end{minipage}\hfill
  \begin{minipage}[t]{0.487\linewidth}
    \begin{tcolorbox}[enhanced,colback=poppurple,colframe=popink,boxrule=2.2pt,arc=10pt,
        drop shadow=popink,left=11pt,right=11pt,top=10pt,bottom=10pt,height=3.1cm,valign=center]
      \tikz[baseline=-0.7ex]\node[circle,fill=white,draw=popink,line width=1.3pt,minimum size=1.6cm,inner sep=0]{\popdisplay\fontsize{24}{24}\selectfont\color{poppurple}+};%
      \hspace{8pt}\begin{minipage}[c]{0.6\linewidth}
        {\popdisplay\fontsize{11}{12.5}\selectfont\color{white}\MakeUppercase{Passe à OnBuch+}}\par\vspace{4pt}
        {\raggedright\popsemi\fontsize{8.4}{11}\selectfont\color{white}Tous les cours signés, Tuteur IA illimité, fiches PDF — onglet OnBuch+ de l'app\par}
      \end{minipage}
    \end{tcolorbox}
  \end{minipage}
  \vspace{8pt}
  \popcontenu
  \vfill
  \signatureonbuch
  \restoregeometry
  \newpage
}

% Rangée « Ce cours contient » (page de garde)
\newcommand{\poptile}[3]{% #1 couleur #2 gros texte #3 légende
  \begin{tcolorbox}[enhanced,colback=#1,colframe=popink,boxrule=2pt,arc=9pt,shadow={2.5pt}{-2.5pt}{0pt}{popink},
      left=6pt,right=6pt,top=9pt,bottom=9pt,halign=center,valign=center,height=1.75cm,width=\linewidth,nobeforeafter]
    {\popdisplay\fontsize{12.5}{13}\selectfont\color{white}\MakeUppercase{#2}}\par\vspace{3pt}
    {\centering\popsemi\fontsize{7.8}{9.5}\selectfont\color{white}#3\par}
  \end{tcolorbox}}
\newcommand{\popcontenu}{%
  \par\noindent{\popxboldls\fontsize{7.5}{7.5}\selectfont\color{popmuted}CE COURS SIGNÉ CONTIENT}\par\vspace{7pt}
  \noindent\begin{minipage}[t]{0.235\linewidth}\poptile{popblue}{Cours}{complet, illustré, pas à pas}\end{minipage}\hfill
  \begin{minipage}[t]{0.235\linewidth}\poptile{poporange}{Méthodes}{et exercices résolus}\end{minipage}\hfill
  \begin{minipage}[t]{0.235\linewidth}\poptile{poppink}{Exercices}{type examen + corrigés}\end{minipage}\hfill
  \begin{minipage}[t]{0.235\linewidth}\poptile{popgreen}{Fiche bilan}{l'essentiel à retenir}\end{minipage}\par}

% Sommaire
\newtcolorbox{sommaire}{enhanced,colback=white,colframe=popink,boxrule=2.2pt,arc=10pt,
  drop shadow=popink,left=18pt,right=18pt,top=13pt,bottom=13pt,before skip=6pt,after skip=20pt,
  before upper={\poppill{Sommaire}{poppurple}{white}\par\vspace{9pt}\popsemi\fontsize{10.6}{14.5}\selectfont\color{popink}}}

% Signature de fin de cours — poussée tout en bas de la dernière page
\newcommand{\findecours}{%
  \par\vfill\nopagebreak
  \begin{center}\popsquiggle{poporange}\end{center}\vspace{4pt}
  \signatureonbuch
}
\AtBeginDocument{\def\llbracket{\mathopen{[\mkern-3.5mu[}}\def\rrbracket{\mathclose{]\mkern-3.5mu]}}} % intervalles d'entiers (glyphe ⟦ absent de la police)
% Glyphes absents de Fira Math : redéfinis à partir de glyphes disponibles
\AtBeginDocument{%
  \def\setminus{\mathbin{\backslash}}\let\smallsetminus\setminus
  \def\vdots{\mathord{\vbox{\baselineskip3.2pt\lineskiplimit0pt\kern4pt\hbox{.}\hbox{.}\hbox{.}}}}%
  \def\ddots{\mathinner{\mkern1mu\raise6.5pt\hbox{.}\mkern2mu\raise3.6pt\hbox{.}\mkern2mu\raise0.7pt\hbox{.}\mkern1mu}}%
  \def\triangle{\mathord{\scalebox{0.9}{$\symup{\Delta}$}}}%
  \def\nabla{\mathord{\raisebox{\depth}{\scalebox{1}[-1]{$\symup{\Delta}$}}}}%
  \def\checkmark{\popcheck{popdark}}%
  \def\star{\mathbin{\ast}}\def\bigstar{\mathord{\ast}}%
  \def\diamond{\mathbin{\scalebox{0.6}{$\Diamond$}}}%
}

%% FIGFIX-BEGIN v5 — correctifs globaux des figures (docs/onbuch-generation/tools/figfix)
\usepackage{adjustbox}\usepackage{etoolbox}\tcbuselibrary{raster}
% toute figure TikZ/pgfplots plus large que la ligne est réduite à la largeur disponible
\BeforeBeginEnvironment{tikzpicture}{\begin{adjustbox}{max width=\linewidth}}
\AfterEndEnvironment{tikzpicture}{\end{adjustbox}}
% légendes pgfplots lisibles par-dessus les courbes
\pgfplotsset{every axis/.append style={legend style={fill=white,fill opacity=0.92,text opacity=1,draw=black!15,inner sep=3pt}}}
% étiquettes d'arêtes (syntaxe edge["..."]) avec fond blanc
\tikzset{every edge quotes/.append style={fill=white,inner sep=1.2pt}}
%% FIGFIX-END
\begin{document}

\pagegarde{Limites et continuité}{Mathématiques}{Tle A}{Relations et opérations fondamentales dans l'ensemble des nombres réels}{A02}{5 h}{Cette leçon introduit les notions fondamentales de l'analyse : la limite d'une fonction en un point et à l'infini, ainsi que la continuité. À travers les limites usuelles, les opérations algébriques et les techniques de levée d'indétermination, l'élève apprend à décrire le comportement local et asymptotique des fonctions et à interpréter graphiquement les résultats.
\vspace{4pt}
\begin{multicols}{2}\small
\begin{itemize}
\item À la fin de cette leçon, je sais définir et reconnaître la limite d'une fonction en un point, à droite, à gauche et en l'infini
\item À la fin de cette leçon, je sais citer et utiliser les limites usuelles (polynômes, rationnelles, racine carrée, valeur absolue)
\item À la fin de cette leçon, je sais calculer une limite en appliquant les théorèmes sur les sommes, produits, quotients et composées
\item À la fin de cette leçon, je sais lever une forme indéterminée simple par factorisation ou par la quantité conjuguée
\item À la fin de cette leçon, je sais déterminer la limite d'une fonction polynomiale ou rationnelle en l'infini par la règle du terme de plus haut degré
\item À la fin de cette leçon, je sais vérifier la continuité d'une fonction en un point et sur un intervalle
\end{itemize}\end{multicols}}

\begin{sommaire}
\begin{multicols}{2}
\begin{enumerate}
\item Limite d'une fonction en un point
\item Limite d'une fonction en l'infini
\item Opérations sur les limites
\item Lever une forme indéterminée
\item Continuité en un point et sur un intervalle
\item Interprétations graphiques et synthèse
\item Activité d'intégration
\item Méthodes et exercices résolus
\item Exercices
\item Corrigés des exercices
\item Fiche bilan
\end{enumerate}
\end{multicols}
\end{sommaire}

\begin{objectifs}
\begin{itemize}
\item À la fin de cette leçon, je sais définir et reconnaître la limite d'une fonction en un point, à droite, à gauche et en l'infini
\item À la fin de cette leçon, je sais citer et utiliser les limites usuelles (polynômes, rationnelles, racine carrée, valeur absolue)
\item À la fin de cette leçon, je sais calculer une limite en appliquant les théorèmes sur les sommes, produits, quotients et composées
\item À la fin de cette leçon, je sais lever une forme indéterminée simple par factorisation ou par la quantité conjuguée
\item À la fin de cette leçon, je sais déterminer la limite d'une fonction polynomiale ou rationnelle en l'infini par la règle du terme de plus haut degré
\item À la fin de cette leçon, je sais vérifier la continuité d'une fonction en un point et sur un intervalle
\item À la fin de cette leçon, je sais interpréter graphiquement une limite (asymptote, comportement) ou une discontinuité (saut, trou)
\item À la fin de cette leçon, je sais rédiger rigoureusement un calcul de limite en justifiant chaque étape
\end{itemize}
\end{objectifs}

\begin{prerequis}
\begin{itemize}
\item Fonctions usuelles de Première : polynômes, fonctions rationnelles, racine carrée, valeur absolue
\item Lecture graphique d'une fonction (image, courbe représentative)
\item Factorisation d'expressions algébriques et identités remarquables
\item Ensemble de définition d'une fonction
\item Calcul algébrique sur les fractions et les radicaux (quantité conjuguée)
\end{itemize}
\end{prerequis}

\newpage

\coursec{Limite d'une fonction en un point}

Au marché Mokolo, Aminatou vend ses beignets chaque matin. En affinant ses comptes, elle remarque quelque chose de curieux : plus elle produit de beignets, plus le coût unitaire de fabrication se rapproche de 25 FCFA, sans jamais l'atteindre exactement. Un autre jour, son fournisseur de farine change brusquement ses tarifs : dès qu'elle dépasse 200 beignets, le coût unitaire \emph{saute} d'une valeur à une autre, sans transition. Deux phénomènes très différents : dans le premier cas, une quantité \emph{s'approche} d'une valeur ; dans le second, elle \emph{saute} brutalement. Comment décrire mathématiquement ces comportements ? C'est toute la problématique de cette leçon : la notion de \cle{limite} permet de parler rigoureusement d'une valeur « dont on s'approche », et la distinction entre limite à droite et limite à gauche permettra de comprendre les « sauts ». Ces outils sont le socle de toute l'analyse du programme de Terminale.

\courssub{Approche intuitive : s'approcher d'une valeur}

Considérons la fonction $f$ définie pour tout réel $x \neq 1$ par
\[ f(x) = \frac{x^2 - 1}{x - 1}. \]
Cette fonction n'est \emph{pas définie} en $x_0 = 1$ : le dénominateur s'y annule. Pourtant, on peut se poser la question : lorsque $x$ se rapproche de $1$ (sans jamais valoir $1$), que devient $f(x)$ ? Calculons quelques valeurs.

\begin{center}
\begin{tabularx}{\linewidth}{|l|X|X|X|X|X|X|}
\hline
\rowcolor{popblueL}
$x$ & 0,9 & 0,99 & 0,999 & 1,001 & 1,01 & 1,1 \\
\hline
$f(x)$ & 1,9 & 1,99 & 1,999 & 2,001 & 2,01 & 2,1 \\
\hline
\end{tabularx}
\end{center}

On observe que plus $x$ se rapproche de $1$, plus $f(x)$ se rapproche de $2$, et ceci \emph{des deux côtés} de $1$. Ce n'est pas un hasard : pour $x \neq 1$, on peut factoriser le numérateur :
\[ f(x) = \frac{(x-1)(x+1)}{x-1} = x + 1, \]
et il est clair que $x+1$ s'approche de $2$ quand $x$ s'approche de $1$. On dit que $f$ admet pour \cle{limite} $2$ en $1$, et l'on écrit :
\[ \lim_{x \to 1} f(x) = 2. \]

\begin{popfigure}
\begin{center}
\begin{tikzpicture}
\begin{axis}[
    width=0.85\linewidth, height=6.5cm,
    axis lines=middle,
    xlabel={$x$}, ylabel={$y$},
    xmin=-1.5, xmax=3.5, ymin=-0.5, ymax=4.5,
    xtick={1}, ytick={2},
    grid=both, grid style={popline!40},
    domain=-1.5:3.5, samples=100
]
\addplot[popcurve] {x+1};
\addplot[only marks, mark=*, mark size=2pt, poporange, fill=white] coordinates {(1,2)};
\draw[dashed, poporange] (axis cs:1,0) -- (axis cs:1,2) -- (axis cs:0,2);
\node[poporange, above right] at (axis cs:1,2) {trou};
\end{axis}
\end{tikzpicture}
\end{center}
\legende{La fonction $f(x)=\dfrac{x^2-1}{x-1}$ coïncide avec la droite $y=x+1$, sauf en $x=1$ où elle présente un « trou ». La limite en $1$ vaut $2$, même si $f(1)$ n'existe pas.}
\end{popfigure}

L'idée à retenir est simple : la limite décrit le \emph{comportement} de $f(x)$ au voisinage de $x_0$, et non la valeur de la fonction en $x_0$.

\courssub{Définition de la limite finie en un point}

\begin{definition}[Limite finie en un point]
Soit $f$ une fonction définie sur un intervalle contenant $x_0$, sauf éventuellement en $x_0$ lui-même, et soit $L$ un nombre réel. On dit que $f$ admet pour \cle{limite} $L$ en $x_0$, et l'on note
\[ \lim_{x \to x_0} f(x) = L, \]
lorsque $f(x)$ peut être rendu aussi proche de $L$ que l'on veut, pourvu que $x$ soit suffisamment proche de $x_0$ (avec $x \neq x_0$).
\end{definition}

Autrement dit : on peut imposer n'importe quelle marge d'erreur autour de $L$ (par exemple $10^{-3}$ ou $10^{-6}$), et l'on trouvera toujours une zone autour de $x_0$ dans laquelle $f(x)$ respecte cette marge. C'est exactement ce qui se passe pour Aminatou : en produisant suffisamment de beignets, le coût unitaire s'approche de 25 FCFA à moins d'un franc près, puis à moins d'un centime près.

\begin{attention}[Limite en $x_0$ et valeur $f(x_0)$ : deux notions distinctes]
L'erreur la plus fréquente consiste à confondre $\displaystyle\lim_{x \to x_0} f(x)$ et $f(x_0)$. Retiens bien :
\begin{itemize}
\item la limite en $x_0$ peut exister \textbf{sans que $f(x_0)$ existe} : c'est le cas de $f(x)=\dfrac{x^2-1}{x-1}$ en $x_0=1$, dont la limite est $2$ alors que $f(1)$ n'est pas défini ;
\item même si $f(x_0)$ existe, la limite peut être \textbf{différente} de $f(x_0)$ (cas des fonctions définies par morceaux, voir plus bas) ;
\item pour calculer une limite en $x_0$, on ne remplace \textbf{jamais} directement $x$ par $x_0$ sans réfléchir : on étudie le comportement au voisinage.
\end{itemize}
\end{attention}

\begin{exemplebox}[Limite par factorisation directe]
Calculons $\displaystyle\lim_{x \to 2} \frac{x^2 - 4}{x - 2}$. La fonction n'est pas définie en $2$ (numérateur et dénominateur s'annulent), mais pour $x \neq 2$ :
\[ \frac{x^2 - 4}{x - 2} = \frac{(x-2)(x+2)}{x-2} = x + 2. \]
Quand $x$ s'approche de $2$, $x+2$ s'approche de $4$. Donc
\[ \lim_{x \to 2} \frac{x^2 - 4}{x - 2} = 4. \]
\end{exemplebox}

\courssub{Limite à droite et limite à gauche}

Il arrive qu'une fonction ne se comporte pas de la même façon selon que l'on s'approche de $x_0$ par valeurs supérieures (par la \emph{droite}) ou par valeurs inférieures (par la \emph{gauche}). C'est précisément le « saut » observé par Aminatou lorsque son fournisseur change de tarif au seuil de 200 beignets.

\begin{definition}[Limites latérales]
Soit $f$ une fonction définie au voisinage de $x_0$, sauf éventuellement en $x_0$.
\begin{itemize}
\item On dit que $f$ admet pour \cle{limite à droite} $L$ en $x_0$, et l'on note
\[ \lim_{\substack{x \to x_0 \\ x > x_0}} f(x) = L \quad \text{ou} \quad \lim_{x \to x_0^{+}} f(x) = L, \]
lorsque $f(x)$ s'approche de $L$ quand $x$ s'approche de $x_0$ par valeurs \textbf{strictement supérieures}.
\item On dit que $f$ admet pour \cle{limite à gauche} $L'$ en $x_0$, et l'on note
\[ \lim_{\substack{x \to x_0 \\ x < x_0}} f(x) = L' \quad \text{ou} \quad \lim_{x \to x_0^{-}} f(x) = L', \]
lorsque $f(x)$ s'approche de $L'$ quand $x$ s'approche de $x_0$ par valeurs \textbf{strictement inférieures}.
\end{itemize}
\end{definition}

\begin{exemplebox}[Une fonction en escalier]
Soit $g$ la fonction définie sur $\mathbb{R}$ par
\[ g(x) = \begin{cases} 1 & \text{si } x < 2 \\ 3 & \text{si } x \geq 2 \end{cases} \]
Quand $x$ s'approche de $2$ par la gauche, $g(x)$ vaut constamment $1$ : donc $\displaystyle\lim_{x \to 2^{-}} g(x) = 1$. Quand $x$ s'approche de $2$ par la droite, $g(x)$ vaut constamment $3$ : donc $\displaystyle\lim_{x \to 2^{+}} g(x) = 3$. Les deux limites latérales existent mais sont \emph{différentes} : la fonction « saute » en $2$, comme le coût unitaire d'Aminatou au seuil des 200 beignets.
\end{exemplebox}

\begin{popfigure}
\begin{center}
\begin{tikzpicture}
\begin{axis}[
    width=0.85\linewidth, height=6cm,
    axis lines=middle,
    xlabel={$x$}, ylabel={$y$},
    xmin=-0.5, xmax=4.5, ymin=-0.5, ymax=4.5,
    xtick={2}, ytick={1,3},
    grid=both, grid style={popline!40}
]
\addplot[popcurve, domain=-0.5:2, samples=2] {1};
\addplot[popcurve, domain=2:4.5, samples=2] {3};
\addplot[only marks, mark=*, mark size=2pt, popblue, fill=white] coordinates {(2,1)};
\addplot[only marks, mark=*, mark size=2pt, popblue, fill=popblue] coordinates {(2,3)};
\node[popblue, below left] at (axis cs:2,1) {$\lim_{x\to 2^-} g = 1$};
\node[popblue, above right] at (axis cs:2,3) {$\lim_{x\to 2^+} g = 3$};
\end{axis}
\end{tikzpicture}
\end{center}
\legende{Fonction en escalier : les limites à gauche ($1$) et à droite ($3$) en $2$ sont différentes. Le point plein marque la valeur $g(2)=3$, le point creux la valeur non atteinte.}
\end{popfigure}

Le lien entre limite globale et limites latérales est résumé par le théorème suivant, que nous admettons et justifions graphiquement : pour que $f(x)$ s'approche d'une \emph{même} valeur $L$ quand $x$ s'approche de $x_0$, il faut et il suffit que ce soit le cas des deux côtés à la fois.

\begin{propriete}[Existence de la limite et limites latérales — admis]
Soit $f$ une fonction définie au voisinage de $x_0$, sauf éventuellement en $x_0$, et $L$ un réel. Alors :
\[ \lim_{x \to x_0} f(x) = L \quad \Longleftrightarrow \quad \lim_{x \to x_0^{-}} f(x) = L \ \text{et} \ \lim_{x \to x_0^{+}} f(x) = L. \]
En particulier, si les limites à gauche et à droite sont différentes (ou si l'une d'elles n'existe pas), alors $f$ \textbf{n'admet pas de limite} en $x_0$.
\end{propriete}

Sur la figure de la fonction en escalier ci-dessus, on lit immédiatement que les limites latérales en $2$ valent $1$ et $3$ : elles diffèrent, donc $g$ n'a pas de limite en $2$. En revanche, pour $f(x)=\dfrac{x^2-1}{x-1}$, les deux limites latérales en $1$ valent $2$, donc la limite existe et vaut $2$.

\begin{methode}[Étudier la limite en un point d'une fonction définie par morceaux]
Pour étudier la limite en un point de jonction $x_0$ d'une fonction définie par morceaux :
\begin{enumerate}
\item Identifier l'expression de $f$ valable pour $x < x_0$ et celle valable pour $x > x_0$.
\item Calculer la \textbf{limite à gauche} $\displaystyle\lim_{x \to x_0^{-}} f(x)$ avec la première expression.
\item Calculer la \textbf{limite à droite} $\displaystyle\lim_{x \to x_0^{+}} f(x)$ avec la seconde expression.
\item \textbf{Comparer} : si les deux limites latérales sont égales à un même réel $L$, conclure que $\displaystyle\lim_{x \to x_0} f(x) = L$ ; sinon, conclure que $f$ n'admet pas de limite en $x_0$.
\end{enumerate}
\end{methode}

\begin{exoresolu}[Fonction définie par morceaux en son point de jonction]
Soit $h$ la fonction définie sur $\mathbb{R}$ par
\[ h(x) = \begin{cases} 2x + 1 & \text{si } x < 3 \\ 10 - x & \text{si } x \geq 3 \end{cases} \]
La fonction $h$ admet-elle une limite en $3$ ?
\tcblower
\textbf{Solution.} On applique la méthode en calculant les deux limites latérales.

\emph{Limite à gauche} : pour $x < 3$, $h(x) = 2x+1$. Quand $x$ s'approche de $3$, $2x+1$ s'approche de $2\times 3 + 1 = 7$. Donc
\[ \lim_{x \to 3^{-}} h(x) = 7. \]

\emph{Limite à droite} : pour $x \geq 3$, $h(x) = 10 - x$. Quand $x$ s'approche de $3$ par valeurs supérieures, $10-x$ s'approche de $10 - 3 = 7$. Donc
\[ \lim_{x \to 3^{+}} h(x) = 7. \]

\emph{Comparaison} : les deux limites latérales sont égales à $7$. D'après le théorème, $h$ admet une limite en $3$ et
\[ \lim_{x \to 3} h(x) = 7. \]
Remarque : ici $h(3) = 10 - 3 = 7$, la valeur et la limite coïncident — on dira dans la suite de la leçon que $h$ est \emph{continue} en $3$.
\end{exoresolu}

\courssub{Limites infinies en un point et asymptote verticale}

Il reste un dernier comportement possible : $f(x)$ peut ne s'approcher d'aucun nombre fini, mais au contraire devenir arbitrairement grand (en valeur absolue) lorsque $x$ s'approche de $x_0$. L'exemple emblématique est la fonction inverse en $0$.

\begin{definition}[Limite infinie en un point]
Soit $f$ une fonction définie au voisinage de $x_0$, sauf éventuellement en $x_0$.
\begin{itemize}
\item On dit que $f$ tend vers $+\infty$ en $x_0$, et l'on note $\displaystyle\lim_{x \to x_0} f(x) = +\infty$, lorsque $f(x)$ devient aussi grand que l'on veut pourvu que $x$ soit suffisamment proche de $x_0$.
\item On définit de même $\displaystyle\lim_{x \to x_0} f(x) = -\infty$ lorsque $f(x)$ devient négatif, aussi grand que l'on veut en valeur absolue.
\end{itemize}
Ces notions s'appliquent aussi aux limites à droite et à gauche.
\end{definition}

\begin{exemplebox}[La fonction inverse en $0$]
Étudions $f(x) = \dfrac{1}{x}$ au voisinage de $0$.
\begin{itemize}
\item Si $x$ s'approche de $0$ par valeurs \emph{positives} ($0,1$ ; $0,01$ ; $0,001$), alors $\dfrac{1}{x}$ vaut $10$ ; $100$ ; $1000$ : les valeurs deviennent arbitrairement grandes. Donc $\displaystyle\lim_{x \to 0^{+}} \frac{1}{x} = +\infty$.
\item Si $x$ s'approche de $0$ par valeurs \emph{négatives} ($-0,1$ ; $-0,01$ ; $-0,001$), alors $\dfrac{1}{x}$ vaut $-10$ ; $-100$ ; $-1000$. Donc $\displaystyle\lim_{x \to 0^{-}} \frac{1}{x} = -\infty$.
\end{itemize}
Les deux limites latérales sont infinies et de signes opposés : la fonction inverse n'admet pas de limite (finie) en $0$.
\end{exemplebox}

\begin{popfigure}
\begin{center}
\begin{tikzpicture}
\begin{axis}[
    width=0.85\linewidth, height=7cm,
    axis lines=middle,
    xlabel={$x$}, ylabel={$y$},
    xmin=-4, xmax=4, ymin=-6, ymax=6,
    xtick={\empty}, ytick={\empty},
    grid=both, grid style={popline!40},
    samples=200
]
\addplot[popcurve, domain=0.15:4] {1/x};
\addplot[popcurve, domain=-4:-0.15] {1/x};
\draw[dashed, poporange, thick] (axis cs:0,-6) -- (axis cs:0,6);
\node[poporange, right] at (axis cs:0.15,5.2) {$x=0$};
\node[popblue, right] at (axis cs:1.6,1.4) {$y=\dfrac{1}{x}$};
\end{axis}
\end{tikzpicture}
\end{center}
\legende{La fonction inverse : quand $x \to 0^{+}$, la courbe monte vers $+\infty$ ; quand $x \to 0^{-}$, elle descend vers $-\infty$. La droite $x=0$ est asymptote verticale.}
\end{popfigure}

\begin{aretenir}[Asymptote verticale]
Si $\displaystyle\lim_{x \to x_0} f(x) = +\infty$ ou $-\infty$ (éventuellement seulement à droite ou à gauche), alors la droite d'équation $x = x_0$ est \cle{asymptote verticale} à la courbe représentative de $f$ : la courbe s'approche indéfiniment de cette droite sans jamais la toucher.
\end{aretenir}

\begin{attention}[Signe du dénominateur : le réflexe indispensable]
Devant un quotient dont le dénominateur tend vers $0$ en $x_0$, ne conclus jamais trop vite. Il faut étudier le \textbf{signe} du dénominateur de chaque côté de $x_0$. Par exemple, pour $f(x)=\dfrac{1}{x-2}$ : si $x \to 2^{+}$, alors $x-2 > 0$ et $f(x) \to +\infty$ ; si $x \to 2^{-}$, alors $x-2 < 0$ et $f(x) \to -\infty$. Oublier cette distinction est une erreur classique au Baccalauréat.
\end{attention}

\courssub{Synthèse de la section}

\begin{aretenir}[Ce qu'il faut retenir]
\begin{itemize}
\item $\displaystyle\lim_{x \to x_0} f(x) = L$ signifie que $f(x)$ s'approche de $L$ quand $x$ s'approche de $x_0$ ; cette limite est \textbf{indépendante} de l'existence et de la valeur de $f(x_0)$.
\item La limite en $x_0$ existe et vaut $L$ \textbf{si et seulement si} les limites à gauche et à droite existent et valent toutes deux $L$.
\item Si les limites latérales sont différentes, la fonction n'a pas de limite en $x_0$ : c'est un « saut » (fonction en escalier).
\item Si $f(x)$ devient infiniment grand au voisinage de $x_0$, on a une limite infinie et la droite $x = x_0$ est asymptote verticale.
\item Pour une fonction définie par morceaux, on calcule toujours séparément les deux limites latérales au point de jonction, puis on compare.
\end{itemize}
\end{aretenir}

Tu sais désormais décrire rigoureusement le comportement d'une fonction au voisinage d'un point : approche d'une valeur finie, saut brutal, ou envolée vers l'infini. Dans la section suivante, nous transposerons ces idées « à l'infini » : que devient $f(x)$ lorsque $x$ devient très grand ? C'est là qu'apparaîtront les asymptotes horizontales — et que nous retrouverons le coût unitaire d'Aminatou, qui se rapproche de 25 FCFA quand la production croît indéfiniment.

\coursec{Limite d'une fonction en l'infini}

Dans la section précédente, nous avons étudié ce qui se passe lorsque $x$ se rapproche d'un point $x_0$ fixé : la fonction peut admettre une limite finie, infinie, ou des limites différentes à gauche et à droite. Nous changeons maintenant de perspective : au lieu de regarder $x$ s'approcher d'un nombre, nous le laissons \emph{partir à l'infini}. Que devient $f(x)$ lorsque $x$ devient arbitrairement grand, par exemple $x = \num{1000}$, $x = \num{1000000}$, $x = 10^{12}$ ? C'est la question des \cle{limites en l'infini}, indispensable pour décrire le comportement global d'une fonction et tracer sa courbe avec intelligence.

\courssub{Approche intuitive : que se passe-t-il quand $x$ devient très grand ?}

Considérons la fonction $f$ définie sur $]1\,;\,+\infty[$ par
\[ f(x) = \frac{2x+1}{x-1}. \]
Calculons quelques valeurs pour des $x$ de plus en plus grands :

\begin{center}
\begin{tabularx}{\linewidth}{|l|X|X|X|X|X|}
\hline
\rowcolor{popblueL}
$x$ & $10$ & $100$ & $\num{1000}$ & $\num{100000}$ & $10^{6}$ \\
\hline
$f(x)$ & $2{,}333$ & $2{,}0303$ & $2{,}003003$ & $2{,}00003$ & $2{,}000003$ \\
\hline
\end{tabularx}
\end{center}

On observe que $f(x)$ se rapproche de plus en plus de $2$ : l'écart $|f(x)-2|$ devient aussi petit que l'on veut, pourvu que $x$ soit suffisamment grand. On dit que $f$ a pour \cle{limite} $2$ en $+\infty$, et on écrit :
\[ \lim_{x \to +\infty} \frac{2x+1}{x-1} = 2. \]

De même, on peut laisser $x$ devenir très grand \emph{en valeur absolue} du côté des négatifs : $x = -10$, $-1000$, $-10^{6}$. On écrit alors $x \to -\infty$. Pour la même fonction $f$, on constate que $f(-1000) \approx 1{,}997$, $f(-10^{6}) \approx 1{,}999997$ : les valeurs se rapprochent encore de $2$.

\courssub{Limite finie en l'infini et asymptote horizontale}

\begin{definition}[Limite finie en $+\infty$]
Soit $f$ une fonction définie sur un intervalle de la forme $[a\,;\,+\infty[$ et $L$ un nombre réel. On dit que $f$ admet pour limite $L$ en $+\infty$, et on note
\[ \lim_{x \to +\infty} f(x) = L, \]
lorsque, pour tout intervalle ouvert $I$ contenant $L$, il existe un réel $A$ tel que pour tout $x > A$, $f(x) \in I$.

Autrement dit : on peut rendre $f(x)$ aussi proche de $L$ que l'on veut, à condition de prendre $x$ assez grand.
\end{definition}

On définit de manière analogue $\lim_{x \to -\infty} f(x) = L$ : pour tout intervalle ouvert $I$ contenant $L$, il existe un réel $A$ tel que pour tout $x < A$, $f(x) \in I$.

\begin{definition}[Asymptote horizontale]
Si $\lim_{x \to +\infty} f(x) = L$ (respectivement $\lim_{x \to -\infty} f(x) = L$), alors la droite d'équation $y = L$ est appelée \cle{asymptote horizontale} à la courbe $\mathscr{C}_f$ en $+\infty$ (respectivement en $-\infty$).

Graphiquement, la courbe se rapproche indéfiniment de cette droite lorsque $x$ part vers l'infini, sans jamais être obligée de la toucher ni de rester d'un seul côté.
\end{definition}

\begin{popfigure}
\begin{tikzpicture}
\begin{axis}[
  width=\linewidth, height=7cm,
  axis lines=middle,
  xlabel={$x$}, ylabel={$y$},
  xmin=-6, xmax=8, ymin=-4, ymax=7,
  xtick={1}, ytick={2},
  legend style={at={(0.98,0.05)}, anchor=south east, font=\small},
  clip=true
]
\addplot[popcurve,domain=-6:0.75,samples=200]{(2*x+1)/(x-1)};
\addplot[popcurve,domain=1.28:8,samples=200]{(2*x+1)/(x-1)};
\addplot[poporange, dashed, domain=-6:8, samples=2]{2};
\addplot[popgreen, dashed] coordinates {(1,-4) (1,7)};
\legend{$y=\dfrac{2x+1}{x-1}$, asymptote $y=2$, asymptote $x=1$}
\end{axis}
\end{tikzpicture}
\legende{La courbe de $f(x)=\dfrac{2x+1}{x-1}$ admet la droite $y=2$ comme asymptote horizontale en $+\infty$ et en $-\infty$, et la droite $x=1$ comme asymptote verticale.}
\end{popfigure}

\begin{attention}[On ne dit pas que $f(x)$ « vaut » $L$ à l'infini]
Une erreur classique consiste à écrire : « en $+\infty$, $\dfrac{1}{x}$ vaut $0$ ». C'est faux : $+\infty$ n'est pas un nombre, on ne peut pas « remplacer $x$ par $+\infty$ ». Pour tout réel $x$, $\dfrac{1}{x} \neq 0$. La formulation correcte est :
\[ \lim_{x \to +\infty} \frac{1}{x} = 0. \]
La limite est $0$ ; la fonction ne prend jamais la valeur $0$. De même, une courbe peut traverser son asymptote horizontale : l'asymptote décrit un comportement \emph{à l'infini}, pas une interdiction de contact.
\end{attention}

\courssub{Limites infinies en l'infini}

Certaines fonctions ne se stabilisent pas : elles \emph{explosent}. Prenons $g(x) = x^2$ : pour $x = 10$, $100$, $\num{1000}$, on obtient $g(x) = 100$, $\num{10000}$, $\num{1000000}$. Les valeurs dépassent n'importe quel seuil fixé à l'avance.

\begin{definition}[Limite infinie en l'infini]
On dit que $f$ a pour limite $+\infty$ en $+\infty$, et on note
\[ \lim_{x \to +\infty} f(x) = +\infty, \]
lorsque, pour tout réel $A$, il existe un réel $B$ tel que pour tout $x > B$, $f(x) > A$.

Autrement dit : on peut rendre $f(x)$ aussi grand que l'on veut, à condition de prendre $x$ assez grand. On définit de façon analogue les trois autres situations :
\[ \lim_{x \to +\infty} f(x) = -\infty, \qquad \lim_{x \to -\infty} f(x) = +\infty, \qquad \lim_{x \to -\infty} f(x) = -\infty. \]
\end{definition}

\begin{attention}[Une limite infinie est une limite qui n'existe pas au sens fini]
Dire $\lim_{x \to +\infty} f(x) = +\infty$ ne signifie pas que $f$ « a une limite » au même titre qu'une limite finie : cela décrit précisément la manière dont $f$ \emph{n'a pas} de limite finie. En revanche, c'est une information précieuse : la courbe monte (ou descend) sans borne, et il n'y a \textbf{pas d'asymptote horizontale} dans ce cas.
\end{attention}

\courssub{Limites usuelles en l'infini (admises)}

Les résultats suivants sont \textbf{admis} au programme de Terminale A. Ils doivent être connus par cœur, car tous les calculs de limites en l'infini s'appuient dessus.

\begin{propriete}[Limites usuelles en l'infini — admises]
Pour tout entier naturel non nul $n$ et tout réel $a > 0$ :
\begin{itemize}
\item $\lim_{x \to +\infty} x^n = +\infty$ ;
\item $\lim_{x \to -\infty} x^n = +\infty$ si $n$ est pair, et $\lim_{x \to -\infty} x^n = -\infty$ si $n$ est impair ;
\item $\lim_{x \to +\infty} \dfrac{1}{x^n} = 0$ et $\lim_{x \to -\infty} \dfrac{1}{x^n} = 0$ ;
\item $\lim_{x \to +\infty} \sqrt{x} = +\infty$ ;
\item $\lim_{x \to +\infty} |x| = +\infty$ et $\lim_{x \to -\infty} |x| = +\infty$ ;
\item \textbf{Exponentielle :} $\lim_{x \to +\infty} e^x = +\infty$, $\lim_{x \to -\infty} e^x = 0$ ; plus généralement pour $a>1$, $\lim_{x \to +\infty} a^x = +\infty$, $\lim_{x \to -\infty} a^x = 0$ ; pour $0<a<1$, $\lim_{x \to +\infty} a^x = 0$, $\lim_{x \to -\infty} a^x = +\infty$ ;
\item \textbf{Logarithme népérien :} $\lim_{x \to +\infty} \ln x = +\infty$, $\lim_{x \to 0^+} \ln x = -\infty$.
\end{itemize}
\end{propriete}

\begin{aretenir}[Tableau récapitulatif des limites usuelles en l'infini]
\begin{center}
\begin{tabularx}{\linewidth}{|l|X|X|}
\hline
\rowcolor{popblueL}
\textbf{Fonction} & \textbf{Limite en $+\infty$} & \textbf{Limite en $-\infty$} \\
\hline
$x^n$ ($n$ pair) & $+\infty$ & $+\infty$ \\
\hline
$x^n$ ($n$ impair) & $+\infty$ & $-\infty$ \\
\hline
$\dfrac{1}{x^n}$ ($n \in \mathbb{N}^*$) & $0$ & $0$ \\
\hline
$\sqrt{x}$ & $+\infty$ & non définie \\
\hline
$|x|$ & $+\infty$ & $+\infty$ \\
\hline
$e^x$ & $+\infty$ & $0$ \\
\hline
$a^x$ ($a>1$) & $+\infty$ & $0$ \\
\hline
$a^x$ ($0<a<1$) & $0$ & $+\infty$ \\
\hline
$\ln x$ & $+\infty$ & non définie (défini sur $]0,+\infty[$) \\
\hline
$\cos x$, $\sin x$ & pas de limite & pas de limite \\
\hline
\end{tabularx}
\end{center}
\end{aretenir}

\begin{popfigure}
\begin{tikzpicture}
\begin{axis}[
  width=\linewidth, height=7.5cm,
  axis lines=middle,
  xlabel={$x$}, ylabel={$y$},
  xmin=0, xmax=4, ymin=0, ymax=16,
  legend style={at={(0.03,0.97)}, anchor=north west, font=\small},
  clip=true
]
\addplot[popcurve,popblue,domain=0:4,samples=100]{x};
\addplot[popcurve,poporange,domain=0:4,samples=100]{x^2};
\addplot[popcurve,poppurple,domain=0:2.55,samples=100]{x^3};
\addplot[popcurve,popgreen,domain=0:4,samples=100]{sqrt(max(0,x))};
\addplot[popcurve,poppink,domain=0:2.8,samples=100]{exp(x)};
\legend{$y=x$, $y=x^2$, $y=x^3$, $y=\sqrt{x}$, $y=e^x$}
\end{axis}
\end{tikzpicture}
\legende{Comparaison des croissances en $+\infty$ : toutes ces fonctions tendent vers $+\infty$, mais $e^x$ croît plus vite que toute puissance $x^n$, qui croît plus vite que $x$, qui croît plus vite que $\sqrt{x}$.}
\end{popfigure}

\courssub{Hiéarchie des croissances (repères pour le Bac)}

\begin{propriete}[Comparaison des croissances en $+\infty$ — admis]
Lorsque $x \to +\infty$, on a l'ordre de grandeur suivant (de la plus lente à la plus rapide) :
\[ \ln x \ll x^a \ll x^b \ll e^x \quad \text{pour } 0 < a < b. \]
Plus précisément : $\lim_{x \to +\infty} \frac{\ln x}{x^a} = 0$, $\lim_{x \to +\infty} \frac{x^a}{e^x} = 0$.
\end{propriete}

\begin{attention}[Hiérarchie et formes indéterminées]
Cette hiérarchie permet de lever immédiatement certaines formes indéterminées de type $\frac{\infty}{\infty}$ ou $\infty - \infty$ sans calcul long. Exemple : $\lim_{x \to +\infty} \frac{x^2 + 3x}{e^x} = 0$ car l'exponentielle domine toute puissance polynomiale.
\end{attention}

\courssub{Fonctions sans limite en l'infini}

Toutes les fonctions n'ont pas nécessairement une limite, finie ou infinie, en l'infini. Les fonctions \cle{périodiques} fournissent les exemples les plus célèbres.

\begin{exemplebox}[Les fonctions sinus et cosinus n'ont pas de limite en l'infini (admis)]
La fonction $x \mapsto \cos x$ prend indéfiniment les valeurs $1$ (en $x = 0$, $2\pi$, $4\pi$, \ldots) et $-1$ (en $x = \pi$, $3\pi$, \ldots). Elle oscille sans cesse entre $-1$ et $1$ : elle ne se rapproche d'aucun nombre $L$, et ne part ni vers $+\infty$ ni vers $-\infty$. On dit que $\cos$ (et de même $\sin$) \textbf{n'a pas de limite} en $+\infty$ ni en $-\infty$. Ce résultat est admis.
\end{exemplebox}

\begin{attention}[Ne pas confondre « pas de limite finie » et « limite infinie »]
Il y a trois situations possibles en $+\infty$ : limite finie $L$ (asymptote horizontale), limite infinie ($\pm\infty$), ou \emph{absence totale de limite} (cas de $\cos$ et $\sin$). Ne concluez jamais trop vite : une fonction qui n'a pas de limite finie n'a pas forcément une limite infinie.
\end{attention}

\courssub{Exemple chiffré détaillé et lecture graphique}

\begin{exemplebox}[Limite de $\dfrac{3x-5}{x+2}$ en $+\infty$]
Soit $f$ définie sur $]-2\,;\,+\infty[$ par $f(x) = \dfrac{3x-5}{x+2}$. L'idée directrice : quand $x$ est très grand, les constantes $-5$ et $+2$ deviennent négligeables devant $3x$ et $x$. On met donc $x$ en facteur au numérateur et au dénominateur :
\[ f(x) = \frac{x\left(3 - \dfrac{5}{x}\right)}{x\left(1 + \dfrac{2}{x}\right)} = \frac{3 - \dfrac{5}{x}}{1 + \dfrac{2}{x}}. \]
Or $\lim_{x \to +\infty} \dfrac{5}{x} = 0$ et $\lim_{x \to +\infty} \dfrac{2}{x} = 0$ (limites usuelles). Donc le numérateur tend vers $3$ et le dénominateur vers $1$, d'où :
\[ \lim_{x \to +\infty} \frac{3x-5}{x+2} = \frac{3}{1} = 3. \]
\textbf{Interprétation graphique :} la droite d'équation $y = 3$ est asymptote horizontale à la courbe de $f$ en $+\infty$.
\end{exemplebox}

\begin{exoresolu}[Lecture d'asymptotes sur une courbe]
On donne la courbe $\mathscr{C}$ d'une fonction $g$ définie sur $\mathbb{R} \setminus \{1\}$. On constate graphiquement que :
\begin{itemize}
\item quand $x \to +\infty$, la courbe se rapproche de la droite $y = -1$ ;
\item quand $x \to -\infty$, la courbe se rapproche aussi de la droite $y = -1$ ;
\item quand $x \to 1$ par valeurs supérieures, la courbe monte indéfiniment.
\end{itemize}
Énoncé : traduire chaque observation par une égalité de limites et nommer les asymptotes éventuelles.
\tcblower
\textbf{Solution.}
\begin{itemize}
\item La courbe se rapproche de $y = -1$ en $+\infty$ : $\lim_{x \to +\infty} g(x) = -1$. La droite $y = -1$ est \textbf{asymptote horizontale} en $+\infty$.
\item De même en $-\infty$ : $\lim_{x \to -\infty} g(x) = -1$. La même droite $y = -1$ est asymptote horizontale en $-\infty$ (une même droite peut être asymptote des deux côtés).
\item La courbe monte indéfiniment quand $x \to 1^+$ : $\lim_{\substack{x \to 1 \\ x > 1}} g(x) = +\infty$. La droite $x = 1$ est \textbf{asymptote verticale} à $\mathscr{C}$.
\end{itemize}
\end{exoresolu}

\begin{methode}[Déterminer graphiquement une asymptote horizontale]
\begin{enumerate}
\item Observer le comportement de la courbe aux extrémités du graphique (à droite pour $+\infty$, à gauche pour $-\infty$).
\item Repérer la droite horizontale $y = L$ dont la courbe se rapproche indéfiniment.
\item Conclure par l'écriture rigoureuse : $\lim_{x \to +\infty} f(x) = L$, donc la droite $y = L$ est asymptote horizontale en $+\infty$.
\item Si la courbe monte ou descend sans se stabiliser, conclure à une limite infinie (pas d'asymptote horizontale).
\end{enumerate}
\end{methode}

\begin{savaistu}[Les asymptotes modélisent la saturation]
La notion d'asymptote horizontale n'est pas qu'un jeu de mathématicien : elle décrit les phénomènes de \emph{saturation}. Lorsqu'on charge un condensateur, la tension à ses bornes croît selon une loi du type $u(t) = E\left(1 - e^{-t/\tau}\right)$ et tend vers la limite $E$ : la droite $u = E$ est asymptote, la tension ne dépassera jamais la tension du générateur. En médecine, la concentration d'un médicament administré en perfusion continue dans le sang se stabilise vers une valeur seuil : là encore, une asymptote horizontale. Au Cameroun, les ingénieurs de CAMTEL ou d'ENEO rencontrent ces modèles dès qu'ils étudient la charge d'un circuit ou la montée en régime d'un réseau : comprendre les limites en l'infini, c'est comprendre le comportement à long terme des systèmes.
\end{savaistu}

\begin{experience}[Découverte : asymptote oblique, au-delà de l'horizontale]
La fonction $f(x) = x + \frac{1}{x}$ n'a pas de limite finie en $+\infty$ (elle tend vers $+\infty$), donc pas d'asymptote horizontale. Pourtant, sa courbe se rapproche d'une droite \emph{non horizontale}.
\begin{enumerate}
\item Calculer $\lim_{x \to +\infty} \frac{f(x)}{x}$.
\item Calculer $\lim_{x \to +\infty} \bigl(f(x) - x\bigr)$.
\item En déduire l'équation de la droite vers laquelle la courbe se rapproche.
\end{enumerate}
\tcblower
\textbf{Solution guidée.}
\begin{enumerate}
\item $\frac{f(x)}{x} = 1 + \frac{1}{x^2} \to 1$. Le coefficient directeur est $a = 1$.
\item $f(x) - x = \frac{1}{x} \to 0$. L'ordonnée à l'origine est $b = 0$.
\item La droite $y = x$ est \textbf{asymptote oblique} en $+\infty$ (et en $-\infty$).
\end{enumerate}
Cette méthode (recherche de $a$ puis $b$) sera formalisée après l'étude des opérations sur les limites.
\end{experience}

\begin{exercice}[Pour s'entraîner]
\begin{enumerate}
\item Déterminer les limites suivantes en justifiant par les limites usuelles :
\[ \lim_{x \to +\infty} x^4, \qquad \lim_{x \to -\infty} x^5, \qquad \lim_{x \to -\infty} \frac{1}{x^3}, \qquad \lim_{x \to +\infty} \sqrt{x}, \qquad \lim_{x \to +\infty} e^{-x}, \qquad \lim_{x \to +\infty} \ln x. \]
\item Soit $h(x) = \dfrac{5x + 3}{2x - 1}$. Calculer $\lim_{x \to +\infty} h(x)$ et interpréter graphiquement le résultat.
\item La fonction $x \mapsto \sin x$ admet-elle une asymptote horizontale en $+\infty$ ? Justifier.
\item \textbf{Défi :} En utilisant la hiérarchie des croissances, donner sans calcul détaillé :
\[ \lim_{x \to +\infty} \frac{x^3 + 2x}{e^x}, \qquad \lim_{x \to +\infty} \frac{\ln x}{x}, \qquad \lim_{x \to +\infty} \left( \sqrt{x^2+x} - x \right) \text{ (indice : quantité conjuguée)}. \]
\end{enumerate}
\end{exercice}

\begin{corrige}[Exercice « Pour s'entraîner »]
\begin{enumerate}
\item $\lim_{x \to +\infty} x^4 = +\infty$ (puissance paire) ; $\lim_{x \to -\infty} x^5 = -\infty$ (puissance impaire) ; $\lim_{x \to -\infty} \dfrac{1}{x^3} = 0$ ; $\lim_{x \to +\infty} \sqrt{x} = +\infty$ ; $\lim_{x \to +\infty} e^{-x} = 0$ ; $\lim_{x \to +\infty} \ln x = +\infty$.
\item On factorise $x$ : $h(x) = \dfrac{5 + \dfrac{3}{x}}{2 - \dfrac{1}{x}}$. Comme $\dfrac{3}{x} \to 0$ et $\dfrac{1}{x} \to 0$ en $+\infty$, on obtient $\lim_{x \to +\infty} h(x) = \dfrac{5}{2}$. La droite $y = \dfrac{5}{2}$ est asymptote horizontale en $+\infty$.
\item Non : $\sin x$ oscille indéfiniment entre $-1$ et $1$ et n'a pas de limite en $+\infty$ (admis). Sans limite finie, il ne peut pas y avoir d'asymptote horizontale.
\item $\frac{x^3}{e^x} \to 0$ (exponentielle domine) donc première limite $0$. $\frac{\ln x}{x} \to 0$ (puissance domine logarithme) donc deuxième limite $0$. Pour la troisième : $\sqrt{x^2+x} - x = \frac{x}{\sqrt{x^2+x}+x} = \frac{1}{\sqrt{1+1/x}+1} \to \frac{1}{2}$.
\end{enumerate}
\end{corrige}

Nous savons désormais décrire le comportement d'une fonction aux extrémités de son domaine, y compris la recherche d'asymptotes horizontales et l'intuition des asymptotes obliques. Mais les fonctions rencontrées au Baccalauréat sont souvent des combinaisons de fonctions simples : sommes, produits, quotients, composées. La section suivante établit les règles opératoires qui permettent de calculer leurs limites sans revenir chaque fois aux définitions.

\coursec{Opérations sur les limites}

Dans la section précédente, tu as appris à lire la limite d'une fonction en un point et en l'infini, et tu as mémorisé les limites usuelles (puissances, racine carrée, inverse). Mais dans la pratique, les fonctions que l'on rencontre au Baccalauréat ne sont presque jamais des fonctions « de base » : ce sont des \cle{sommes}, des \cle{produits}, des \cle{quotients} et des \cle{composées} de fonctions simples. Par exemple, le chiffre d'affaires d'une boutique de Douala modélisé par $f(x) = 2x^3 - x + 1$, ou un taux moyen $\dfrac{2x^3 - x + 1}{x^3 + 4x}$. Comment déterminer leurs limites sans repartir de la définition ? C'est toute l'ambition de cette section : des \cle{théorèmes d'opérations} qui permettent de calculer des limites \emph{par morceaux}, puis une règle redoutablement efficace pour les polynômes et les fractions rationnelles.

\courssub{Limite d'une somme de fonctions}

\textbf{Intuition.} Si, lorsque $x$ se rapproche de $a$, la quantité $f(x)$ se rapproche de $3$ et $g(x)$ se rapproche de $-5$, il est naturel que $f(x) + g(x)$ se rapproche de $3 + (-5) = -2$. De même, si $f(x)$ devient immense (tend vers $+\infty$) et que $g(x)$ reste proche d'une valeur finie, la somme devient immense. Le seul cas délicat est celui où deux infinis se « combattent » : $+\infty$ d'un côté, $-\infty$ de l'autre. Qui l'emporte ? Tout dépend des fonctions : on parlera de \cle{forme indéterminée}.

\begin{propriete}[Limite d'une somme (admise)]
Soient $f$ et $g$ deux fonctions admettant chacune une limite (finie ou infinie) en un point $a$ (fini ou infini). La limite de $f + g$ en $a$ est donnée par le tableau suivant :
\begin{center}
\begin{tabularx}{\linewidth}{|c|X|X|X|}
\hline
\rowcolor{popblueL}
$\lim f$ & $\ell$ (finie) & $+\infty$ & $-\infty$ \\
\hline
$\ell'$ (finie) & $\ell + \ell'$ & $+\infty$ & $-\infty$ \\
\hline
$+\infty$ & $+\infty$ & $+\infty$ & \cellcolor{poppinkL}\textbf{F.I.} \\
\hline
$-\infty$ & $-\infty$ & \cellcolor{poppinkL}\textbf{F.I.} & $-\infty$ \\
\hline
\end{tabularx}
\end{center}
La case \textbf{F.I.} (forme indéterminée) correspond au cas $(+\infty) + (-\infty)$, noté « $\infty - \infty$ » : aucune conclusion générale n'est possible, il faut transformer l'expression.
\end{propriete}

\begin{exemplebox}[Sommes de limites]
\begin{enumerate}
\item $\displaystyle\lim_{x \to +\infty} \left(x^2 + \dfrac{1}{x}\right)$ : on a $\lim_{x\to+\infty} x^2 = +\infty$ et $\lim_{x\to+\infty} \dfrac{1}{x} = 0$. Par somme : la limite vaut $+\infty$.
\item $\displaystyle\lim_{x \to 0^+} \left(x^3 + \sqrt{x} + 2\right) = 0 + 0 + 2 = 2$.
\item $\displaystyle\lim_{x \to +\infty} \left(x^2 - x\right)$ : c'est le cas $(+\infty) + (-\infty)$, une forme indéterminée. On ne peut \textbf{pas} conclure directement (nous la lèverons plus loin par factorisation).
\end{enumerate}
\end{exemplebox}

\courssub{Limite d'un produit de fonctions}

\textbf{Intuition.} Multiplier des quantités qui se stabilisent revient à multiplier leurs valeurs limites. Multiplier une quantité énorme par une quantité proche d'un nombre non nul donne une quantité énorme (le signe dépend du signe du nombre). En revanche, multiplier une quantité énorme par une quantité qui s'écrase vers $0$ est un duel imprévisible : c'est la forme indéterminée « $0 \times \infty$ ».

\begin{propriete}[Limite d'un produit (admise)]
Avec les mêmes hypothèses, la limite de $f \times g$ est donnée par :
\begin{center}
\begin{tabularx}{\linewidth}{|c|X|X|X|}
\hline
\rowcolor{popblueL}
$\lim f$ & $\ell \neq 0$ & $0$ & $\pm\infty$ \\
\hline
$\ell'$ (finie) & $\ell \times \ell'$ & $0$ & $\pm\infty$ (règle des signes) \\
\hline
$0$ & $0$ & $0$ & \cellcolor{poppinkL}\textbf{F.I.} \\
\hline
$\pm\infty$ & $\pm\infty$ (règle des signes) & \cellcolor{poppinkL}\textbf{F.I.} & $\pm\infty$ (règle des signes) \\
\hline
\end{tabularx}
\end{center}
Pour les cases infinies, on applique la \cle{règle des signes} habituelle : $(+)\times(+) = (+)$, $(+)\times(-) = (-)$, etc. La forme indéterminée est « $0 \times \infty$ ».
\end{propriete}

\begin{exemplebox}[Produits de limites]
\begin{enumerate}
\item $\displaystyle\lim_{x \to +\infty} x^2 \times \left(2 + \dfrac{1}{x}\right)$ : $\lim x^2 = +\infty$ et $\lim \left(2 + \tfrac{1}{x}\right) = 2 > 0$. Par produit : $+\infty$.
\item $\displaystyle\lim_{x \to +\infty} (-3x) \times \dfrac{1}{x^2}$ : attention, $-\infty \times 0$ est une forme indéterminée ! En simplifiant : $-3x \times \dfrac{1}{x^2} = \dfrac{-3}{x}$, dont la limite est $0$. La F.I. cachait en réalité une limite nulle.
\end{enumerate}
\end{exemplebox}

\courssub{Limite d'un quotient de fonctions}

\textbf{Intuition.} Diviser par une quantité qui devient énorme écrase le résultat vers $0$ ; diviser par une quantité qui s'approche de $0$ fait exploser le résultat (vers l'infini, avec un signe à déterminer). Les duels « $0/0$ » et « $\infty/\infty$ » sont indéterminés.

\begin{propriete}[Limite d'un quotient (admise)]
La limite de $\dfrac{f}{g}$ est donnée par :
\begin{center}
\begin{tabularx}{\linewidth}{|c|X|X|X|}
\hline
\rowcolor{popblueL}
$\lim f$ & $\ell'$ (finie) $\neq 0$ & $0$ & $\pm\infty$ \\
\hline
$\ell$ (finie) & $\dfrac{\ell}{\ell'}$ & $\dfrac{0}{0}$ : \cellcolor{poppinkL}\textbf{F.I.} & $0$ \\
\hline
$\pm\infty$ & $\pm\infty$ (signe par la règle des signes) & $\pm\infty$ (signe à étudier) & $\dfrac{\infty}{\infty}$ : \cellcolor{poppinkL}\textbf{F.I.} \\
\hline
\end{tabularx}
\end{center}
Précision importante : si le dénominateur tend vers $0$ \emph{en gardant un signe constant} et que le numérateur tend vers une limite finie non nulle, le quotient tend vers $+\infty$ ou $-\infty$ selon les signes. C'est pourquoi l'étude du \cle{signe du dénominateur} est indispensable (limite à gauche / à droite).
\end{propriete}

\begin{exemplebox}[Quotients de limites]
\begin{enumerate}
\item $\displaystyle\lim_{x \to +\infty} \dfrac{5}{x^2 + 1}$ : le numérateur tend vers $5$, le dénominateur vers $+\infty$, donc le quotient tend vers $0$.
\item $\displaystyle\lim_{\substack{x \to 2 \\ x > 2}} \dfrac{3}{x - 2}$ : le numérateur tend vers $3 > 0$, le dénominateur tend vers $0$ par valeurs \emph{positives}, donc la limite est $+\infty$. À gauche de $2$, elle serait $-\infty$ : la fonction n'a pas de limite en $2$, seulement des limites unilatérales.
\end{enumerate}
\end{exemplebox}

\courssub{Les quatre formes indéterminées}

\begin{definition}[Formes indéterminées]
On appelle \cle{forme indéterminée} (F.I.) chacune des quatre situations suivantes, où les théorèmes d'opérations ne permettent \emph{pas} de conclure :
\[ \infty - \infty \qquad ; \qquad 0 \times \infty \qquad ; \qquad \dfrac{0}{0} \qquad ; \qquad \dfrac{\infty}{\infty}. \]
\end{definition}

\begin{aretenir}[Une F.I. n'est pas une impossibilité]
Une forme indéterminée ne signifie \textbf{pas} « la limite n'existe pas » ni « on ne peut pas savoir ». Elle signifie : \emph{« l'écriture actuelle de la fonction ne permet pas de conclure ; il faut la transformer »}. Les techniques de transformation (factorisation du terme dominant, quantité conjuguée, simplification) font l'objet de la section suivante. Deux fonctions présentant la même F.I. peuvent avoir des limites totalement différentes :
\begin{itemize}
\item $\dfrac{x^2}{x} \to 0$, $\dfrac{x}{x^2} \to +\infty$, $\dfrac{3x}{x} \to 3$ en $+\infty$ : pourtant ce sont trois fois la forme « $\infty/\infty$ ».
\end{itemize}
\end{aretenir}

\begin{attention}[Erreur fréquente : conclure hâtivement sur une F.I.]
L'erreur la plus grave — et la plus sanctionnée au Baccalauréat — consiste à écrire par exemple : « $\lim_{x\to+\infty}(x^2 - x) = +\infty - \infty = 0$ » ou « $\dfrac{\infty}{\infty} = 1$ ». Ces « calculs » sur le symbole $\infty$ sont \textbf{faux} et interdits : $\infty$ n'est pas un nombre réel, on ne peut pas lui appliquer les opérations usuelles. La bonne démarche est :
\begin{enumerate}
\item repérer la forme indéterminée et \textbf{l'annoncer} dans la rédaction ;
\item transformer l'expression (factoriser, simplifier, multiplier par le conjugué) ;
\item appliquer les théorèmes d'opérations à la nouvelle écriture.
\end{enumerate}
\end{attention}

\courssub{Limite d'une fonction composée}

\textbf{Intuition.} Une fonction composée, c'est une machine à deux étages : $x$ entre dans la première machine $u$, le résultat $u(x)$ entre dans la seconde machine $g$. Si la sortie du premier étage se rapproche d'une valeur $a$, et que la seconde machine envoie les entrées proches de $a$ vers $L$, alors la sortie finale se rapproche de $L$. C'est le principe du \emph{relais} : la limite se transmet d'étage en étage.

\begin{propriete}[Limite d'une composée (admise)]
Soient $u$ et $g$ deux fonctions. Si
\[ \lim_{x \to x_0} u(x) = a \qquad \text{et} \qquad \lim_{y \to a} g(y) = L, \]
alors
\[ \lim_{x \to x_0} g\big(u(x)\big) = L. \]
Ce résultat reste valable lorsque $x_0$, $a$ ou $L$ sont infinis. On dit qu'on effectue le \cle{changement de variable} $y = u(x)$.
\end{propriete}

\begin{popfigure}
\begin{tikzpicture}[
  box/.style={draw=popblue, thick, rounded corners, fill=popblueL, minimum width=2.6cm, minimum height=1.1cm, align=center, font=\small},
  arr/.style={-{Stealth[length=3mm]}, very thick, poporange}
]
\node[box] (x) at (0,0) {$x$};
\node[box] (u) at (4.2,0) {$y = u(x)$};
\node[box] (g) at (8.6,0) {$g(y) = g(u(x))$};
\draw[arr] (x) -- (u) node[midway, above, font=\small, popdark]{$u$};
\draw[arr] (u) -- (g) node[midway, above, font=\small, popdark]{$g$};
\node[below=0.35cm of x, font=\small, poppurple] {$x \to x_0$};
\node[below=0.35cm of u, font=\small, poppurple] {$y \to a$};
\node[below=0.35cm of g, font=\small, poppurple] {$g(y) \to L$};
\end{tikzpicture}
\legende{La composition en arbre : la limite se transmet à chaque étage, de $x$ vers $u(x)$, puis vers $g(u(x))$.}
\end{popfigure}

\begin{methode}[Appliquer le théorème de composition]
Pour calculer $\lim_{x \to x_0} g(u(x))$ :
\begin{enumerate}
\item \textbf{Identifier} la fonction « intérieure » $u$ et la fonction « extérieure » $g$.
\item \textbf{Calculer} $a = \lim_{x \to x_0} u(x)$.
\item \textbf{Poser} $y = u(x)$ et calculer $L = \lim_{y \to a} g(y)$ (limite usuelle).
\item \textbf{Conclure} : $\lim_{x \to x_0} g(u(x)) = L$, en rédigeant la chaîne complète.
\end{enumerate}
\end{methode}

\begin{exemplebox}[Composées avec $\sqrt{u}$ et $\dfrac{1}{u}$]
\begin{enumerate}
\item Calculons $\displaystyle\lim_{x \to +\infty} \sqrt{x^2 + 1}$. On pose $u(x) = x^2 + 1$ : $\lim_{x\to+\infty} u(x) = +\infty$. Or $\lim_{y \to +\infty} \sqrt{y} = +\infty$. Donc $\lim_{x\to+\infty} \sqrt{x^2+1} = +\infty$.
\item Calculons $\displaystyle\lim_{x \to +\infty} \dfrac{1}{x^2 + 3}$. On pose $u(x) = x^2 + 3 \to +\infty$ ; or $\lim_{y\to+\infty} \dfrac{1}{y} = 0$. Donc la limite cherchée est $0$.
\item Calculons $\displaystyle\lim_{x \to 3} \sqrt{2x + 3}$. On a $u(x) = 2x+3 \to 9$ et $\lim_{y \to 9} \sqrt{y} = 3$. Donc $\lim_{x \to 3} \sqrt{2x+3} = 3$.
\end{enumerate}
\end{exemplebox}

\courssub{Règle du terme de plus haut degré}

Voici le résultat pratique le plus utilisé de tout le chapitre : il transforme instantanément une forme indéterminée « $\infty - \infty$ » ou « $\infty/\infty$ » en un calcul immédiat.

\begin{propriete}[Limite d'un polynôme et d'une fraction rationnelle en l'infini]
\begin{itemize}
\item En $+\infty$ et en $-\infty$, un \cle{polynôme} a la même limite que son \cle{terme de plus haut degré}.
\item En $+\infty$ et en $-\infty$, une \cle{fraction rationnelle} (quotient de deux polynômes) a la même limite que le \cle{quotient des termes de plus haut degré}.
\end{itemize}
\end{propriete}

\begin{experience}[Démonstration guidée : pourquoi le terme dominant commande]
Considérons le polynôme $P(x) = 2x^3 - x + 1$ en $+\infty$. Directement, on obtient « $+\infty - \infty + 1$ » : forme indéterminée. L'idée est de \textbf{factoriser le terme dominant} $x^3$ :
\begin{enumerate}
\item \textbf{Factorise} : pour $x \neq 0$,
\[ P(x) = x^3\left(2 - \dfrac{1}{x^2} + \dfrac{1}{x^3}\right). \]
Vérifie en développant : on retombe bien sur $2x^3 - x + 1$.
\item \textbf{Passe à la limite dans la parenthèse} : par les limites usuelles, $\dfrac{1}{x^2} \to 0$ et $\dfrac{1}{x^3} \to 0$, donc par somme :
\[ \lim_{x \to +\infty} \left(2 - \dfrac{1}{x^2} + \dfrac{1}{x^3}\right) = 2. \]
\item \textbf{Applique le théorème du produit} : $x^3 \to +\infty$ et la parenthèse tend vers $2 > 0$, donc
\[ \lim_{x \to +\infty} P(x) = +\infty, \]
qui est bien la limite du seul terme $2x^3$.
\end{enumerate}
Le même raisonnement fonctionne pour tout polynôme : après factorisation par $x^n$, la parenthèse tend toujours vers le coefficient dominant, qui est non nul. Pour une fraction rationnelle, on factorise le terme dominant au numérateur \emph{et} au dénominateur, puis on simplifie les puissances de $x$. Cette démonstration est \textbf{formatrice} : retiens la technique de factorisation du terme dominant, elle sera exigée dans les exercices.
\end{experience}

\begin{exoresolu}[Limite d'une fraction rationnelle en $+\infty$ et en $-\infty$]
Énoncé. Déterminer $\displaystyle\lim_{x \to +\infty} f(x)$ et $\displaystyle\lim_{x \to -\infty} f(x)$, où $f(x) = \dfrac{2x^3 - x + 1}{x^3 + 4x}$.
\tcblower
\textbf{Solution détaillée.}

\emph{Étape 1 : observation.} En $+\infty$, le numérateur et le dénominateur tendent vers $+\infty$ : c'est la forme indéterminée « $\dfrac{\infty}{\infty}$ ». On ne peut pas conclure directement.

\emph{Étape 2 : application de la règle.} D'après la règle du terme de plus haut degré, $f(x)$ a en l'infini la même limite que le quotient des termes dominants :
\[ \lim_{x \to +\infty} f(x) = \lim_{x \to +\infty} \dfrac{2x^3}{x^3} = \lim_{x \to +\infty} 2 = 2. \]

\emph{Étape 3 : justification rigoureuse par factorisation} (rédaction attendue au Bac). Pour $x \neq 0$ :
\[ f(x) = \dfrac{x^3\left(2 - \dfrac{1}{x^2} + \dfrac{1}{x^3}\right)}{x^3\left(1 + \dfrac{4}{x^2}\right)} = \dfrac{2 - \dfrac{1}{x^2} + \dfrac{1}{x^3}}{1 + \dfrac{4}{x^2}}. \]
Par les limites usuelles et les opérations : le numérateur tend vers $2$, le dénominateur vers $1 \neq 0$. Par quotient :
\[ \lim_{x \to +\infty} f(x) = \dfrac{2}{1} = 2. \]

\emph{Étape 4 : en $-\infty$.} Les termes dominants sont les mêmes : $\dfrac{2x^3}{x^3} = 2$ pour tout $x \neq 0$, donc
\[ \lim_{x \to -\infty} f(x) = 2. \]
(Le signe de l'infini ne change rien ici car les degrés sont égaux ; il aurait été décisif pour $\dfrac{2x^3}{x^2} = 2x$, qui tend vers $-\infty$ en $-\infty$.)

\textbf{Conclusion} : $\lim_{x \to +\infty} f(x) = \lim_{x \to -\infty} f(x) = 2$.
\end{exoresolu}

\begin{attention}[La règle du terme dominant ne s'applique qu'en l'infini]
La règle du terme de plus haut degré est valable \textbf{uniquement} en $+\infty$ et en $-\infty$. En un point fini $x_0$, elle est fausse en général ! Par exemple, pour $P(x) = 2x^3 - x + 1$, on a $P(0) = 1$, alors que le terme dominant $2x^3$ vaut $0$ en $0$. En un point fini, on calcule simplement l'image (ou l'on lève la F.I. par factorisation par $(x - x_0)$, technique vue dans la section suivante).
\end{attention}

\begin{savaistu}[Pourquoi ces théorèmes sont-ils « admis » ?]
Les théorèmes d'opérations sur les limites se démontrent rigoureusement à l'université, à l'aide de la définition « epsilon-delta » introduite par le mathématicien allemand Karl Weierstrass au XIX\textsuperscript{e} siècle. Cette définition, qui a mis fin à deux mille ans de débats sur l'infiniment petit depuis Archimède et Newton, dépasse le programme de Terminale. Au Baccalauréat, tu dois donc \emph{citer} ces théorèmes (« par somme », « par produit », « par quotient », « par composition ») à chaque étape de tes calculs de limites : c'est la justification qui fait la différence entre une copie correcte et une copie excellente.
\end{savaistu}

\begin{exercice}[À toi de jouer]
Déterminer les limites suivantes en justifiant chaque étape par le théorème utilisé :
\begin{enumerate}
\item $\displaystyle\lim_{x \to +\infty} \left(3x^2 - 5x + 7\right)$ ;
\item $\displaystyle\lim_{x \to -\infty} \dfrac{-4x^2 + x}{x^2 + 1}$ ;
\item $\displaystyle\lim_{x \to +\infty} \sqrt{\dfrac{x^2 + 1}{4x^2 - 3}}$ ;
\item $\displaystyle\lim_{x \to 1} \dfrac{2x + 1}{x + 3}$.
\end{enumerate}
\end{exercice}

\begin{corrige}[Corrigé des exercices]
\begin{enumerate}
\item Forme « $\infty - \infty$ ». Par la règle du terme dominant : $\lim_{x\to+\infty} 3x^2 = +\infty$, donc la limite est $+\infty$.
\item Par la règle du terme dominant : $\dfrac{-4x^2}{x^2} = -4$, donc la limite est $-4$.
\item Par la règle du terme dominant, la fraction intérieure tend vers $\dfrac{1}{4}$. Par composition avec $y \mapsto \sqrt{y}$, continue en $\tfrac14$ : la limite est $\sqrt{\tfrac14} = \dfrac{1}{2}$.
\item Pas d'indétermination : le numérateur tend vers $3$, le dénominateur vers $4 \neq 0$. Par quotient : la limite est $\dfrac{3}{4}$.
\end{enumerate}
\end{corrige}

Tu disposes maintenant de toute la « boîte à outils » algébrique des limites. Reste à apprendre les gestes techniques qui lèvent les formes indéterminées — factorisation, quantité conjuguée — : c'est précisément l'objet de la section suivante.

\coursec{Lever une forme indéterminée}

Dans la section précédente, nous avons établi les règles opératoires sur les limites : somme, produit, quotient, composition. Mais nous avons buté sur quatre cas rebelles, les \cle{formes indéterminées} : $\dfrac{0}{0}$, $\dfrac{\infty}{\infty}$, $0 \times \infty$ et $\infty - \infty$. Pour ces formes, aucune règle générale ne permet de conclure directement. Cette section est la boîte à outils du candidat au Baccalauréat : tu vas apprendre à \emph{transformer} l'expression pour faire disparaître l'indétermination, puis à conclure proprement. C'est l'une des compétences les plus rentables de l'épreuve.

\courssub{Qu'est-ce qu'une forme indéterminée ? Stratégie générale}

\begin{definition}[Forme indéterminée]
On dit qu'un calcul de limite présente une \cle{forme indéterminée} (on écrit souvent F.I.) lorsque l'application directe des théorèmes sur les opérations conduit à l'une des quatre écritures suivantes, qui ne permettent pas de conclure :
\[ \frac{0}{0} \qquad ; \qquad \frac{\infty}{\infty} \qquad ; \qquad 0 \times \infty \qquad ; \qquad \infty - \infty. \]
\end{definition}

\begin{attention}[Ce que « indéterminée » ne veut pas dire]
Une forme indéterminée ne signifie \textbf{pas} que la limite n'existe pas, ni qu'elle vaut $0$ ou $1$. Elle signifie seulement : « avec les seules règles opératoires, on ne peut pas décider ». Il faut alors \emph{transformer l'expression} de la fonction, sans changer sa valeur, pour faire apparaître une forme déterminée.
\end{attention}

\begin{savaistu}[Pourquoi dit-on « indéterminée » ?]
Prenons la forme $\dfrac{0}{0}$ en $0$. Selon le numérateur et le dénominateur, tout peut arriver :
\[ \lim_{x \to 0} \frac{x^2}{x} = 0 \quad ; \quad \lim_{x \to 0} \frac{7x}{x} = 7 \quad ; \quad \lim_{x \to 0} \frac{x}{x^2} = +\infty \ \text{à droite} \quad ; \quad \lim_{x \to 0} \frac{|x|}{x} \ \text{n'existe pas.} \]
La même forme $\dfrac{0}{0}$ donne $0$, $7$, l'infini, ou rien du tout ! Le résultat dépend de \emph{chaque cas particulier} : voilà pourquoi on dit « indéterminée ». Ton travail consiste précisément à déterminer ce cas particulier.
\end{savaistu}

\begin{methode}[Stratégie générale face à une forme indéterminée]
\begin{enumerate}
\item \textbf{Identifier} la forme indéterminée en appliquant naïvement les opérations, et l'\textbf{annoncer} dans la rédaction.
\item \textbf{Transformer} l'expression par une technique algébrique adaptée : factoriser, réduire au même dénominateur, multiplier par la quantité conjuguée, mettre le terme dominant en facteur, ou réécrire un produit comme un quotient.
\item \textbf{Simplifier} puis \textbf{conclure} en appliquant les règles opératoires à la nouvelle expression, qui n'est plus indéterminée.
\end{enumerate}
\end{methode}

\begin{popfigure}
\begin{tikzpicture}[
  node distance=0.55cm and 0.9cm,
  decision/.style={diamond, draw=popdark, fill=popgoldL, aspect=2.2, inner sep=1pt, align=center, font=\small},
  block/.style={rectangle, rounded corners, draw=popdark, fill=popblueL, align=center, font=\small, text width=4.6cm},
  result/.style={rectangle, rounded corners, draw=popdark, fill=popgreenL, align=center, font=\small},
  fleche/.style={-{Stealth[length=2.5mm]}, thick, popdark}
]
\node[decision, align=center] (fi) {Forme\\indéterminée ?};
\node[block, below left=0.9cm and 2.2cm of fi, align=center] (frac) {$\dfrac{0}{0}$ en $x_0$, fraction\\de polynômes :\\ \textbf{factoriser par $(x-x_0)$}\\ puis simplifier};
\node[block, below=0.9cm of fi, align=center] (rac) {$\dfrac{0}{0}$ ou $\infty-\infty$\\avec des racines carrées :\\ \textbf{quantité conjuguée}\\ $(a-b)(a+b)=a^2-b^2$};
\node[block, below right=0.9cm and 2.2cm of fi, align=center] (inf) {$\dfrac{\infty}{\infty}$ ou $\infty-\infty$\\en l'infini :\\ \textbf{mettre le terme dominant}\\ \textbf{en facteur}};
\node[result, below=1.1cm of rac, align=center] (fin) {Appliquer les règles opératoires\\ et conclure};
\draw[fleche] (fi) -- node[above left, font=\small]{fraction rationnelle} (frac);
\draw[fleche] (fi) -- node[right, font=\small]{radicaux} (rac);
\draw[fleche] (fi) -- node[above right, font=\small]{polynômes en $\pm\infty$} (inf);
\draw[fleche] (frac) |- (fin);
\draw[fleche] (rac) -- (fin);
\draw[fleche] (inf) |- (fin);
\end{tikzpicture}
\legende{Organigramme de décision : quelle technique choisir selon la forme de l'expression et la forme indéterminée rencontrée.}
\end{popfigure}

\courssub{Technique 1 : factorisation par $(x - x_0)$ — cas $\dfrac{0}{0}$ des fractions rationnelles}

Commençons par une intuition. Si un polynôme $P$ s'annule en $x_0$, c'est-à-dire si $P(x_0) = 0$, alors $(x - x_0)$ est un facteur de $P$ : c'est le théorème de factorisation vu en Première. On peut donc écrire $P(x) = (x - x_0)\,Q(x)$, où $Q$ est un polynôme de degré inférieur d'une unité. Quand numérateur \emph{et} dénominateur s'annulent en $x_0$, le facteur commun $(x - x_0)$ se simplifie, et l'indétermination disparaît.

\begin{methode}[Lever une F.I. $\dfrac{0}{0}$ en $x_0$ par factorisation]
\begin{enumerate}
\item Vérifier que le numérateur et le dénominateur s'annulent tous les deux en $x_0$.
\item Factoriser chacun par $(x - x_0)$ : pour un trinôme $ax^2+bx+c$ admettant $x_0$ comme racine, on a $ax^2+bx+c = a(x-x_0)(x-x_1)$ où $x_1$ est l'autre racine (trouvée grâce au produit des racines $\dfrac{c}{a}$ ou à la somme $-\dfrac{b}{a}$).
\item Simplifier par $(x - x_0)$ (licite car on étudie la limite pour $x \neq x_0$).
\item Calculer la limite de l'expression simplifiée par substitution directe.
\end{enumerate}
\end{methode}

\begin{exemplebox}[Limite de $\dfrac{x^2-5x+6}{x-2}$ en $2$, rédigée en trois étapes]
\textbf{Étape 1 — Annoncer la F.I.} En $2$ : le numérateur vaut $4 - 10 + 6 = 0$ et le dénominateur vaut $0$. On a une forme indéterminée $\dfrac{0}{0}$.

\textbf{Étape 2 — Factoriser.} $x_0 = 2$ est racine de $x^2 - 5x + 6$. Le produit des racines vaut $6$, donc l'autre racine est $3$ :
\[ x^2 - 5x + 6 = (x-2)(x-3). \]

\textbf{Étape 3 — Simplifier et conclure.} Pour $x \neq 2$ :
\[ \frac{x^2-5x+6}{x-2} = \frac{(x-2)(x-3)}{x-2} = x - 3 \quad \xrightarrow[x \to 2]{} \quad 2 - 3 = -1. \]
Donc $\displaystyle\lim_{x \to 2} \frac{x^2-5x+6}{x-2} = -1$.
\end{exemplebox}

\begin{attention}[On ne simplifie que pour $x \neq x_0$]
La simplification par $(x - x_0)$ est valable parce que, dans un calcul de limite en $x_0$, on considère des valeurs de $x$ \emph{proches} de $x_0$ mais \emph{différentes} de $x_0$. La fonction initiale n'est pas définie en $x_0$ : la fonction simplifiée, elle, l'est souvent. C'est exactement l'idée du \cle{prolongement par continuité} que nous verrons en section 6.
\end{attention}

\courssub{Technique 2 : la quantité conjuguée — expressions avec radicaux}

Quand l'expression contient une différence de racines carrées, la factorisation classique est impossible. L'outil miracle est l'identité remarquable :
\[ (a-b)(a+b) = a^2 - b^2, \]
qui fait \emph{disparaître les racines} car $\left(\sqrt{A}\right)^2 = A$ pour $A \geq 0$.

\begin{definition}[Quantité conjuguée]
La \cle{quantité conjuguée} de $\sqrt{A} - \sqrt{B}$ (avec $A, B \geq 0$) est $\sqrt{A} + \sqrt{B}$, et réciproquement. Multiplier par la quantité conjuguée donne :
\[ \left(\sqrt{A} - \sqrt{B}\right)\left(\sqrt{A} + \sqrt{B}\right) = A - B. \]
De même, le conjugué de $\sqrt{A} - b$ est $\sqrt{A} + b$, et $\left(\sqrt{A}-b\right)\left(\sqrt{A}+b\right) = A - b^2$.
\end{definition}

\begin{methode}[Lever une indétermination avec radicaux par la quantité conjuguée]
\begin{enumerate}
\item Repérer la différence contenant une racine carrée (au numérateur ou au dénominateur).
\item Multiplier le numérateur \textbf{et} le dénominateur par la quantité conjuguée de cette différence.
\item Développer le produit obtenu grâce à $(a-b)(a+b) = a^2 - b^2$ : les racines disparaissent.
\item Simplifier, puis calculer la limite de la nouvelle expression.
\end{enumerate}
\end{methode}

\begin{attention}[L'erreur classique : le conjugué « à moitié »]
L'erreur la plus fréquente au Bac est de multiplier \textbf{uniquement le numérateur} par la quantité conjuguée, en oubliant le dénominateur (ou l'inverse). On change alors la valeur de l'expression : le calcul est faux dès la première ligne. Retiens : on multiplie toujours par $\dfrac{\text{conjugué}}{\text{conjugué}}$, c'est-à-dire par $1$, ce qui ne change rien à la valeur de l'expression. Autre piège : $(\sqrt{A} - b)(\sqrt{A} + b) = A - b^2$, et non $A - b$.
\end{attention}

\begin{exemplebox}[Limite de $\dfrac{\sqrt{x+1}-2}{x-3}$ en $3$]
\textbf{Étape 1 — F.I.} En $3$ : $\sqrt{4} - 2 = 0$ au numérateur et $3 - 3 = 0$ au dénominateur. Forme $\dfrac{0}{0}$.

\textbf{Étape 2 — Quantité conjuguée.} Le conjugué de $\sqrt{x+1} - 2$ est $\sqrt{x+1} + 2$. Pour $x \neq 3$ :
\begin{align*}
\frac{\sqrt{x+1}-2}{x-3} &= \frac{\left(\sqrt{x+1}-2\right)\left(\sqrt{x+1}+2\right)}{(x-3)\left(\sqrt{x+1}+2\right)} \\
&= \frac{(x+1) - 4}{(x-3)\left(\sqrt{x+1}+2\right)} = \frac{x-3}{(x-3)\left(\sqrt{x+1}+2\right)} = \frac{1}{\sqrt{x+1}+2}.
\end{align*}

\textbf{Étape 3 — Conclusion.} $\displaystyle\lim_{x \to 3} \frac{1}{\sqrt{x+1}+2} = \frac{1}{\sqrt{4}+2} = \frac{1}{4}$. Donc $\displaystyle\lim_{x \to 3} \frac{\sqrt{x+1}-2}{x-3} = \frac{1}{4}$.
\end{exemplebox}

\begin{popfigure}
\begin{tikzpicture}
\begin{axis}[
  width=\linewidth, height=7cm,
  axis lines=middle,
  xlabel={$x$}, ylabel={$y$},
  xmin=-1.3, xmax=7, ymin=-0.6, ymax=1.4,
  xtick={-1,1,2,3,4,5,6}, ytick={-0.5,0.25,0.5,1},
  tick label style={font=\small},
  label style={font=\small},
  samples=200,
  unbounded coords=discard
]
\addplot[popcurve, domain=-1:6.5] {(sqrt(max(0,x+1))-2)/(x-3)};
\addplot[popblue, dashed, thin] coordinates {(3,-0.6) (3,1.4)};
\addplot[only marks, mark=o, mark size=2.2pt, popdark, thick] coordinates {(3,0.25)};
\node[font=\small, popdark, anchor=west] at (axis cs:3.15,0.45) {trou en $(3\,;\,0{,}25)$};
\node[font=\small, popblue, anchor=south west] at (axis cs:-1.2,0.9) {$y=\dfrac{\sqrt{x+1}-2}{x-3}$};
\end{axis}
\end{tikzpicture}
\legende{La fonction $f(x)=\dfrac{\sqrt{x+1}-2}{x-3}$ n'est pas définie en $3$ (le point est « creux »), mais la courbe se prolonge naturellement : la limite $\dfrac{1}{4}$ comble le trou. C'est un prolongement par continuité apparent.}
\end{popfigure}

\courssub{Technique 3 : mise en facteur du terme dominant — cas en l'infini}

En $+\infty$ ou $-\infty$, les formes $\dfrac{\infty}{\infty}$ et $\infty - \infty$ sont légion. L'idée directrice : dans une somme de puissances de $x$, c'est la \cle{puissance la plus élevée} qui « écrase » toutes les autres quand $|x|$ devient immense. On la met en facteur pour faire apparaître des termes qui tendent vers $0$.

\begin{propriete}[Terme dominant d'un polynôme ou d'une fraction rationnelle]
En $+\infty$ ou $-\infty$, un polynôme a la même limite que son \cle{terme de plus haut degré}, et une fraction rationnelle a la même limite que le quotient des termes de plus haut degré :
\[ \lim_{x \to \pm\infty} \frac{a_n x^n + \cdots + a_0}{b_p x^p + \cdots + b_0} = \lim_{x \to \pm\infty} \frac{a_n x^n}{b_p x^p} \qquad (a_n \neq 0,\ b_p \neq 0). \]
Ce résultat se démontre précisément par la mise en facteur du terme dominant ; la technique est donc exigible même si le raccourci est admis.
\end{propriete}

\begin{exemplebox}[Mise en facteur du terme dominant]
Calculons $\displaystyle\lim_{x \to +\infty} \frac{3x^2 - 5x + 1}{2x^2 + 7}$. Forme $\dfrac{\infty}{\infty}$. On factorise par $x^2$ en haut et en bas :
\[ \frac{3x^2 - 5x + 1}{2x^2 + 7} = \frac{x^2\left(3 - \dfrac{5}{x} + \dfrac{1}{x^2}\right)}{x^2\left(2 + \dfrac{7}{x^2}\right)} = \frac{3 - \dfrac{5}{x} + \dfrac{1}{x^2}}{2 + \dfrac{7}{x^2}} \xrightarrow[x \to +\infty]{} \frac{3}{2}, \]
car $\dfrac{5}{x}$, $\dfrac{1}{x^2}$ et $\dfrac{7}{x^2}$ tendent vers $0$.
\end{exemplebox}

Pour une différence du type $\sqrt{x^2 + ax} - x$ en $+\infty$ (forme $\infty - \infty$), la mise en facteur seule ne suffit pas : on la combine avec la quantité conjuguée. C'est le grand classique du Baccalauréat.

\begin{exoresolu}[Limite de $\sqrt{x^2 + 3x} - x$ en $+\infty$]
Soit $f(x) = \sqrt{x^2 + 3x} - x$, définie en particulier pour $x \geq 0$. Déterminer $\displaystyle\lim_{x \to +\infty} f(x)$.
\tcblower
\textbf{Étape 1 — Annoncer la F.I.} Quand $x \to +\infty$, $\sqrt{x^2+3x} \to +\infty$ et $x \to +\infty$ : forme indéterminée $\infty - \infty$.

\textbf{Étape 2 — Quantité conjuguée.} Pour $x > 0$, on multiplie et on divise par $\sqrt{x^2+3x} + x$ :
\begin{align*}
\sqrt{x^2+3x} - x &= \frac{\left(\sqrt{x^2+3x} - x\right)\left(\sqrt{x^2+3x} + x\right)}{\sqrt{x^2+3x} + x} \\
&= \frac{(x^2 + 3x) - x^2}{\sqrt{x^2+3x} + x} = \frac{3x}{\sqrt{x^2+3x} + x}.
\end{align*}

\textbf{Étape 3 — Terme dominant.} Pour $x > 0$, on a $\sqrt{x^2} = |x| = x$, donc $\sqrt{x^2+3x} = \sqrt{x^2\left(1+\dfrac{3}{x}\right)} = x\sqrt{1+\dfrac{3}{x}}$. On factorise par $x$ :
\[ \frac{3x}{\sqrt{x^2+3x} + x} = \frac{3x}{x\sqrt{1 + \dfrac{3}{x}} + x} = \frac{3x}{x\left(\sqrt{1 + \dfrac{3}{x}} + 1\right)} = \frac{3}{\sqrt{1 + \dfrac{3}{x}} + 1}. \]

\textbf{Étape 4 — Conclusion.} Comme $\dfrac{3}{x} \to 0$, on obtient $\sqrt{1 + \dfrac{3}{x}} \to 1$, d'où :
\[ \lim_{x \to +\infty} \left(\sqrt{x^2+3x} - x\right) = \frac{3}{1 + 1} = \frac{3}{2}. \]
\end{exoresolu}

\begin{attention}[$\sqrt{x^2} = |x|$, pas toujours $x$ !]
Dans la factorisation ci-dessus, on a utilisé $\sqrt{x^2} = x$ parce que $x > 0$. Si la limite était cherchée en $-\infty$, il faudrait écrire $\sqrt{x^2} = |x| = -x$, ce qui change complètement le calcul. Vérifie toujours le signe de $x$ avant de « sortir » un $x$ d'une racine carrée.
\end{attention}

\courssub{Et la forme $0 \times \infty$ ? On la ramène à un quotient}

La forme $0 \times \infty$ ne possède pas de technique propre : on la \textbf{réécrit} pour se ramener à $\dfrac{0}{0}$ ou $\dfrac{\infty}{\infty}$, en « faisant descendre » l'un des deux facteurs au dénominateur :
\[ f \times g = \frac{f}{1/g} \qquad \text{ou} \qquad f \times g = \frac{g}{1/f}. \]

\begin{exemplebox}[Transformer $0 \times \infty$ en $\dfrac{\infty}{\infty}$]
Calculons $\displaystyle\lim_{x \to +\infty} x\left(\sqrt{1+\tfrac{2}{x}} - 1\right)$. Quand $x \to +\infty$, le facteur $\sqrt{1+\frac{2}{x}} - 1 \to 0$ et $x \to +\infty$ : forme $0 \times \infty$. On écrit $x = \dfrac{1}{1/x}$ pour faire apparaître un quotient, puis on utilise la quantité conjuguée :
\[ x\left(\sqrt{1+\tfrac{2}{x}} - 1\right) = \frac{\sqrt{1+\frac{2}{x}} - 1}{\frac{1}{x}} = \frac{\left(1+\frac{2}{x}\right) - 1}{\frac{1}{x}\left(\sqrt{1+\frac{2}{x}} + 1\right)} = \frac{\frac{2}{x}}{\frac{1}{x}\left(\sqrt{1+\frac{2}{x}} + 1\right)} = \frac{2}{\sqrt{1+\frac{2}{x}} + 1}. \]
Cette dernière expression tend vers $\dfrac{2}{1+1} = 1$. Donc $\displaystyle\lim_{x \to +\infty} x\left(\sqrt{1+\tfrac{2}{x}} - 1\right) = 1$.
\end{exemplebox}

\courssub{Complément utile au Bac et méthode de rédaction}

\begin{propriete}[Limite remarquable (admise)]
\[ \lim_{x \to 0} \frac{\sin x}{x} = 1 \qquad (x \text{ en radians}). \]
Ce résultat est \textbf{admis} au programme de Terminale A. Il permet de traiter certaines formes $\dfrac{0}{0}$ trigonométriques : par exemple $\displaystyle\lim_{x \to 0} \frac{\sin(3x)}{x} = \lim_{x \to 0} 3 \times \frac{\sin(3x)}{3x} = 3 \times 1 = 3$, car $3x \to 0$ quand $x \to 0$.
\end{propriete}

\begin{methode}[Rédaction type d'une levée d'indétermination au Bac]
\begin{enumerate}
\item \textbf{Annoncer} : « En $x_0$, le numérateur tend vers $0$ et le dénominateur tend vers $0$ : forme indéterminée $\dfrac{0}{0}$. »
\item \textbf{Transformer} : « Pour $x \neq x_0$, on a \ldots » puis écrire la suite d'égalités (factorisation, conjugué, etc.). Chaque ligne doit être une égalité exacte.
\item \textbf{Conclure} : « Or $\displaystyle\lim_{x \to x_0} \big(\text{expression simplifiée}\big) = \ell$, donc $\displaystyle\lim_{x \to x_0} f(x) = \ell$. » Encadrer ou souligner le résultat.
\end{enumerate}
Ne jamais écrire $\dfrac{0}{0}$ comme un \emph{résultat} : c'est un diagnostic, pas une réponse.
\end{methode}

\begin{exercice}[À toi de jouer]
\begin{enumerate}
\item Calculer $\displaystyle\lim_{x \to 1} \frac{x^2 + x - 2}{x^2 - 3x + 2}$.
\item Calculer $\displaystyle\lim_{x \to +\infty} \left(\sqrt{x^2 + x} - x\right)$.
\item Calculer $\displaystyle\lim_{x \to 0} \frac{\sqrt{x+4} - 2}{x}$.
\end{enumerate}
\end{exercice}

\begin{corrige}[Exercice : À toi de jouer]
\textbf{1.} En $1$ : numérateur $= 1 + 1 - 2 = 0$, dénominateur $= 1 - 3 + 2 = 0$ : F.I. $\dfrac{0}{0}$. On factorise : $x^2 + x - 2 = (x-1)(x+2)$ (l'autre racine vaut $-2$, car le produit des racines est $-2$) et $x^2 - 3x + 2 = (x-1)(x-2)$ (l'autre racine vaut $2$). Pour $x \neq 1$ :
\[ \frac{x^2+x-2}{x^2-3x+2} = \frac{x+2}{x-2} \xrightarrow[x \to 1]{} \frac{3}{-1} = -3. \]
Donc $\displaystyle\lim_{x \to 1} \frac{x^2 + x - 2}{x^2 - 3x + 2} = -3$.

\textbf{2.} F.I. $\infty - \infty$. Quantité conjuguée : pour $x > 0$,
\[ \sqrt{x^2+x} - x = \frac{(x^2+x) - x^2}{\sqrt{x^2+x} + x} = \frac{x}{\sqrt{x^2+x} + x} = \frac{x}{x\left(\sqrt{1+\frac{1}{x}} + 1\right)} = \frac{1}{\sqrt{1+\frac{1}{x}} + 1} \xrightarrow[x \to +\infty]{} \frac{1}{2}. \]
Donc $\displaystyle\lim_{x \to +\infty} \left(\sqrt{x^2 + x} - x\right) = \dfrac{1}{2}$.

\textbf{3.} F.I. $\dfrac{0}{0}$. Conjugué de $\sqrt{x+4} - 2$ : pour $x \neq 0$ (et $x \geq -4$),
\[ \frac{\sqrt{x+4}-2}{x} = \frac{(x+4)-4}{x\left(\sqrt{x+4}+2\right)} = \frac{x}{x\left(\sqrt{x+4}+2\right)} = \frac{1}{\sqrt{x+4}+2} \xrightarrow[x \to 0]{} \frac{1}{4}. \]
Donc $\displaystyle\lim_{x \to 0} \frac{\sqrt{x+4} - 2}{x} = \dfrac{1}{4}$.
\end{corrige}

\begin{aretenir}[L'essentiel de la section]
\begin{itemize}
\item Une forme indéterminée ($\frac{0}{0}$, $\frac{\infty}{\infty}$, $0 \times \infty$, $\infty - \infty$) se \textbf{lève en transformant l'expression}, jamais en concluant directement.
\item $\frac{0}{0}$ avec des polynômes en $x_0$ : \textbf{factoriser par $(x - x_0)$} et simplifier.
\item Radicaux : \textbf{quantité conjuguée}, en multipliant numérateur \emph{et} dénominateur ; identité $(a-b)(a+b) = a^2 - b^2$.
\item En l'infini : \textbf{mettre le terme dominant en facteur} ; attention à $\sqrt{x^2} = |x|$.
\item $0 \times \infty$ : \textbf{réécrire le produit comme un quotient} pour se ramener à $\frac{0}{0}$ ou $\frac{\infty}{\infty}$.
\item Rédaction en trois temps : annoncer la F.I., transformer par égalités exactes, conclure.
\end{itemize}
\end{aretenir}

\coursec{Continuité en un point et sur un intervalle}

Dans la section précédente, nous avons appris à lever les formes indéterminées et à calculer des limites en un point $x_0$, y compris lorsque la fonction n'est pas définie en ce point. Une question naturelle se pose alors : la limite trouvée correspond-elle à la valeur réelle de la fonction en $x_0$ ? Autrement dit, la courbe « passe-t-elle bien » par le point attendu, ou présente-t-elle une cassure, un trou, un saut ? C'est exactement l'objet de la \cle{continuité}, notion centrale de cette section.

\courssub{Approche intuitive : tracer sans lever le crayon}

Imagine que tu dessines la courbe d'une fonction sur ton cahier. Si tu peux la tracer d'un bout à l'autre \emph{sans jamais lever le crayon}, on dira que la fonction est \cle{continue}. Si tu es obligé de lever le crayon — parce que la courbe saute brusquement, parce qu'il manque un point, ou parce qu'elle part vers l'infini — alors la fonction est \cle{discontinue} à cet endroit.

\begin{exemplebox}[Situations concrètes]
\begin{itemize}
\item Le compteur d'un taximètre à Douala augmente par paliers : le prix affiché \emph{saute} d'un coup à chaque tranche. C'est une fonction \emph{discontinue} du temps.
\item La température à Yaoundé au cours d'une journée varie progressivement, sans saut brutal : c'est une fonction \emph{continue} du temps.
\end{itemize}
\end{exemplebox}

Cette intuition du « crayon qui ne se lève pas » va maintenant être traduite en langage mathématique rigoureux, grâce à la notion de limite étudiée dans les sections précédentes.

\courssub{Continuité d'une fonction en un point}

\begin{definition}[Continuité en un point]
Soit $f$ une fonction définie sur un intervalle contenant $x_0$. On dit que $f$ est \cle{continue en $x_0$} lorsque :
\begin{enumerate}
\item $f$ est définie en $x_0$ (c'est-à-dire $f(x_0)$ existe) ;
\item $\displaystyle\lim_{x \to x_0} f(x) = f(x_0)$.
\end{enumerate}
Dans le cas contraire, on dit que $f$ est \emph{discontinue} en $x_0$.
\end{definition}

Autrement dit, une fonction est continue en $x_0$ lorsque la limite « annoncée » par les valeurs voisines coïncide exactement avec la valeur réelle de la fonction en ce point. Il n'y a ni surprise, ni cassure.

\begin{propriete}[Condition équivalente par les limites latérales]
$f$ est continue en $x_0$ si et seulement si :
\[
\lim_{\substack{x \to x_0 \\ x < x_0}} f(x) \;=\; \lim_{\substack{x \to x_0 \\ x > x_0}} f(x) \;=\; f(x_0).
\]
On dit alors que la \emph{limite à gauche}, la \emph{limite à droite} et la \emph{valeur en $x_0$} sont égales.
\end{propriete}

Cette propriété découle directement du fait qu'une fonction admet une limite en $x_0$ si et seulement si ses limites à gauche et à droite existent et sont égales. Elle est l'outil privilégié pour étudier les fonctions \emph{définies par morceaux}.

\begin{definition}[Continuité à gauche et à droite]
\begin{itemize}
\item $f$ est \cle{continue à gauche} en $x_0$ si $\displaystyle\lim_{\substack{x \to x_0 \\ x < x_0}} f(x) = f(x_0)$ ;
\item $f$ est \cle{continue à droite} en $x_0$ si $\displaystyle\lim_{\substack{x \to x_0 \\ x > x_0}} f(x) = f(x_0)$.
\end{itemize}
$f$ est continue en $x_0$ si et seulement si elle est continue à gauche \emph{et} à droite en $x_0$.
\end{definition}

\begin{attention}[Erreur fréquente : la limite existe, donc c'est continu ?]
Non ! Vérifier que $\displaystyle\lim_{x \to x_0} f(x)$ existe \textbf{ne suffit pas} pour conclure à la continuité. Il faut en plus :
\begin{enumerate}
\item que $f(x_0)$ soit définie ;
\item que la limite soit \emph{égale} à $f(x_0)$.
\end{enumerate}
Par exemple, si $\displaystyle\lim_{x \to 2} f(x) = 5$ mais $f(2) = 3$, alors $f$ est \emph{discontinue} en $2$ : la courbe a un « trou » comblé au mauvais endroit. Rédige toujours la comparaison explicite entre la limite et $f(x_0)$.
\end{attention}

\courssub{Les trois visages de la discontinuité}

\begin{popfigure}
\begin{tikzpicture}
\begin{axis}[width=0.52\linewidth,height=4.5cm,axis lines=middle,xmin=-0.5,xmax=3.5,ymin=-0.5,ymax=4,xtick=\empty,ytick=\empty,title={Continue},title style={font=\small}]
\addplot[popcurve,domain=0:3,samples=100]{0.5*x^2 - x + 1.5};
\end{axis}
\end{tikzpicture}\hfill
\begin{tikzpicture}
\begin{axis}[width=0.52\linewidth,height=4.5cm,axis lines=middle,xmin=-0.5,xmax=3.5,ymin=-0.5,ymax=4,xtick=\empty,ytick=\empty,title={Saut},title style={font=\small}]
\addplot[popcurve,domain=0:1.5,samples=50]{x+0.5};
\addplot[popcurve,domain=1.5:3,samples=50]{x+1.5};
\addplot[mark=*,color=popink] coordinates {(1.5,2)};
\addplot[mark=*,color=popink,fill=white] coordinates {(1.5,3)};
\end{axis}
\end{tikzpicture}

\vspace{0.4cm}

\begin{tikzpicture}
\begin{axis}[width=0.52\linewidth,height=4.5cm,axis lines=middle,xmin=-0.5,xmax=3.5,ymin=-0.5,ymax=4,xtick=\empty,ytick=\empty,title={Trou},title style={font=\small}]
\addplot[popcurve,domain=0:3,samples=100]{x+0.5};
\addplot[mark=*,color=popink,fill=white] coordinates {(1.5,2)};
\end{axis}
\end{tikzpicture}\hfill
\begin{tikzpicture}
\begin{axis}[width=0.52\linewidth,height=4.5cm,axis lines=middle,xmin=-0.5,xmax=3.5,ymin=-0.5,ymax=4,xtick=\empty,ytick=\empty,title={Asymptote verticale},title style={font=\small}]
\addplot[popcurve,domain=0:1.35,samples=100]{1/(1.5-x)+0.2};
\addplot[popcurve,domain=1.65:3,samples=100]{1/(x-1.5)-1.5};
\addplot[dashed,color=popmuted] coordinates {(1.5,-0.5) (1.5,4)};
\end{axis}
\end{tikzpicture}
\legende{Une fonction continue et les trois types classiques de discontinuité en $x_0$ : saut (limites latérales différentes), trou (limite finie mais $f(x_0)$ absente ou différente), asymptote verticale (limite infinie).}
\end{popfigure}

\courssub{Continuité sur un intervalle}

\begin{definition}[Continuité sur un intervalle]
On dit que $f$ est \cle{continue sur un intervalle} $I$ lorsqu'elle est continue en \emph{tout point} de $I$. (Aux bornes éventuellement incluses dans $I$, on exige seulement la continuité à droite ou à gauche selon le cas.)
\end{definition}

Graphiquement, la courbe d'une fonction continue sur $I$ se trace d'un seul tenant, sans lever le crayon, sur tout l'intervalle $I$.

\begin{propriete}[Continuité des fonctions usuelles et stabilité par opérations (admis)]
\begin{enumerate}
\item Les \textbf{fonctions polynômes} sont continues sur $\mathbb{R}$.
\item Les \textbf{fonctions rationnelles} sont continues sur chaque intervalle de leur ensemble de définition.
\item La fonction \textbf{racine carrée} $x \mapsto \sqrt{x}$ est continue sur $[0\,;\,+\infty[$.
\item La fonction \textbf{valeur absolue} $x \mapsto |x|$ est continue sur $\mathbb{R}$.
\item Toute \textbf{somme}, tout \textbf{produit}, tout \textbf{quotient} (là où le dénominateur ne s'annule pas) et toute \textbf{composée} de fonctions continues sont continus.
\end{enumerate}
Ce théorème est admis ; il permet de conclure immédiatement à la continuité de la quasi-totalité des fonctions rencontrées au Bac, \emph{sur leur ensemble de définition}.
\end{propriete}

\begin{exemplebox}[Application directe du théorème]
La fonction $f : x \mapsto \dfrac{x^2 + 3x - 1}{x - 2}$ est une fonction rationnelle, définie sur $\mathbb{R} \setminus \{2\}$. Elle est donc continue sur $]-\infty\,;\,2[$ et sur $]2\,;\,+\infty[$. En $x = 2$, elle n'est pas définie : on ne peut pas parler de continuité en $2$ (et l'on vérifierait que la limite y est infinie : asymptote verticale).
\end{exemplebox}

\courssub{Étude de la continuité d'une fonction définie par morceaux}

C'est la situation la plus fréquente au Baccalauréat. La fonction change d'expression en un point de « jonction » $x_0$ : c'est là, et seulement là, qu'un travail est nécessaire (ailleurs, le théorème précédent s'applique).

\begin{methode}[Continuité d'une fonction définie par morceaux au point de jonction $x_0$]
\begin{enumerate}
\item Calculer $f(x_0)$ à l'aide du morceau qui contient $x_0$.
\item Calculer la limite à gauche $\displaystyle\lim_{\substack{x \to x_0 \\ x < x_0}} f(x)$ avec l'expression valable pour $x < x_0$.
\item Calculer la limite à droite $\displaystyle\lim_{\substack{x \to x_0 \\ x > x_0}} f(x)$ avec l'expression valable pour $x > x_0$.
\item Comparer les trois nombres : s'ils sont égaux, $f$ est continue en $x_0$ ; sinon, elle est discontinue (préciser le type de discontinuité).
\end{enumerate}
\end{methode}

\begin{exoresolu}[Étude de continuité en $1$]
On considère la fonction $f$ définie sur $\mathbb{R}$ par :
\[
f(x) = \begin{cases} x^2 + 1 & \text{si } x \leq 1 \\ 3x - 1 & \text{si } x > 1 \end{cases}
\]
Étudier la continuité de $f$ en $1$.
\tcblower
\textbf{Solution détaillée.}

\textbf{Étape 1 : valeur en $1$.} Le point $x_0 = 1$ appartient au premier morceau ($x \leq 1$) :
\[
f(1) = 1^2 + 1 = 2.
\]

\textbf{Étape 2 : limite à gauche.} Pour $x < 1$, on utilise $f(x) = x^2 + 1$ :
\[
\lim_{\substack{x \to 1 \\ x < 1}} f(x) = \lim_{x \to 1} (x^2 + 1) = 1 + 1 = 2.
\]

\textbf{Étape 3 : limite à droite.} Pour $x > 1$, on utilise $f(x) = 3x - 1$ :
\[
\lim_{\substack{x \to 1 \\ x > 1}} f(x) = \lim_{x \to 1} (3x - 1) = 3 - 1 = 2.
\]

\textbf{Étape 4 : comparaison.} On a :
\[
\lim_{\substack{x \to 1 \\ x < 1}} f(x) = \lim_{\substack{x \to 1 \\ x > 1}} f(x) = f(1) = 2.
\]
Les trois quantités sont égales : \textbf{$f$ est continue en $1$}.

De plus, $f$ est continue sur $]-\infty\,;\,1[$ et sur $]1\,;\,+\infty[$ comme fonction polynôme sur chaque intervalle. Donc $f$ est continue sur $\mathbb{R}$ tout entier.
\end{exoresolu}

\begin{exercice}[À toi de jouer]
Soit $g$ définie sur $\mathbb{R}$ par $g(x) = x + 2$ si $x \leq 2$ et $g(x) = x^2$ si $x > 2$. Étudier la continuité de $g$ en $2$.
\end{exercice}

\begin{corrige}[Exercice 1]
$g(2) = 2 + 2 = 4$. Limite à gauche : $\displaystyle\lim_{\substack{x \to 2 \\ x < 2}} (x+2) = 4$. Limite à droite : $\displaystyle\lim_{\substack{x \to 2 \\ x > 2}} x^2 = 4$. Les trois valeurs sont égales à $4$ : $g$ est continue en $2$.
\end{corrige}

\courssub{Un exemple célèbre de discontinuité : la fonction partie entière}

\begin{definition}[Partie entière]
La \cle{partie entière} d'un réel $x$, notée $E(x)$, est le plus grand entier inférieur ou égal à $x$. Par exemple : $E(2{,}7) = 2$, $E(3) = 3$, $E(-1{,}2) = -2$.
\end{definition}

\begin{popfigure}
\begin{tikzpicture}
\begin{axis}[width=0.75\linewidth,height=5.5cm,axis lines=middle,xmin=-2.5,xmax=3.5,ymin=-2.5,ymax=3.5,xtick={-2,-1,1,2,3},ytick={-2,-1,1,2,3},xlabel={$x$},ylabel={$E(x)$},grid=both,grid style={popline!40}]
\addplot[popcurve,domain=-2:-1,samples=2]{-2};
\addplot[popcurve,domain=-1:0,samples=2]{-1};
\addplot[popcurve,domain=0:1,samples=2]{0};
\addplot[popcurve,domain=1:2,samples=2]{1};
\addplot[popcurve,domain=2:3,samples=2]{2};
\addplot[mark=*,color=popink] coordinates {(-2,-2) (-1,-1) (0,0) (1,1) (2,2) (3,3)};
\addplot[mark=*,color=popink,fill=white] coordinates {(-1,-2) (0,-1) (1,0) (2,1) (3,2)};
\end{axis}
\end{tikzpicture}
\legende{La fonction partie entière : une courbe « en escalier ». En chaque entier $n$, la limite à gauche vaut $n-1$ et la limite à droite vaut $n = E(n)$ : la fonction est continue à droite mais discontinue en tout entier.}
\end{popfigure}

En effet, pour tout entier $n$ : $\displaystyle\lim_{\substack{x \to n \\ x < n}} E(x) = n - 1$ alors que $\displaystyle\lim_{\substack{x \to n \\ x > n}} E(x) = n = E(n)$. Les limites latérales diffèrent : c'est une discontinuité de type \emph{saut}. C'est le modèle mathématique du taximètre ou de la facturation par tranches.

\courssub{Prolongement par continuité (complément utile au Bac)}

Il arrive qu'une fonction ne soit pas définie en $x_0$ mais admette une \emph{limite finie} en $x_0$ — par exemple après avoir levé une forme indéterminée, comme nous l'avons appris dans la section précédente. On peut alors « réparer » la fonction en comblant le trou.

\begin{definition}[Prolongement par continuité]
Soit $f$ une fonction non définie en $x_0$, mais telle que $\displaystyle\lim_{x \to x_0} f(x) = \ell$ avec $\ell$ un réel \emph{fini}. La fonction $\tilde{f}$ définie par :
\[
\tilde{f}(x) = \begin{cases} f(x) & \text{si } x \neq x_0 \\ \ell & \text{si } x = x_0 \end{cases}
\]
est continue en $x_0$. On l'appelle le \cle{prolongement par continuité} de $f$ en $x_0$.
\end{definition}

\begin{exoresolu}[Prolongement par continuité de $f(x) = \dfrac{x^2 - 1}{x - 1}$ en $1$]
Soit $f(x) = \dfrac{x^2 - 1}{x - 1}$, définie sur $\mathbb{R} \setminus \{1\}$. Peut-on prolonger $f$ par continuité en $1$ ?
\tcblower
\textbf{Solution détaillée.} En $x = 1$, on obtient la forme indéterminée $\dfrac{0}{0}$. On factorise le numérateur (identité remarquable) :
\[
f(x) = \dfrac{(x-1)(x+1)}{x-1} = x + 1 \quad \text{pour tout } x \neq 1.
\]
Donc :
\[
\lim_{x \to 1} f(x) = \lim_{x \to 1} (x + 1) = 2.
\]
La limite est finie : on peut prolonger $f$ par continuité en posant $\tilde{f}(1) = 2$. Le prolongement est simplement $\tilde{f}(x) = x + 1$ pour tout $x \in \mathbb{R}$, fonction affine continue sur $\mathbb{R}$. Graphiquement, la courbe de $f$ est la droite d'équation $y = x + 1$ \emph{privée du point} $(1\,;\,2)$ ; le prolongement « bouche le trou ».
\end{exoresolu}

\begin{popfigure}
\begin{tikzpicture}
\begin{axis}[width=0.7\linewidth,height=5.5cm,axis lines=middle,xmin=-1.5,xmax=3.5,ymin=-0.5,ymax=4.5,xtick={-1,1,2,3},ytick={1,2,3,4},xlabel={$x$},ylabel={$y$},grid=both,grid style={popline!40}]
\addplot[popcurve,domain=-1.2:3.2,samples=100]{x+1};
\addplot[mark=*,color=poporange,fill=white,mark size=2.5pt] coordinates {(1,2)};
\end{axis}
\end{tikzpicture}
\legende{La courbe de $f(x) = \dfrac{x^2-1}{x-1}$ est la droite $y = x+1$ avec un trou (point creux orange) en $(1\,;\,2)$. Le prolongement par continuité $\tilde{f}$ comble le trou : $\tilde{f}(1) = 2$.}
\end{popfigure}

\begin{attention}[Conditions du prolongement]
On ne peut prolonger par continuité que si la limite en $x_0$ est \textbf{finie}. Si $\displaystyle\lim_{x \to x_0} f(x) = +\infty$ (asymptote verticale) ou si les limites à gauche et à droite diffèrent (saut), aucun prolongement continu n'est possible. Par exemple, $x \mapsto \dfrac{1}{x}$ ne peut pas être prolongée par continuité en $0$.
\end{attention}

\begin{aretenir}[L'essentiel de la section]
\begin{itemize}
\item $f$ est \cle{continue en $x_0$} si et seulement si $\displaystyle\lim_{x \to x_0} f(x) = f(x_0)$, c'est-à-dire limite à gauche $=$ limite à droite $= f(x_0)$.
\item $f$ est \cle{continue sur un intervalle} si elle est continue en chacun de ses points.
\item Polynômes, fonctions rationnelles (sur leur ensemble de définition), $\sqrt{x}$, $|x|$, ainsi que les sommes, produits, quotients et composées de fonctions continues, sont continus.
\item Trois types de discontinuité : le \emph{saut}, le \emph{trou}, l'\emph{asymptote verticale}.
\item Si la limite en $x_0$ existe et est finie mais $f(x_0)$ n'est pas définie, on peut \cle{prolonger par continuité} en posant $f(x_0) = \ell$.
\end{itemize}
\end{aretenir}

\begin{savaistu}[Pourquoi la continuité est-elle si importante ?]
La continuité est la porte d'entrée de deux théorèmes majeurs que tu rencontreras : le \emph{théorème des valeurs intermédiaires} (une fonction continue qui change de signe s'annule forcément — principe utilisé pour localiser les solutions d'équations) et l'existence de \emph{primitives}. C'est aussi l'hypothèse silencieuse derrière tous les capteurs physiques : quand CAMTEL mesure le débit d'une ligne ou qu'un thermomètre suit la température à Bafoussam, on suppose que les grandeurs varient continûment, sans téléportation d'une valeur à une autre !
\end{savaistu}

\coursec{Interprétations graphiques et synthèse}

Dans la section précédente, tu as appris à reconnaître une fonction continue en un point ou sur un intervalle, et tu as vu qu'une fonction continue se dessine \og sans lever le crayon \fg{}. Nous allons maintenant faire le chemin inverse : partir d'une \emph{courbe} pour y lire des limites et des discontinuités, puis relier chaque résultat de calcul de limite à une conséquence géométrique précise sur la représentation graphique. C'est exactement la démarche attendue au Baccalauréat : un calcul de limite n'est jamais une fin en soi, il se traduit toujours par une propriété de la courbe.

\courssub{Lecture graphique des limites : asymptotes horizontales et verticales}

\textbf{Intuition.} Imagine que tu suis la courbe d'une fonction avec ton stylo, en allant de plus en plus loin vers la droite. Si la courbe finit par \og coller \fg{} à une droite horizontale, c'est que la fonction a une limite finie en $+\infty$. Si au contraire, en s'approchant d'une certaine abscisse $x_0$, la courbe monte ou descend indéfiniment, c'est que la fonction a une limite infinie en $x_0$. Ces droites qui \og attirent \fg{} la courbe sans jamais être atteintes s'appellent des \cle{asymptotes}.

\begin{definition}[Asymptote horizontale]
Soit $f$ une fonction et $\mathcal{C}_f$ sa courbe représentative. Si $\lim\limits_{x\to +\infty} f(x) = \ell$ avec $\ell \in \mathbb{R}$ (respectivement en $-\infty$), alors la droite d'équation $y = \ell$ est une \cle{asymptote horizontale} à $\mathcal{C}_f$ en $+\infty$ (respectivement en $-\infty$).
\end{definition}

\begin{definition}[Asymptote verticale]
Si $\lim\limits_{x\to x_0} f(x) = +\infty$ ou $-\infty$ (ou si l'une des limites à droite ou à gauche en $x_0$ est infinie), alors la droite d'équation $x = x_0$ est une \cle{asymptote verticale} à $\mathcal{C}_f$.
\end{definition}

\begin{propriete}[Lien entre limites et asymptotes]
Soit $f$ une fonction de courbe représentative $\mathcal{C}_f$.
\begin{itemize}
\item $\lim\limits_{x\to x_0} f(x) = \pm\infty$ (ou limite infinie à droite ou à gauche) $\iff$ la droite $x = x_0$ est asymptote verticale à $\mathcal{C}_f$.
\item $\lim\limits_{x\to \pm\infty} f(x) = \ell \in \mathbb{R}$ $\iff$ la droite $y = \ell$ est asymptote horizontale à $\mathcal{C}_f$ en $\pm\infty$.
\item \emph{Complément utile au Bac :} si $\lim\limits_{x\to \pm\infty} \big[f(x) - (ax+b)\big] = 0$, alors la droite d'équation $y = ax+b$ est une \cle{asymptote oblique} à $\mathcal{C}_f$ en $\pm\infty$.
\end{itemize}
\end{propriete}

\begin{attention}[Le trou n'est pas une asymptote]
Erreur très fréquente au Bac : affirmer que la droite $x = x_0$ est asymptote verticale alors que la limite de $f$ en $x_0$ est \emph{finie}. Si $\lim\limits_{x\to x_0} f(x) = \ell$ avec $\ell$ réel, la courbe ne s'envole pas : elle présente simplement un \cle{trou} au point $(x_0\,;\,\ell)$ lorsque $f$ n'est pas définie en $x_0$. Avant de conclure à une asymptote verticale, vérifie toujours que la limite est bien \emph{infinie}. C'est le cas du facteur commun qui se simplifie : la forme indéterminée cache souvent un trou, pas une asymptote.
\end{attention}

\courssub{Exemple chiffré : étude complète d'une fonction rationnelle}

\begin{exemplebox}[Limites et asymptotes de $f(x)=\dfrac{x^2-3x+2}{x-1}$]
Soit $f(x) = \dfrac{x^2-3x+2}{x-1}$, définie sur $\mathbb{R}\setminus\{1\}$.

\textbf{Étude en $x_0 = 1$.} On a une forme $\dfrac{0}{0}$ : c'est une forme indéterminée. On factorise le numérateur : $x^2-3x+2 = (x-1)(x-2)$. Donc pour $x \neq 1$ :
\[ f(x) = \frac{(x-1)(x-2)}{x-1} = x-2 \quad\text{et}\quad \lim_{x\to 1} f(x) = \lim_{x\to 1} (x-2) = -1. \]
La limite est \emph{finie} : il n'y a \textbf{pas} d'asymptote verticale en $x=1$, mais un \textbf{trou} au point $(1\,;\,-1)$.

\textbf{Étude en l'infini.} Pour $x \neq 1$, $f(x) = x-2$, donc $\lim\limits_{x\to +\infty} f(x) = +\infty$ et $\lim\limits_{x\to -\infty} f(x) = -\infty$. Pas d'asymptote horizontale ; la courbe est la droite $y = x-2$ privée du point $(1\,;\,-1)$.

\textbf{Conclusion graphique.} La courbe de $f$ est une droite trouée : aucune asymptote, un seul point manquant. Ce cas d'école illustre parfaitement le piège signalé plus haut.
\end{exemplebox}

Étudions maintenant un cas où les asymptotes existent réellement.

\begin{popfigure}
\begin{tikzpicture}
\begin{axis}[
  width=\linewidth, height=7.5cm,
  axis lines=middle, xlabel=$x$, ylabel=$y$,
  xmin=-6, xmax=8, ymin=-5, ymax=9,
  xtick={-4,-2,2,4,6}, ytick={-4,-2,2,4,6,8},
  domain=-6:8, samples=200, clip=true]
\addplot[popcurve, domain=-6:1.6] {(2*x+1)/(x-2)};
\addplot[popcurve, domain=2.4:8] {(2*x+1)/(x-2)};
\addplot[poporange, dashed, thick] coordinates {(2,-5) (2,9)};
\addplot[popgreen, dashed, thick] coordinates {(-6,2) (8,2)};
\node[poporange, anchor=west] at (axis cs:2.2,7.5) {$x=2$};
\node[popgreen, anchor=south west] at (axis cs:-5.8,2.1) {$y=2$};
\end{axis}
\end{tikzpicture}
\legende{Courbe de $g(x)=\dfrac{2x+1}{x-2}$ : asymptote verticale $x=2$ (car $\lim\limits_{x\to 2^+}g(x)=+\infty$ et $\lim\limits_{x\to 2^-}g(x)=-\infty$) et asymptote horizontale $y=2$ (car $\lim\limits_{x\to\pm\infty}g(x)=2$).}
\end{popfigure}

Sur cette figure, on \emph{lit} directement : quand $x$ s'approche de $2$ par la droite, la courbe monte vers $+\infty$ ; quand $x$ grandit, la courbe se rapproche de la droite $y=2$. Le calcul de limite et la lecture graphique se confirment mutuellement : c'est cette double compétence que le correcteur du Bac attend.

\courssub{Lecture graphique de la continuité et des discontinuités}

\textbf{Intuition.} Une fonction continue sur un intervalle a une courbe \og d'un seul morceau \fg{} sur cet intervalle. À l'inverse, une discontinuité se voit : la courbe est \emph{cassée}. Mais toutes les cassures ne se ressemblent pas.

\begin{definition}[Nature des discontinuités]
Soit $f$ définie au voisinage de $x_0$.
\begin{itemize}
\item \cle{Discontinuité de première espèce} (ou saut) : les limites à gauche et à droite en $x_0$ existent, sont finies, mais différentes. La courbe fait un \og saut \fg{}.
\item \cle{Discontinuité de seconde espèce} : au moins l'une des limites à gauche ou à droite est infinie ou n'existe pas. Cas particulier fréquent : l'asymptote verticale.
\item \cle{Discontinuité par trou} (prolongeable par continuité) : la limite en $x_0$ existe et est finie, mais $f$ n'est pas définie en $x_0$ (ou $f(x_0)$ diffère de cette limite).
\end{itemize}
\end{definition}

\begin{exemplebox}[Identifier les discontinuités sur une courbe]
Considère la fonction partie entière $E(x)$ (le plus grand entier inférieur ou égal à $x$). En $x_0 = 2$ : $\lim\limits_{x\to 2^-} E(x) = 1$ et $\lim\limits_{x\to 2^+} E(x) = 2$. Les deux limites sont finies mais distinctes : discontinuité de première espèce (saut de $1$). En revanche, $g(x)=\dfrac{2x+1}{x-2}$ a en $2$ une discontinuité de seconde espèce (limites infinies, asymptote verticale), et $f(x)=\dfrac{x^2-3x+2}{x-1}$ a en $1$ une discontinuité par trou, \emph{prolongeable par continuité} en posant $\tilde{f}(1) = -1$.
\end{exemplebox}

\begin{attention}[Continuité et valeur en un point]
Pour conclure que $f$ est continue en $x_0$, il faut vérifier \textbf{trois} choses : $f(x_0)$ existe, $\lim\limits_{x\to x_0} f(x)$ existe, et les deux sont \emph{égales}. Oublier de comparer la limite à $f(x_0)$ est une erreur classique : une fonction définie par morceaux peut avoir une limite en $x_0$ différente de $f(x_0)$ (point isolé décalé).
\end{attention}

\courssub{Complément Bac : l'asymptote oblique}

Certaines fonctions n'ont pas de limite finie en l'infini, mais leur courbe se rapproche d'une droite inclinée. C'est le cas des fonctions de la forme $f(x) = ax + b + \dfrac{\varphi(x)}{1}$ où $\varphi(x) \to 0$.

\begin{propriete}[Asymptote oblique — complément utile au Bac]
La droite d'équation $y = ax+b$ est asymptote oblique à $\mathcal{C}_f$ en $+\infty$ (ou $-\infty$) si et seulement si :
\[ \lim_{x\to \pm\infty} \big[f(x) - (ax+b)\big] = 0. \]
Le signe de $f(x)-(ax+b)$ donne en outre la \emph{position relative} de la courbe par rapport à son asymptote.
\end{propriete}

\begin{exemplebox}[Mise en évidence d'une asymptote oblique]
Soit $h(x) = x + 1 + \dfrac{2}{x}$ pour $x > 0$. On calcule :
\[ h(x) - (x+1) = \frac{2}{x} \xrightarrow[x\to +\infty]{} 0. \]
Donc la droite $y = x+1$ est asymptote oblique à $\mathcal{C}_h$ en $+\infty$. Comme $\dfrac{2}{x} > 0$ pour $x>0$, la courbe est \emph{au-dessus} de son asymptote. De plus $\lim\limits_{x\to 0^+} h(x) = +\infty$ : la droite $x=0$ (l'axe des ordonnées) est asymptote verticale.
\end{exemplebox}

\begin{popfigure}
\begin{tikzpicture}
\begin{axis}[
  width=\linewidth, height=7.5cm,
  axis lines=middle, xlabel=$x$, ylabel=$y$,
  xmin=-0.5, xmax=8, ymin=-0.5, ymax=10,
  xtick={2,4,6}, ytick={2,4,6,8},
  clip=true]
\addplot[popcurve, domain=0.22:8, samples=200] {x+1+2/x};
\addplot[poporange, dashed, thick, domain=0:8] {x+1};
\node[poporange, anchor=south west, rotate=45] at (axis cs:3,4) {$y=x+1$};
\node[popblue, anchor=west] at (axis cs:3.2,7.6) {$\mathcal{C}_h$};
\end{axis}
\end{tikzpicture}
\legende{Courbe de $h(x)=x+1+\dfrac{2}{x}$ et son asymptote oblique $y=x+1$ : l'écart vertical $\frac{2}{x}$ tend vers $0$ quand $x\to+\infty$.}
\end{popfigure}

\begin{savaistu}[Les asymptotes obliques en économie]
En économie, le \emph{coût moyen de production} d'une entreprise est souvent modélisé par une fonction du type $C_M(q) = aq + b + \dfrac{c}{q}$, où $q$ est la quantité produite et $\dfrac{c}{q}$ représente les coûts fixes répartis. Quand la production $q$ devient très grande (production de masse, par exemple dans une cimenterie ou une brasserie camerounaise), le coût moyen se rapproche de la droite $y = aq+b$ : l'asymptote oblique représente le \emph{coût moyen à long terme}, vers lequel l'entreprise tend sans jamais l'atteindre exactement. Les mathématiques que tu apprends aujourd'hui servent à optimiser de vraies décisions industrielles !
\end{savaistu}

\courssub{Méthode : tableau des limites aux bornes et déduction des asymptotes}

\begin{methode}[Dresser le bilan des limites et des asymptotes]
Pour une fonction $f$ d'ensemble de définition $\mathcal{D}_f$ :
\begin{enumerate}
\item \textbf{Repérer les bornes} de $\mathcal{D}_f$ : $-\infty$, $+\infty$, et chaque valeur interdite $x_0$ (dénominateur nul, racine d'une quantité négative, etc.).
\item \textbf{Calculer la limite en chaque borne} : en une valeur interdite, calculer les limites à gauche \emph{et} à droite séparément.
\item \textbf{Interpréter chaque résultat} : limite infinie en $x_0$ $\Rightarrow$ asymptote verticale $x=x_0$ ; limite finie $\ell$ en $\pm\infty$ $\Rightarrow$ asymptote horizontale $y=\ell$ ; limite infinie en l'infini $\Rightarrow$ chercher éventuellement une asymptote oblique (complément) ; limite finie en $x_0$ avec $f$ non définie $\Rightarrow$ trou (prolongement par continuité).
\item \textbf{Rédiger la conclusion} en reliant explicitement le calcul à la conséquence graphique : \og donc la droite d'équation \ldots est asymptote à $\mathcal{C}_f$ \fg{}.
\end{enumerate}
\end{methode}

\begin{aretenir}[Tableau bilan : limite, continuité, interprétation graphique]
\begin{center}
\renewcommand{\arraystretch}{1.4}
\begin{tabularx}{\linewidth}{|l|X|}
\hline
\rowcolor{popblueL} \textbf{Résultat du calcul} & \textbf{Conséquence graphique} \\
\hline
$\lim\limits_{x\to x_0} f(x) = \pm\infty$ & Asymptote verticale $x = x_0$ ; discontinuité de seconde espèce \\
\hline
$\lim\limits_{x\to x_0} f(x) = \ell \in \mathbb{R}$, $f$ non définie en $x_0$ & Trou au point $(x_0\,;\,\ell)$ ; prolongement par continuité possible \\
\hline
$\lim\limits_{x\to x_0^-} f(x) \neq \lim\limits_{x\to x_0^+} f(x)$ (finies) & Saut : discontinuité de première espèce \\
\hline
$\lim\limits_{x\to x_0} f(x) = f(x_0)$ & $f$ continue en $x_0$ : courbe sans cassure en ce point \\
\hline
$\lim\limits_{x\to \pm\infty} f(x) = \ell \in \mathbb{R}$ & Asymptote horizontale $y = \ell$ en $\pm\infty$ \\
\hline
$\lim\limits_{x\to \pm\infty} f(x) = \pm\infty$ & Pas d'asymptote horizontale ; branche infinie (chercher éventuellement une asymptote oblique) \\
\hline
$\lim\limits_{x\to \pm\infty} [f(x)-(ax+b)] = 0$ & Asymptote oblique $y = ax+b$ (complément Bac) \\
\hline
\end{tabularx}
\end{center}
\end{aretenir}

\courssub{Entraînement à la rédaction type Bac}

\begin{exoresolu}[Limites aux bornes, asymptotes et continuité]
Soit $f$ la fonction définie par $f(x) = \dfrac{x^2 + x - 2}{x + 2}$.
\begin{enumerate}
\item Déterminer l'ensemble de définition $\mathcal{D}_f$ de $f$.
\item Calculer les limites de $f$ aux bornes de $\mathcal{D}_f$ et en déduire les asymptotes éventuelles à la courbe $\mathcal{C}_f$.
\item $f$ est-elle continue en $-2$ ? Justifier.
\item (Complément) Montrer que la droite $y = x - 1$ est asymptote oblique à $\mathcal{C}_f$.
\end{enumerate}
\tcblower
\textbf{1. Ensemble de définition.} $f$ est définie si et seulement si $x+2 \neq 0$, donc $\mathcal{D}_f = \mathbb{R}\setminus\{-2\} = \left]-\infty\,;\,-2\right[ \cup \left]-2\,;\,+\infty\right[$.

\textbf{2. Limites aux bornes.} Les bornes sont $-\infty$, $-2$ (à gauche et à droite) et $+\infty$.

\emph{En $-2$ :} le numérateur vaut $(-2)^2 + (-2) - 2 = 0$ en $-2$ : forme indéterminée $\dfrac{0}{0}$. On factorise : $x^2+x-2 = (x+2)(x-1)$. Donc pour $x \neq -2$ :
\[ f(x) = \frac{(x+2)(x-1)}{x+2} = x-1 \quad\text{et}\quad \lim_{x\to -2} f(x) = \lim_{x\to -2} (x-1) = -3. \]
La limite est \emph{finie} : il n'y a \textbf{pas d'asymptote verticale} en $-2$, mais un \textbf{trou} au point $(-2\,;\,-3)$.

\emph{En l'infini :} pour $x \neq -2$, $f(x) = x-1$, donc $\lim\limits_{x\to +\infty} f(x) = +\infty$ et $\lim\limits_{x\to -\infty} f(x) = -\infty$. Pas d'asymptote horizontale.

\textbf{3. Continuité en $-2$.} $f$ n'est \emph{pas définie} en $-2$, donc $f$ n'est pas continue en $-2$. Cependant, comme $\lim\limits_{x\to -2} f(x) = -3$ existe et est finie, $f$ admet un \textbf{prolongement par continuité} en $-2$, défini par $\tilde{f}(-2) = -3$. Sur $\mathcal{D}_f$, $f$ est continue comme quotient de polynômes dont le dénominateur ne s'annule pas.

\textbf{4. Asymptote oblique.} Pour $x \neq -2$, $f(x) - (x-1) = 0$, donc $\lim\limits_{x\to \pm\infty} [f(x)-(x-1)] = 0$ : la droite $y = x-1$ est asymptote oblique à $\mathcal{C}_f$ en $-\infty$ et en $+\infty$ (ici, la courbe est même \emph{confondue} avec cette droite, au trou près).
\end{exoresolu}

\begin{exercice}[Pour t'entraîner]
Soit $g(x) = \dfrac{3x+1}{x-1}$ et $h(x) = \dfrac{x^2-4}{x-2}$.
\begin{enumerate}
\item Détermine les ensembles de définition de $g$ et $h$.
\item Calcule les limites aux bornes de chaque fonction et déduis-en toutes les asymptotes (ou les trous) des courbes.
\item $h$ admet-elle un prolongement par continuité en $2$ ? Si oui, le définir.
\end{enumerate}
\end{exercice}

\begin{corrige}[Exercice : éléments de correction]
\textbf{1.} $\mathcal{D}_g = \mathbb{R}\setminus\{1\}$ ; $\mathcal{D}_h = \mathbb{R}\setminus\{2\}$.

\textbf{2.} Pour $g$ : $\lim\limits_{x\to 1^-} g(x) = -\infty$ et $\lim\limits_{x\to 1^+} g(x) = +\infty$ (numérateur tend vers $4>0$, dénominateur tend vers $0$ en changeant de signe) : asymptote verticale $x=1$. $\lim\limits_{x\to \pm\infty} g(x) = 3$ : asymptote horizontale $y=3$.

Pour $h$ : $x^2-4 = (x-2)(x+2)$, donc $h(x) = x+2$ pour $x\neq 2$ et $\lim\limits_{x\to 2} h(x) = 4$ : pas d'asymptote verticale, mais un trou en $(2\,;\,4)$. $\lim\limits_{x\to\pm\infty} h(x) = \pm\infty$ : pas d'asymptote horizontale.

\textbf{3.} Oui : le prolongement par continuité de $h$ en $2$ est $\tilde{h}(2) = 4$.
\end{corrige}

\begin{aretenir}[L'essentiel de la leçon]
\begin{itemize}
\item Toute limite calculée se traduit géométriquement : limite infinie en un point $\Rightarrow$ asymptote verticale ; limite finie en l'infini $\Rightarrow$ asymptote horizontale ; $f(x)-(ax+b)\to 0$ $\Rightarrow$ asymptote oblique.
\item Une limite finie en une valeur interdite signale un \emph{trou}, jamais une asymptote verticale.
\item $f$ continue en $x_0$ $\iff$ $\lim\limits_{x\to x_0} f(x) = f(x_0)$ : graphiquement, la courbe n'a ni trou ni saut en ce point.
\item Au Bac, conclus toujours chaque calcul de limite par son interprétation graphique rédigée en une phrase complète.
\end{itemize}
\end{aretenir}

\newpage

\coursec{Activité d'intégration}

\coursec{Limites et continuité}

\courssub{Synthèse théorique : rappels et outils fondamentaux}

\begin{definition}[Limite d'une fonction en un point]
Soit $f$ une fonction définie sur un intervalle $I$ (sauf éventuellement en $a$) et $a \in \mathbb{R}$ une valeur d'adhérence de $I$. On dit que $f$ admet la limite $\ell \in \mathbb{R}$ en $a$ si, pour tout voisinage $V$ de $\ell$, il existe un voisinage $U$ de $a$ tel que pour tout $x \in U \cap I$, $x \neq a$, on a $f(x) \in V$. On note :
\[ \lim_{x \to a} f(x) = \ell. \]
On définit de même les \cle{limites à droite} ($x \to a^+$, $x > a$) et \cle{à gauche} ($x \to a^-$, $x < a$). La limite en $a$ existe si et seulement si les limites à droite et à gauche existent et sont égales.
\end{definition}

\begin{definition}[Limite en $+\infty$ (ou $-\infty$)]
Soit $f$ définie sur un intervalle de la forme $]A, +\infty[$. On dit que $f$ admet la limite $\ell \in \mathbb{R}$ en $+\infty$ si, pour tout $\varepsilon > 0$, il existe $M > 0$ tel que pour tout $x > M$, $|f(x) - \ell| < \varepsilon$. On note :
\[ \lim_{x \to +\infty} f(x) = \ell. \]
Les limites infinies ($\lim f(x) = +\infty$ ou $-\infty$) se définissent par des inégalités du type $f(x) > M$ (ou $< -M$).
\end{definition}

\begin{propriete}[Limites usuelles (admises)]
Pour tout $\alpha > 0$ :
\[ \lim_{x \to +\infty} \frac{1}{x^\alpha} = 0, \qquad \lim_{\substack{x \to 0 \\ x > 0}} \frac{1}{x^\alpha} = +\infty. \]
Pour toute fonction polynôme $P$ de degré $n$ et coefficient dominant $a_n \neq 0$ :
\[ \lim_{x \to +\infty} P(x) = \operatorname{sgn}(a_n) \cdot (+\infty), \qquad \lim_{x \to +\infty} \frac{P(x)}{Q(x)} = \begin{cases} 0 & \text{si } \deg P < \deg Q \\ \frac{a_n}{b_m} & \text{si } \deg P = \deg Q \\ \operatorname{sgn}(a_n/b_m) \cdot (+\infty) & \text{si } \deg P > \deg Q \end{cases} \]
\end{propriete}

\begin{propriete}[Opérations sur les limites]
Soient $f$ et $g$ deux fonctions admettant des limites en $a$ (ou $+\infty$).
\begin{itemize}
\item \textbf{Somme :} $\lim (f+g) = \lim f + \lim g$ (si pas de forme $+\infty - \infty$).
\item \textbf{Produit :} $\lim (f \times g) = \lim f \times \lim g$ (si pas de forme $0 \times \infty$).
\item \textbf{Quotient :} $\lim \frac{f}{g} = \frac{\lim f}{\lim g}$ si $\lim g \neq 0$ (si $\lim g = 0$ et $\lim f \neq 0$, limite infinie ; si $\lim f = \lim g = 0$, forme indéterminée $\frac{0}{0}$).
\item \textbf{Composée :} Si $\lim_{x \to a} g(x) = \ell$ et $f$ continue en $\ell$, alors $\lim_{x \to a} f(g(x)) = f(\ell)$.
\end{itemize}
Les cas exclus (\emph{formes indéterminées}) sont : $\dfrac{0}{0}$, $\dfrac{\infty}{\infty}$, $0 \times \infty$, $+\infty - \infty$, $1^\infty$, $0^0$, $\infty^0$. Ils nécessitent une transformation de l'expression.
\end{propriete}

\begin{methode}[Lever une forme indéterminée $\dfrac{0}{0}$]
\begin{enumerate}
\item Identifier la forme indéterminée (numérateur $\to 0$, dénominateur $\to 0$).
\item \textbf{Factorisation :} Si $f$ et $g$ sont des polynômes ou expressions factorisables, factoriser par $(x-a)$ (ou l'expression qui s'annule) et simplifier.
\item \textbf{Quantité conjuguée :} Si l'expression contient une différence de racines carrées $\sqrt{u(x)} - \sqrt{v(x)}$, multiplier numérateur et dénominateur par $\sqrt{u(x)} + \sqrt{v(x)}$ en utilisant l'identité $(\sqrt{A}-\sqrt{B})(\sqrt{A}+\sqrt{B}) = A-B$.
\item Calculer la limite de l'expression simplifiée (qui n'est plus indéterminée).
\item Conclure.
\end{enumerate}
\end{methode}

\begin{definition}[Continuité en un point et sur un intervalle]
Une fonction $f$ est \cle{continue en $a$} (appartenant à son ensemble de définition) si :
\[ \lim_{x \to a} f(x) = f(a). \]
Cela implique trois conditions : $f$ définie en $a$, limite existante (finie), égalité.
$f$ est \cle{continue sur un intervalle $I$} si elle est continue en tout point de $I$.
Une fonction définie par morceaux est continue aux points de raccordement si les limites à gauche et à droite sont égales à la valeur de la fonction.
\end{definition}

\begin{propriete}[Prolongement par continuité]
Si $f$ n'est pas définie en $a$ mais admet une limite finie $\ell$ en $a$, on peut définir une nouvelle fonction $\tilde{f}$ sur $I \cup \{a\}$ par $\tilde{f}(x) = f(x)$ pour $x \neq a$ et $\tilde{f}(a) = \ell$. La fonction $\tilde{f}$ est continue en $a$ ; c'est le \cle{prolongement par continuité} de $f$ en $a$.
\end{propriete}

\begin{aretenir}[Interprétations graphiques]
\begin{itemize}
\item $\lim_{x \to +\infty} f(x) = \ell \in \mathbb{R}$ $\iff$ la courbe $\mathcal{C}_f$ admet la droite $y = \ell$ comme \cle{asymptote horizontale} en $+\infty$.
\item $\lim_{x \to a} f(x) = +\infty$ $\iff$ la courbe $\mathcal{C}_f$ admet la droite $x = a$ comme \cle{asymptote verticale} (branche montante).
\item $f$ non continue en $a$ (limites latérales finies et différentes) $\iff$ \cle{saut de discontinuité} (ou discontinuité de première espèce) : la courbe présente un « trou » et un « bond ».
\end{itemize}
\end{aretenir}

\courssub{Activité d'intégration : Le forage de la coopérative de Pitoa}

\courssub{Objectifs de l'activité}

À l'issue de cette activité d'intégration, tu seras capable de :

\begin{itemize}
\item \cle{calculer} la limite d'une fonction en un point et en l'infini en utilisant les opérations sur les limites et les limites usuelles ;
\item \cle{lever} une forme indéterminée du type $\dfrac{0}{0}$ par la technique de la quantité conjuguée ;
\item \cle{prolonger} une fonction par continuité en un point ;
\item \cle{étudier} la continuité d'une fonction définie par morceaux en un point de raccordement ;
\item \cle{interpréter} graphiquement et concrètement une limite (asymptote horizontale, discontinuité) dans une situation économique réelle ;
\item \cle{rédiger} un court rapport argumenté à destination de décideurs non mathématiciens.
\end{itemize}

\courssub{Situation}

Dans la région du Nord, aux environs de Garoua, la \textbf{SODECOTON} (Société de Développement du Coton) accompagne les coopératives agricoles dans l'accès à l'eau, indispensable à l'irrigation pendant la longue saison sèche. La coopérative cotonnière de Pitoa souhaite faire forer un puits et mandate une entreprise spécialisée.

L'entreprise propose la tarification suivante : le coût total de forage, exprimé en \textbf{milliers de FCFA} (kFCFA), pour une profondeur de $x$ mètres ($x > 0$) est modélisé par :
\[ C(x) = 50x + 200. \]
Le terme $200$ correspond aux frais fixes (déplacement du matériel, installation de la foreuse) et le terme $50x$ au coût proportionnel à la profondeur (soit $50\,000$ FCFA par mètre foré).

Le trésorier de la coopérative s'intéresse au \textbf{coût moyen par mètre foré}, défini par :
\[ M(x) = \frac{C(x)}{x} = \frac{50x + 200}{x}, \qquad x > 0. \]

Par ailleurs, le sous-sol de la zone contient une couche rocheuse. L'entreprise prévoit un \textbf{supplément de pénalité} (en milliers de FCFA par mètre) lié à la dureté du sol, modélisé pour $x > 0$ par :
\[ P(x) = \frac{\sqrt{x+9} - 3}{x}. \]

Enfin, pour fidéliser ses gros clients, l'entreprise accorde une \textbf{remise} (en milliers de FCFA) qui dépend de la profondeur forée, avec un palier à $50$ mètres :
\[ R(x) = \begin{cases} \dfrac{x}{5} & \text{si } 0 < x < 50, \\[6pt] \dfrac{x}{5} + 8 & \text{si } x \geq 50. \end{cases} \]

La coopérative hésite : faut-il forer un puits peu profond (autour de $45$ m) ou franchir le palier des $50$ m ? Le coût moyen devient-il raisonnable si l'on fore très profond ? Que vaut la pénalité au tout début du forage ? Tu es le(la) conseiller(ère) technique de la coopérative : à toi de trancher avec les mathématiques.

\courssub{Consignes}

Travail en groupes de trois ou quatre élèves. Durée conseillée : $1$ h $30$, suivies d'une restitution orale.

\begin{enumerate}
\item \textbf{Coût moyen en grande profondeur.} Calculer $\displaystyle\lim_{x \to +\infty} M(x)$. Interpréter concrètement le résultat : que devient le coût moyen par mètre lorsque le puits est très profond ? Quelle droite est asymptote à la courbe de $M$ ?
\item \textbf{Coût moyen au voisinage de zéro.} Calculer $\displaystyle\lim_{\substack{x \to 0 \\ x > 0}} M(x)$. Expliquer pourquoi ce résultat signifie qu'un forage de très faible profondeur est économiquement absurde.
\item \textbf{La pénalité de roche.} Montrer que le calcul direct de la limite de $P(x)$ en $0$ conduit à une forme indéterminée. En multipliant numérateur et dénominateur par la quantité conjuguée $\sqrt{x+9}+3$, lever cette indétermination et déterminer $\displaystyle\lim_{x \to 0} P(x)$. La fonction $P$ est-elle prolongeable par continuité en $0$ ? Si oui, donner ce prolongement.
\item \textbf{Le palier de remise.} Calculer $\displaystyle\lim_{\substack{x \to 50 \\ x < 50}} R(x)$ et $\displaystyle\lim_{\substack{x \to 50 \\ x > 50}} R(x)$, ainsi que $R(50)$. La fonction $R$ est-elle continue en $50$ ? Interpréter le « saut » observé pour la coopérative.
\item \textbf{Représentation graphique.} Sur papier millimétré (ou à l'aide d'un logiciel), tracer dans un même repère les courbes de $M$ et de $R$ pour $x \in [1\,;\,100]$. Faire apparaître l'asymptote horizontale de $M$ et le point de discontinuité de $R$.
\item \textbf{Rapport de conseil.} Rédiger en cinq à huit lignes un rapport à l'intention du président de la coopérative, répondant à la question : « Vaut-il mieux forer juste en dessous ou juste au-dessus de $50$ mètres ? », en mobilisant au moins deux résultats de limites obtenus ci-dessus.
\end{enumerate}

\courssub{Ressources à mobiliser}

\begin{aretenir}[Boîte à outils de l'activité]
\begin{itemize}
\item \textbf{Limites usuelles :} $\displaystyle\lim_{x \to +\infty} \frac{1}{x} = 0$ ; $\displaystyle\lim_{\substack{x \to 0 \\ x > 0}} \frac{1}{x} = +\infty$.
\item \textbf{Opérations :} la limite d'une somme, d'un produit ou d'un quotient (lorsque le dénominateur a une limite non nulle) est la somme, le produit ou le quotient des limites.
\item \textbf{Quantité conjuguée :} $(\sqrt{a}-b)(\sqrt{a}+b) = a - b^2$. Utile pour lever une forme $\dfrac{0}{0}$ avec des racines carrées.
\item \textbf{Continuité en un point :} $f$ est continue en $a$ si et seulement si $\displaystyle\lim_{x \to a} f(x) = f(a)$, ce qui exige l'égalité des limites à gauche et à droite.
\item \textbf{Prolongement par continuité :} si $\displaystyle\lim_{x \to a} f(x) = \ell$ (finie) alors que $f$ n'est pas définie en $a$, on pose $\tilde{f}(a) = \ell$.
\end{itemize}
\end{aretenir}

\courssub{Démarche et corrigé commenté}

\begin{exoresolu}[Corrigé complet de l'activité]
Reprendre successivement les six consignes : limites de $M$ en $+\infty$ et en $0^+$, levée d'indétermination pour $P$, continuité de $R$ en $50$, graphique et rapport.
\tcblower
\textbf{1. Limite de $M$ en $+\infty$.} Pour $x > 0$, on écrit :
\[ M(x) = \frac{50x + 200}{x} = 50 + \frac{200}{x}. \]
Or $\displaystyle\lim_{x \to +\infty} \frac{200}{x} = 0$ (limite usuelle). Par somme de limites :
\[ \lim_{x \to +\infty} M(x) = 50 + 0 = 50. \]
\textbf{Interprétation :} quand le puits devient très profond, les frais fixes de $200\,000$ FCFA se « diluent » sur un grand nombre de mètres, et le coût moyen se stabilise vers $50$ milliers de FCFA par mètre, c'est-à-dire le coût proportionnel pur. La droite d'équation $y = 50$ est \cle{asymptote horizontale} à la courbe de $M$ en $+\infty$.

\textbf{2. Limite de $M$ en $0^+$.} On a $\displaystyle\lim_{\substack{x \to 0 \\ x>0}} 50 = 50$ et $\displaystyle\lim_{\substack{x \to 0 \\ x>0}} \frac{200}{x} = +\infty$. Par somme :
\[ \lim_{\substack{x \to 0 \\ x > 0}} M(x) = +\infty. \]
\textbf{Interprétation :} pour un forage de quelques mètres seulement, les $200\,000$ FCFA de frais fixes sont répartis sur très peu de mètres : le coût moyen explose. Forer un puits très peu profond est donc économiquement absurde, ce qui conforte l'intuition du trésorier.

\textbf{3. Levée d'indétermination pour $P$.} Quand $x \to 0$ : $\sqrt{x+9} - 3 \to \sqrt{9} - 3 = 0$ et $x \to 0$. C'est une forme indéterminée $\dfrac{0}{0}$. Multiplions par la quantité conjuguée :
\begin{align*}
P(x) &= \frac{(\sqrt{x+9}-3)(\sqrt{x+9}+3)}{x(\sqrt{x+9}+3)} = \frac{(x+9) - 9}{x(\sqrt{x+9}+3)} \\
&= \frac{x}{x(\sqrt{x+9}+3)} = \frac{1}{\sqrt{x+9}+3} \quad (x > 0).
\end{align*}
La fonction $x \mapsto \sqrt{x+9}$ est continue en $0$, donc :
\[ \lim_{x \to 0} P(x) = \frac{1}{\sqrt{9}+3} = \frac{1}{6}. \]
$P$ n'est pas définie en $0$ mais admet une limite finie : elle est \cle{prolongeable par continuité} en $0$ par $\tilde{P}(0) = \dfrac{1}{6} \approx 0,167$ millier de FCFA par mètre, soit environ $167$ FCFA par mètre. La pénalité au démarrage du forage est donc faible et parfaitement maîtrisée.

\textbf{4. Continuité de $R$ en $50$.}
\[ \lim_{\substack{x \to 50 \\ x < 50}} R(x) = \frac{50}{5} = 10, \qquad \lim_{\substack{x \to 50 \\ x > 50}} R(x) = \frac{50}{5} + 8 = 18, \qquad R(50) = 18. \]
Les limites à gauche ($10$) et à droite ($18$) sont différentes : $R$ \textbf{n'est pas continue en $50$}. La fonction est continue à droite en $50$ (car $R(50) = 18$), mais pas à gauche. \textbf{Interprétation :} dès que l'on atteint $50$ m, la remise bondit instantanément de $10$ à $18$ milliers de FCFA : un « saut » de $8\,000$ FCFA. Un puits de $49$ m rapporte $9,8$ milliers de remise ; un puits de $50$ m en rapporte $18$ !

\textbf{5. Graphique.} Voir la figure ci-dessous : la courbe de $M$ décroît et se colle à son asymptote $y = 50$ ; la courbe de $R$ présente un saut visible en $x = 50$ (point plein en $(50\,;\,18)$, point creux en $(50\,;\,10)$).

\textbf{6. Rapport de conseil (modèle).} « Monsieur le Président, nos calculs montrent que le coût moyen par mètre diminue avec la profondeur et se stabilise vers $50\,000$ FCFA/m : forer profond ne pénalise pas le coût moyen. La pénalité de roche reste négligeable au démarrage ($\approx 167$ FCFA/m). Enfin, la remise subit un saut de $8\,000$ FCFA au palier de $50$ m : un puits de $49$ m fait perdre presque toute cette remise. Nous recommandons donc de forer \textbf{au moins $50$ mètres}, la contrainte technique de la nappe phréatique le permettant. »
\end{exoresolu}

\begin{popfigure}
\begin{tikzpicture}
\begin{axis}[width=0.92\linewidth, height=7cm, xlabel={Profondeur $x$ (m)}, ylabel={Milliers de FCFA}, xmin=0, xmax=105, ymin=0, ymax=120, legend pos=north east, grid=both, grid style={popline!40}, samples=200]
\addplot[popcurve, domain=2:100]{50 + 200/x};
\addlegendentry{$M(x) = 50 + \frac{200}{x}$}
\addplot[dashed, poporange, thick, domain=0:105]{50};
\addlegendentry{Asymptote $y = 50$}
\addplot[popgreen, thick, domain=1:49.9]{x/5};
\addplot[popgreen, thick, domain=50:100]{x/5 + 8};
\addlegendentry{$R(x)$}
\addplot[mark=*, mark size=1.5pt, popgreen, only marks] coordinates {(50,18)};
\addplot[mark=o, mark size=1.5pt, popgreen, only marks] coordinates {(50,10)};
\end{axis}
\end{tikzpicture}
\legende{Coût moyen $M(x)$ avec son asymptote horizontale $y = 50$, et remise $R(x)$ présentant un saut de discontinuité en $x = 50$.}
\end{popfigure}

\courssub{Bilan de l'activité}

Cette activité t'a permis de mettre en évidence que :

\begin{itemize}
\item les \textbf{limites en $+\infty$} décrivent des phénomènes de \cle{stabilisation} : ici, le coût moyen tend vers $50$ milliers de FCFA/m, ce qui se traduit graphiquement par une asymptote horizontale ;
\item les \textbf{limites en un point} (y compris à droite ou à gauche) décrivent des comportements locaux parfois spectaculaires : explosion du coût moyen près de $0$, saut de la remise en $50$ ;
\item une \textbf{forme indéterminée} n'est pas une impasse : la quantité conjuguée a transformé un quotient $\dfrac{0}{0}$ en une expression dont la limite se lit directement, et a permis un \cle{prolongement par continuité} ;
\item la \textbf{continuité} n'est pas une propriété abstraite : son absence se traduit ici par un « saut de tarif » de $8\,000$ FCFA, argument décisif dans une décision économique réelle ;
\item les mathématiques de Terminale fournissent des outils concrets d'\textbf{aide à la décision} pour des acteurs économiques camerounais comme les coopératives de la SODECOTON.
\end{itemize}

\begin{savaistu}[L'eau, un enjeu mathématique]
Dans le Grand Nord camerounais, la profondeur des nappes phréatiques varie fortement selon les localités, et le choix de la profondeur d'un forage est un véritable problème d'optimisation sous contrainte. Les ingénieurs hydrogéologues utilisent des modèles mathématiques bien plus complexes que celui de cette activité, mais l'idée de départ est la même : modéliser un coût par une fonction, puis étudier son comportement à l'aide de limites pour prendre la meilleure décision possible.
\end{savaistu}

\newpage

\coursec{Méthodes et exercices résolus}

\coursec{Limites et continuité}

\courssub{Intuition et définitions : limite en un point}

\begin{experience}[Approche intuitive de la limite]
Imaginons un taxi qui roule sur l'axe routier Yaoundé--Douala. Sa position $x$ (en km depuis Yaoundé) évolue avec le temps. On veut savoir où il se trouve \emph{au moment précis} $t_0 = 1$ h, mais le GPS bugge exactement à cet instant. On regarde alors les positions à $0{,}99$ h, $0{,}999$ h, $1{,}001$ h, $1{,}0001$ h\ldots Si toutes ces positions se rapprochent d'une même valeur $L$ (disons $150$ km), on dit que \textbf{la limite de la position à l'instant $1$ h est $150$ km}, peu importe la position réelle (ou l'absence de position) à $1$ h pile.
\end{experience}

\begin{definition}[Limite finie en un point]
Soit $f$ une fonction définie sur un intervalle pointé $I\setminus\{x_0\}$ (autour de $x_0$ mais pas nécessairement en $x_0$). On dit que $f$ admet la limite \textbf{finie} $\ell \in \mathbb{R}$ en $x_0$ si, quand $x$ se rapproche de $x_0$ (sans y être égal), $f(x)$ se rapproche arbitrairement de $\ell$.
\[
\lim_{x\to x_0} f(x) = \ell
\]
\end{definition}

\begin{definition}[Limites à droite et à gauche]
\begin{itemize}
\item \textbf{Limite à droite} en $x_0$ : $x$ tend vers $x_0$ par valeurs \textbf{supérieures} à $x_0$. Notation : $\displaystyle\lim_{x\to x_0^+} f(x)$ ou $\displaystyle\lim_{\substack{x\to x_0\\x>x_0}} f(x)$.
\item \textbf{Limite à gauche} en $x_0$ : $x$ tend vers $x_0$ par valeurs \textbf{inférieures} à $x_0$. Notation : $\displaystyle\lim_{x\to x_0^-} f(x)$ ou $\displaystyle\lim_{\substack{x\to x_0\\x<x_0}} f(x)$.
\end{itemize}
\end{definition}

\begin{propriete}[Existence d'une limite en un point]
Une fonction $f$ admet une limite (finie ou infinie) en $x_0$ \textbf{si et seulement si} les limites à droite et à gauche existent et sont \textbf{égales} :
\[
\lim_{x\to x_0} f(x) = \ell \quad\Longleftrightarrow\quad \lim_{x\to x_0^-} f(x) = \lim_{x\to x_0^+} f(x) = \ell.
\]
\end{propriete}

\begin{definition}[Limite infinie en un point]
On dit que $f$ a pour limite $+\infty$ (resp. $-\infty$) en $x_0$ si, quand $x\to x_0$, $f(x)$ devient arbitrairement grand (resp. petit) en valeur absolue et reste positif (resp. négatif).
Notations : $\displaystyle\lim_{x\to x_0} f(x) = +\infty$ ou $-\infty$.
\end{definition}

\courssub{Limite en l'infini}

\begin{definition}[Limite en $+\infty$ et $-\infty$]
\begin{itemize}
\item $\displaystyle\lim_{x\to +\infty} f(x) = \ell$ signifie que $f(x)$ se rapproche de $\ell$ quand $x$ devient arbitrairement grand.
\item $\displaystyle\lim_{x\to -\infty} f(x) = \ell$ signifie que $f(x)$ se rapproche de $\ell$ quand $x$ devient arbitrairement petit (grand en valeur absolue, négatif).
\item Les limites infinies en $\pm\infty$ se définissent de même ($f(x)\to +\infty$ ou $-\infty$).
\end{itemize}
\end{definition}

\begin{propriete}[Limites usuelles admises au Bac]
\begin{itemize}
\item \textbf{Puissances :} $\forall n\in\mathbb{N}^*$, $\displaystyle\lim_{x\to +\infty} x^n = +\infty$ ; $\displaystyle\lim_{x\to -\infty} x^n = \begin{cases} +\infty & \text{si $n$ pair} \\ -\infty & \text{si $n$ impair} \end{cases}$.
\item \textbf{Inverse :} $\displaystyle\lim_{x\to \pm\infty} \frac{1}{x} = 0$ ; $\displaystyle\lim_{x\to 0^+} \frac{1}{x} = +\infty$ ; $\displaystyle\lim_{x\to 0^-} \frac{1}{x} = -\infty$.
\item \textbf{Trigonométrique (radian) :} $\displaystyle\lim_{x\to 0} \frac{\sin x}{x} = 1$ \quad (et par conséquent $\displaystyle\lim_{x\to 0} \frac{\tan x}{x} = 1$).
\item \textbf{Exponentielle/Logarithme (rappel Tle D/S, culture pour A) :} $\displaystyle\lim_{x\to -\infty} e^x = 0$, $\displaystyle\lim_{x\to +\infty} e^x = +\infty$, $\displaystyle\lim_{x\to 0^+} \ln x = -\infty$, $\displaystyle\lim_{x\to +\infty} \ln x = +\infty$.
\end{itemize}
\end{propriete}

\begin{aretenir}[Comportement dominant en $\pm\infty$]
\begin{itemize}
\item Un \textbf{polynôme} $P(x) = a_n x^n + \dots + a_0$ ($a_n\neq 0$) se comporte en $\pm\infty$ comme son \cle{monôme de plus haut degré} $a_n x^n$.
\item Une \textbf{fraction rationnelle} $\frac{P(x)}{Q(x)}$ se comporte en $\pm\infty$ comme le \cle{quotient des monômes de plus haut degré} $\frac{a_n x^n}{b_m x^m}$.
\end{itemize}
\end{aretenir}

\courssub{Opérations sur les limites}

\begin{propriete}[Théorèmes d'opérations sur les limites]
Soient $f$ et $g$ deux fonctions et $x_0 \in \mathbb{R} \cup \{+\infty, -\infty\}$. On suppose que $\lim_{x\to x_0} f(x)$ et $\lim_{x\to x_0} g(x)$ \textbf{existent} (finies ou infinies).
\begin{enumerate}
\item \textbf{Somme :} $\lim (f+g) = \lim f + \lim g$ \quad (interdit : $+\infty + (-\infty)$).
\item \textbf{Produit :} $\lim (f\times g) = \lim f \times \lim g$ \quad (interdit : $0 \times \infty$).
\item \textbf{Quotient :} $\lim \frac{f}{g} = \frac{\lim f}{\lim g}$ \quad \textbf{si $\lim g \neq 0$} (interdit : $\frac{0}{0}$, $\frac{\infty}{\infty}$, $\frac{\ell}{0}$).
\item \textbf{Composée :} Si $\lim_{x\to x_0} g(x) = \ell$ et $f$ continue en $\ell$, alors $\lim_{x\to x_0} f(g(x)) = f(\ell)$.
\end{enumerate}
\textbf{Attention :} Les cas interdits sont les \cle{formes indéterminées} (F.I.) : $\infty-\infty$, $0\times\infty$, $\frac{0}{0}$, $\frac{\infty}{\infty}$. Elles ne donnent \textbf{pas} le résultat ; il faut transformer l'expression.
\end{propriete}

\begin{methode}[Calculer une limite par opérations]
Pour déterminer $\displaystyle\lim_{x\to a} f(x)$ où $f$ est obtenue par opérations sur des fonctions usuelles :
\begin{enumerate}
\item \textbf{Identifier} le point $a$ (réel, $+\infty$ ou $-\infty$) et \textbf{décomposer} $f$ en morceaux simples.
\item \textbf{Calculer la limite de chaque morceau} à l'aide des limites usuelles et du comportement dominant.
\item \textbf{Appliquer les théorèmes d'opérations} en \textbf{citant le théorème utilisé} (« par somme », « par produit », « par quotient car la limite du dénominateur est non nulle », « par composition car $f$ est continue en $\ell$ »).
\item \textbf{Vérifier} qu'on n'obtient pas une forme indéterminée. Si c'est le cas, \emph{annoncer explicitement la F.I.} et passer à la méthode de levée.
\item \textbf{Conclure} par une phrase complète : « Donc $\displaystyle\lim_{x\to a} f(x) = \ldots$ ».
\end{enumerate}
\end{methode}

\begin{exoresolu}[Limites par opérations — type Bac]
Calculer les limites suivantes :
\begin{enumerate}
\item $\displaystyle\lim_{x\to +\infty}\left(3x^2 - 5x + 7\right)$ ;
\item $\displaystyle\lim_{x\to -\infty}\frac{2x^2 + 3x - 1}{x^2 + 4}$ ;
\item $\displaystyle\lim_{x\to 0}\frac{\sin(3x)}{x}$.
\end{enumerate}
\tcblower
\textbf{Solution détaillée.}

\textbf{1.} $f(x)=3x^2-5x+7$ est un polynôme. En $+\infty$, il a même limite que son monôme dominant $3x^2$.
\[\lim_{x\to+\infty} x^2 = +\infty \quad\text{et}\quad 3>0 \quad\Longrightarrow\quad \lim_{x\to+\infty} 3x^2 = +\infty.\]
\textbf{Par somme} (limites finies ou $+\infty$), $\displaystyle\lim_{x\to+\infty}\left(3x^2-5x+7\right)=+\infty$.

\emph{Rédaction rigoureuse par factorisation (modèle Bac)} : pour $x\neq 0$,
\[3x^2-5x+7 = x^2\left(3-\frac{5}{x}+\frac{7}{x^2}\right).\]
$\displaystyle\lim_{x\to+\infty}\frac{5}{x}=0$, $\displaystyle\lim_{x\to+\infty}\frac{7}{x^2}=0$ \textbf{(limites usuelles)}.
\textbf{Par somme} : $\displaystyle\lim_{x\to+\infty}\left(3-\frac{5}{x}+\frac{7}{x^2}\right)=3$.
\textbf{Par produit} : $(+\infty)\times 3 = +\infty$.

\textbf{2.} Fraction rationnelle en $-\infty$. On factorise par les monômes dominants ($x^2$). Pour $x\neq 0$ :
\[\frac{2x^2+3x-1}{x^2+4}=\frac{x^2\left(2+\dfrac{3}{x}-\dfrac{1}{x^2}\right)}{x^2\left(1+\dfrac{4}{x^2}\right)}=\frac{2+\dfrac{3}{x}-\dfrac{1}{x^2}}{1+\dfrac{4}{x^2}}.\]
Limites usuelles : $\displaystyle\lim_{x\to-\infty}\frac{3}{x}=0$, $\displaystyle\lim_{x\to-\infty}\frac{1}{x^2}=0$, $\displaystyle\lim_{x\to-\infty}\frac{4}{x^2}=0$.
\textbf{Par somme} : numérateur $\to 2$, dénominateur $\to 1 \neq 0$.
\textbf{Par quotient} : $\displaystyle\lim_{x\to-\infty}\frac{2x^2+3x-1}{x^2+4}=\frac{2}{1}=2$.

\textbf{3.} On utilise $\displaystyle\lim_{u\to 0}\frac{\sin u}{u}=1$. Posons $u=3x$ ; quand $x\to 0$, $u\to 0$. Pour $x\neq 0$ :
\[\frac{\sin(3x)}{x}=\frac{\sin(3x)}{3x}\times 3 = 3\,\frac{\sin u}{u}.\]
\textbf{Par composition} ($u=3x$ tend vers $0$) \textbf{et par produit} :
\[\lim_{x\to 0}\frac{\sin(3x)}{x}=3\times 1 = 3.\]

\textbf{Conclusion} : $\displaystyle\lim_{x\to+\infty}(3x^2-5x+7)=+\infty$ ; $\displaystyle\lim_{x\to-\infty}\frac{2x^2+3x-1}{x^2+4}=2$ ; $\displaystyle\lim_{x\to 0}\frac{\sin(3x)}{x}=3$.
\end{exoresolu}

\courssub{Lever une forme indéterminée}

\begin{methode}[Lever une forme indéterminée simple]
Face à une F.I. ($\frac{0}{0}$, $\frac{\infty}{\infty}$, $\infty-\infty$, $0\times\infty$) :
\begin{enumerate}
\item \textbf{Reconnaître et annoncer} la F.I. : « On obtient une forme indéterminée $\frac{0}{0}$ ».
\item \textbf{Choisir la technique} adaptée :
\begin{itemize}
\item F.I. $\frac{0}{0}$ avec polynômes en $x_0$ : \textbf{factoriser} numérateur et dénominateur par $(x-x_0)$, simplifier.
\item F.I. avec racines carrées : multiplier par la \textbf{quantité conjuguée} $\sqrt{a}-\sqrt{b}=\dfrac{a-b}{\sqrt{a}+\sqrt{b}}$.
\item F.I. $\frac{\infty}{\infty}$ en $\pm\infty$ : factoriser par le \textbf{monôme dominant}.
\item F.I. $\infty-\infty$ : mettre le terme dominant en facteur ou utiliser le conjugué.
\end{itemize}
\item \textbf{Simplifier} (l'expression n'est plus indéterminée).
\item \textbf{Calculer} la limite de l'expression simplifiée et conclure.
\end{enumerate}
\end{methode}

\begin{exoresolu}[Levée d'indétermination par factorisation et quantité conjuguée]
\begin{enumerate}
\item Calculer $\displaystyle\lim_{x\to 2}\frac{x^2-5x+6}{x^2-4}$.
\item Calculer $\displaystyle\lim_{x\to +\infty}\left(\sqrt{x^2+3x}-x\right)$.
\end{enumerate}
\tcblower
\textbf{Solution détaillée.}

\textbf{1.} Pour $x\to 2$ : numérateur $\to 4-10+6=0$, dénominateur $\to 4-4=0$. \textbf{F.I. $\frac{0}{0}$}.
$2$ annule le numérateur, donc $(x-2)$ est facteur. L'autre racine de $x^2-5x+6$ est $3$ (produit $=6$).
\[x^2-5x+6=(x-2)(x-3),\quad x^2-4=(x-2)(x+2).\]
Pour $x\neq 2$, on simplifie :
\[\frac{x^2-5x+6}{x^2-4}=\frac{(x-2)(x-3)}{(x-2)(x+2)}=\frac{x-3}{x+2}.\]
Expression simplifiée : numérateur $\to -1$, dénominateur $\to 4\neq 0$.
\textbf{Par quotient} : $\displaystyle\lim_{x\to 2}\frac{x^2-5x+6}{x^2-4}=\frac{-1}{4}=-\frac{1}{4}$.

\textbf{2.} Pour $x\to+\infty$ : $\sqrt{x^2+3x}\to+\infty$, $x\to+\infty$. \textbf{F.I. $\infty-\infty$}.
On multiplie par la quantité conjuguée ($x>0$) :
\begin{align*}
\sqrt{x^2+3x}-x &= \frac{\left(\sqrt{x^2+3x}-x\right)\left(\sqrt{x^2+3x}+x\right)}{\sqrt{x^2+3x}+x} \\
&= \frac{(x^2+3x)-x^2}{\sqrt{x^2+3x}+x} = \frac{3x}{\sqrt{x^2+3x}+x}.
\end{align*}
Factorisation par le dominant $x$ (car $x>0 \Rightarrow \sqrt{x^2}=x$) :
\[\frac{3x}{\sqrt{x^2+3x}+x}=\frac{3x}{x\sqrt{1+\frac{3}{x}}+x}=\frac{3x}{x\left(\sqrt{1+\frac{3}{x}}+1\right)}=\frac{3}{\sqrt{1+\frac{3}{x}}+1}.\]
$\displaystyle\lim_{x\to+\infty}\frac{3}{x}=0 \Rightarrow$ \textbf{par composition} (racine continue en $1$) $\sqrt{1+\frac{3}{x}}\to 1$.
Dénominateur $\to 2 \neq 0$. \textbf{Par quotient} : $\displaystyle\lim_{x\to+\infty}\left(\sqrt{x^2+3x}-x\right)=\frac{3}{2}$.

\textbf{Conclusion} : $\displaystyle\lim_{x\to 2}\frac{x^2-5x+6}{x^2-4}=-\frac{1}{4}$ ; $\displaystyle\lim_{x\to+\infty}\left(\sqrt{x^2+3x}-x\right)=\dfrac{3}{2}$.
\end{exoresolu}

\courssub{Limites latérales et asymptotes verticales}

\begin{methode}[Limites à droite, à gauche et existence]
Pour étudier la limite de $f$ en $x_0$ (dénominateur nul ou fonction par morceaux) :
\begin{enumerate}
\item \textbf{Repérer} si $x_0$ annule un dénominateur ou sépare deux expressions.
\item \textbf{Calculer séparément} $\displaystyle\lim_{x\to x_0^-} f(x)$ et $\displaystyle\lim_{x\to x_0^+} f(x)$.
\item Pour un quotient $\frac{N(x)}{D(x)}$ avec $\lim D(x)=0$ et $\lim N(x)=\ell\neq 0$ : \textbf{étudier le signe de $D(x)$} de chaque côté ($0^+$ ou $0^-$). Règle : $\frac{\ell}{0^+}=\operatorname{signe}(\ell)\times\infty$, $\frac{\ell}{0^-}=-\operatorname{signe}(\ell)\times\infty$.
\item \textbf{Conclure} : limite existante $\iff$ limites latérales égales.
\end{enumerate}
\end{methode}

\begin{exoresolu}[Limites latérales — fonction définie par morceaux]
Soit $f$ définie sur $\mathbb{R}$ par :
\[f(x)=\begin{cases} \dfrac{x^2-1}{x-1} & \text{si } x<1 \\[6pt] 3x-1 & \text{si } x\geq 1 \end{cases}\]
\begin{enumerate}
\item Calculer $\displaystyle\lim_{x\to 1^-} f(x)$.
\item Calculer $\displaystyle\lim_{x\to 1^+} f(x)$.
\item $f$ admet-elle une limite en $1$ ?
\end{enumerate}
\tcblower
\textbf{Solution détaillée.}

\textbf{1. Limite à gauche ($x\to 1^-$).} Pour $x<1$, $f(x)=\dfrac{x^2-1}{x-1}$. F.I. $\frac{0}{0}$.
Factorisation : $x^2-1=(x-1)(x+1)$. Pour $x<1$ ($x\neq 1$) :
\[f(x)=\frac{(x-1)(x+1)}{x-1}=x+1.\]
$\displaystyle\lim_{x\to 1^-} f(x)=\lim_{x\to 1}(x+1)=2.$

\textbf{2. Limite à droite ($x\to 1^+$).} Pour $x\geq 1$, $f(x)=3x-1$ (polynôme).
$\displaystyle\lim_{x\to 1^+} f(x)=3\times 1-1=2.$

\textbf{3.} $\displaystyle\lim_{x\to 1^-} f(x)=2=\lim_{x\to 1^+} f(x)$. Les limites latérales sont égales, donc
\[\lim_{x\to 1} f(x)=2.\]
\end{exoresolu}

\begin{exoresolu}[Limite infinie en un point — signe du dénominateur]
Soit $g(x)=\dfrac{2x+1}{x-3}$. Calculer les limites en $3$ à gauche et à droite, interpréter graphiquement.
\tcblower
\textbf{Solution détaillée.}

$g$ n'est pas définie en $3$ (dénominateur nul). Numérateur $\to 2\times 3+1=7>0$.

\textbf{À droite ($x\to 3^+$) :} $x>3 \Rightarrow x-3>0 \Rightarrow$ dénominateur $\to 0^+$.
$\displaystyle\lim_{x\to 3^+}\frac{2x+1}{x-3}=\frac{7}{0^+}=+\infty.$

\textbf{À gauche ($x\to 3^-$) :} $x<3 \Rightarrow x-3<0 \Rightarrow$ dénominateur $\to 0^-$.
$\displaystyle\lim_{x\to 3^-}\frac{2x+1}{x-3}=\frac{7}{0^-}=-\infty.$

\textbf{Conclusion.} Limites latérales différentes ($-\infty\neq+\infty$), pas de limite en $3$.
\textbf{Graphiquement :} la droite $x=3$ est une \cle{asymptote verticale}. La courbe « descend » vers $-\infty$ à gauche, « monte » vers $+\infty$ à droite.
\end{exoresolu}

\courssub{Continuité en un point}

\begin{definition}[Continuité en un point]
Une fonction $f$ est \textbf{continue en $x_0$} si et seulement si les trois conditions suivantes sont réunies :
\begin{enumerate}
\item $f$ est \textbf{définie} en $x_0$ ($x_0 \in D_f$).
\item $f$ admet une \textbf{limite finie} en $x_0$.
\item $\displaystyle\lim_{x\to x_0} f(x) = f(x_0)$.
\end{enumerate}
Équivalent (utile pour fonctions par morceaux) :
\[\lim_{x\to x_0^-} f(x) = \lim_{x\to x_0^+} f(x) = f(x_0).\]
\end{definition}

\begin{propriete}[Fonctions continues sur leur domaine]
Les fonctions suivantes sont continues sur tout leur ensemble de définition :
\begin{itemize}
\item Polynômes (sur $\mathbb{R}$).
\item Fonctions rationnelles (sur $\mathbb{R}\setminus\{\text{zéros du dénominateur}\}$).
\item $\sin$, $\cos$ (sur $\mathbb{R}$).
\item $\sqrt{\cdot}$ (sur $[0;+\infty[$).
\item Valeur absolue $x\mapsto |x|$ (sur $\mathbb{R}$).
\end{itemize}
Sommes, produits, quotients (dénominateur non nul) et composées de fonctions continues sont continues.
\end{propriete}

\begin{methode}[Étudier la continuité en $x_0$]
\begin{enumerate}
\item \textbf{Vérifier que $f(x_0)$ existe} (calculer $f(x_0)$). Si non $\Rightarrow$ discontinue (stop).
\item \textbf{Calculer $\lim_{x\to x_0^-} f(x)$ et $\lim_{x\to x_0^+} f(x)$} (ou limite unique si expression unique).
\item \textbf{Comparer} : continue $\iff$ limite gauche = limite droite = $f(x_0)$.
\item Si discontinue : préciser le type (saut si limites finies différentes, asymptote verticale si limite infinie, trou si limite finie $\neq f(x_0)$).
\end{enumerate}
\end{methode}

\begin{exoresolu}[Continuité en un point — contexte camerounais]
Une entreprise de livraison à Yaoundé facture selon le poids $x$ (kg) :
\[f(x)=\begin{cases} 500x & \text{si } 0\leq x\leq 2 \\ 1000+300(x-2) & \text{si } x>2 \end{cases}\quad (f(x) \text{ en F CFA}).\]
\begin{enumerate}
\item $f$ est-elle continue en $x_0=2$ ?
\item Interpréter pour le client.
\end{enumerate}
\tcblower
\textbf{Solution détaillée.}

\textbf{1.} \emph{Étape 1 :} $2\in[0;2] \Rightarrow f(2)=500\times 2=1000$. Définie.

\emph{Étape 2 :} Limites latérales.
$x\to 2^-$ ($x<2$) : $f(x)=500x \to 1000$.
$x\to 2^+$ ($x>2$) : $f(x)=1000+300(x-2) \to 1000$.

\emph{Étape 3 :} $\lim_{x\to 2^-}f(x)=\lim_{x\to 2^+}f(x)=f(2)=1000$.
\textbf{$f$ est continue en $2$.}

\textbf{2.} Pas de « saut » de prix à $2$ kg. $1{,}99$ kg $\approx 995$ F, $2{,}01$ kg $\approx 1003$ F. Tarification équitable, pas de rupture brutale.
\end{exoresolu}

\begin{exoresolu}[Discontinuité et paramètre]
Soit $h$ définie sur $\mathbb{R}$ par :
\[h(x)=\begin{cases} x^2+1 & \text{si } x<1 \\ ax+2 & \text{si } x\geq 1 \end{cases}\quad (a\in\mathbb{R}).\]
\begin{enumerate}
\item Montrer que $h$ est continue en $1$ si et seulement si $a=0$.
\item Pour $a=3$, nature de la discontinuité en $1$ ?
\end{enumerate}
\tcblower
\textbf{Solution détaillée.}

\textbf{1.} \emph{Étape 1 :} $h(1)=a\times 1+2=a+2$.

\emph{Étape 2 :} Limites latérales.
$x\to 1^-$ : $h(x)=x^2+1 \to 2$.
$x\to 1^+$ : $h(x)=ax+2 \to a+2$.

\emph{Étape 3 :} Continuité $\iff 2 = a+2 = a+2 \iff a=0$.
\textbf{Conclusion :} $h$ continue en $1 \iff a=0$.

\textbf{2.} Pour $a=3$ : $\lim_{x\to 1^-}h(x)=2$, $\lim_{x\to 1^+}h(x)=h(1)=5$.
Limites finies mais différentes ($2\neq 5$) $\Rightarrow$ \cle{discontinuité par saut} (hauteur $3$). Graphiquement : coupure, il faut « lever le crayon ».
\end{exoresolu}

\courssub{Continuité sur un intervalle et interprétation graphique}

\begin{methode}[Continuité sur un intervalle et lecture graphique]
\begin{enumerate}
\item \textbf{Continuité sur $I$} : $f$ continue en \emph{tout} point de $I$.
\begin{itemize}
\item Citer les théorèmes : polynômes sur $\mathbb{R}$, rationnelles sur intervalles de $D_f$, $\sin/\cos$ sur $\mathbb{R}$, $\sqrt{\cdot}$ sur $[0;+\infty[$.
\item Opérations (somme, produit, quotient, composée) conservent la continuité.
\item Traiter \textbf{séparément} les points de jonction des fonctions par morceaux.
\end{itemize}
\item \textbf{Lecture graphique} :
\begin{itemize}
\item Fonction continue $\Leftrightarrow$ courbe tracée \textbf{sans lever le crayon}.
\item Coupure, saut, point isolé $\Rightarrow$ discontinuité.
\item Limite finie $\ell$ en $+\infty$ $\Leftrightarrow$ \cle{asymptote horizontale} $y=\ell$.
\item Limite infinie en $x_0$ $\Leftrightarrow$ \cle{asymptote verticale} $x=x_0$.
\end{itemize}
\end{enumerate}
\end{methode}

\begin{exoresolu}[Continuité sur intervalle et graphique]
$f(x)=\dfrac{x^2+1}{x-2}$.
\begin{enumerate}
\item Déterminer $D_f$.
\item Justifier la continuité sur $]-\infty;2[$ et $]2;+\infty[$.
\item Limites en $2^-$, $2^+$, $+\infty$ et interprétation graphique.
\end{enumerate}
\tcblower
\textbf{Solution détaillée.}

\textbf{1.} $x-2=0 \iff x=2$. $D_f=\mathbb{R}\setminus\{2\}=]-\infty;2[\cup]2;+\infty[$.

\textbf{2.} $f$ fraction rationnelle $\Rightarrow$ continue sur chaque intervalle de $D_f$.
Donc continue sur $]-\infty;2[$ et sur $]2;+\infty[$.
\textit{Remarque :} On ne dit pas « continue sur $D_f$ » (réunion disjointe), mais « continue \emph{sur chacun des intervalles}... ».

\textbf{3.} \emph{En $2$ :} Numérateur $\to 5>0$, dénominateur $\to 0$.
$x\to 2^+$ : $x-2>0 \Rightarrow 0^+ \Rightarrow \lim_{x\to 2^+}f(x)=+\infty$.
$x\to 2^-$ : $x-2<0 \Rightarrow 0^- \Rightarrow \lim_{x\to 2^-}f(x)=-\infty$.
\textbf{Asymptote verticale :} $x=2$.

\emph{En $+\infty$ :} Comportement dominant $\frac{x^2}{x}=x$.
$f(x)=\frac{x^2(1+1/x^2)}{x(1-2/x)}=x\frac{1+1/x^2}{1-2/x}$.
Par produit : $x\to+\infty$, fraction $\to 1 \Rightarrow \lim_{x\to+\infty}f(x)=+\infty$.
\textbf{Pas d'asymptote horizontale} en $+\infty$ (la courbe monte indéfiniment).
\end{exoresolu}

\begin{savaistu}[Le théorème des valeurs intermédiaires (TVI) — Outil puissant]
Si $f$ est continue sur $[a;b]$, alors pour toute valeur $k$ entre $f(a)$ et $f(b)$, il existe au moins un $c\in[a;b]$ tel que $f(c)=k$.
\textbf{Application concrète :} Un camion part de Yaoundé (altitude $\approx 700$ m) vers Douala (niveau de la mer, $0$ m). Son altitude varie continûment. À un moment donné, il a \emph{nécessairement} franchi l'altitude $350$ m. Le TVI garantit l'existence de solutions d'équations $f(x)=k$ sans les calculer explicitement. C'est la base de la méthode de dichotomie (recherche de zéros).
\end{savaistu}

\begin{attention}[Les 5 erreurs qui coûtent des points au Bac]
\begin{enumerate}
\item \textbf{Écrire « $\lim = \frac{0}{0}$ » comme résultat final.} Une F.I. n'est jamais une réponse. L'annoncer (« on obtient une F.I. $\frac{0}{0}$ »), la lever, puis donner la limite.
\item \textbf{Oublier $0^+$ et $0^-$.} Dénominateur $\to 0$ $\Rightarrow$ signe crucial : $\frac{5}{0^+}=+\infty$, $\frac{5}{0^-}=-\infty$. Toujours étudier le signe de chaque côté.
\item \textbf{Conclure la continuité sans les 3 conditions.} Il faut : $f(x_0)$ existe, limites latérales existent, \textbf{égalité des trois}. Pour fonction par morceaux, l'étude des limites latérales est obligatoire.
\item \textbf{Factoriser mal en $\pm\infty$.} Facteur = monôme de \emph{plus haut degré}. Piège $\sqrt{x^2}=|x|$ : vaut $x$ si $x>0$, $-x$ si $x<0$ (attention en $-\infty$ !).
\item \textbf{Appliquer les opérations hors conditions.} Quotient des limites seulement si limite dénominateur $\neq 0$. Somme seulement si pas $\infty-\infty$. \textbf{Citer le théorème} (« par quotient car... ») sécurise la copie.
\end{enumerate}
\end{attention}

\begin{exercice}[Entraînement Bac — Limites et continuité]
\begin{enumerate}
\item Calculer : $\lim_{x\to -\infty}\frac{3x^3-2x+1}{x^2+5}$ ; $\lim_{x\to 0}\frac{\sin(5x)}{2x}$ ; $\lim_{x\to 1}\frac{x^3-1}{x^2-1}$.
\item Soit $f(x)=\frac{x+1}{x-1}$. Étudier les limites en $1^-$, $1^+$, $+\infty$. Tracer les asymptotes.
\item Soit $g(x)=\begin{cases} x^2 & x\leq 0 \\ \sqrt{x} & x>0 \end{cases}$. Étudier la continuité en $0$.
\item Déterminer $a$ pour que $h(x)=\begin{cases} 2x+a & x<2 \\ x^2 & x\geq 2 \end{cases}$ soit continue sur $\mathbb{R}$.
\end{enumerate}
\end{exercice}

\begin{corrige}[Exercice n°1 -- Entraînement Bac]
\textbf{1.}
\begin{itemize}
\item $\frac{3x^3}{x^2}=3x \to -\infty$ en $-\infty$ (degré num > degré den, signe $-$). \textbf{Réponse : $-\infty$}.
\item $\frac{\sin(5x)}{2x}=\frac{5}{2}\frac{\sin(5x)}{5x} \to \frac{5}{2}\times 1 = 2,5$. \textbf{Réponse : $2,5$}.
\item F.I. $\frac{0}{0}$ en $1$. $x^3-1=(x-1)(x^2+x+1)$, $x^2-1=(x-1)(x+1)$. Simplification $\frac{x^2+x+1}{x+1} \to \frac{3}{2}$. \textbf{Réponse : $1,5$}.
\end{itemize}
\textbf{2.} $f(x)=\frac{x+1}{x-1}$.
$1^-$ : num $\to 2$, den $\to 0^- \Rightarrow -\infty$.
$1^+$ : num $\to 2$, den $\to 0^+ \Rightarrow +\infty$.
$+\infty$ : $\frac{x}{x}=1 \Rightarrow$ limite $1$ (asymptote horizontale $y=1$).
Asymptotes : verticale $x=1$, horizontale $y=1$.
\textbf{3.} $g(0)=0$. $\lim_{x\to 0^-}x^2=0$. $\lim_{x\to 0^+}\sqrt{x}=0$. Triple égalité $\Rightarrow$ \textbf{continue en $0$}.
\textbf{4.} $h(2)=4$. $\lim_{x\to 2^-}(2x+a)=4+a$. $\lim_{x\to 2^+}x^2=4$.
Continuité $\iff 4+a=4 \iff a=0$.
\end{corrige}

\newpage

\coursec{Exercices}

\coursec{Limites et continuité}

\begin{definition}[Limite d'une fonction en un point]
Soit $f$ une fonction définie sur un intervalle $I$ (sauf éventuellement en $x_0$) et $x_0 \in \mathbb{R}$ une valeur d'adhérence de $I$. On dit que $f$ admet la limite $\ell \in \mathbb{R}$ en $x_0$, et on note
\[ \lim_{x\to x_0}f(x)=\ell, \]
si pour tout $\varepsilon>0$, il existe $\eta>0$ tel que pour tout $x \in I$, $0<|x-x_0|<\eta \implies |f(x)-\ell|<\varepsilon$.
\end{definition}

\begin{definition}[Limites à droite et à gauche]
\begin{itemize}
\item \textbf{Limite à droite :} $\displaystyle\lim_{x\to x_0^+}f(x)=\ell$ si $\forall \varepsilon>0,\ \exists \eta>0,\ \forall x \in I,\ x_0<x<x_0+\eta \implies |f(x)-\ell|<\varepsilon$.
\item \textbf{Limite à gauche :} $\displaystyle\lim_{x\to x_0^-}f(x)=\ell$ si $\forall \varepsilon>0,\ \exists \eta>0,\ \forall x \in I,\ x_0-\eta<x<x_0 \implies |f(x)-\ell|<\varepsilon$.
\end{itemize}
\end{definition}

\begin{propriete}[Théorème de la limite bilatère]
La fonction $f$ admet une limite $\ell$ en $x_0$ \textbf{si et seulement si} elle admet une limite à droite et une limite à gauche en $x_0$, et ces deux limites sont égales à $\ell$ :
\[ \lim_{x\to x_0}f(x)=\ell \iff \lim_{x\to x_0^-}f(x)=\lim_{x\to x_0^+}f(x)=\ell. \]
\end{propriete}

\begin{definition}[Limite en l'infini]
\begin{itemize}
\item $\displaystyle\lim_{x\to +\infty}f(x)=\ell$ si $\forall \varepsilon>0,\ \exists A>0,\ \forall x>A,\ |f(x)-\ell|<\varepsilon$.
\item $\displaystyle\lim_{x\to +\infty}f(x)=+\infty$ si $\forall M>0,\ \exists A>0,\ \forall x>A,\ f(x)>M$.
\item Les définitions pour $x\to -\infty$ sont symétriques.
\end{itemize}
\end{definition}

\begin{aretenir}[Limites usuelles (admises au Bac)]
\begin{itemize}
\item $\displaystyle\lim_{x\to 0}\frac{\sin x}{x}=1$, \quad $\displaystyle\lim_{x\to 0}\frac{\tan x}{x}=1$, \quad $\displaystyle\lim_{x\to 0}\frac{e^x-1}{x}=1$, \quad $\displaystyle\lim_{x\to 0}\frac{\ln(1+x)}{x}=1$.
\item $\displaystyle\lim_{x\to 0}\frac{\sqrt{x+1}-1}{x}=\frac{1}{2}$.
\item Polynômes : $\displaystyle\lim_{x\to \pm\infty}P(x)=\pm\infty \cdot \text{signe}(a_n)$ (terme de plus haut degré).
\item Fonctions rationnelles : $\displaystyle\lim_{x\to \pm\infty}\frac{P(x)}{Q(x)}=\frac{a_n}{b_m}$ si $\deg P = \deg m$ ; $0$ si $\deg P < \deg Q$ ; $\pm\infty$ si $\deg P > \deg Q$.
\end{itemize}
\end{aretenir}

\begin{propriete}[Opérations sur les limites (limites finies)]
Si $\lim_{x\to x_0}f(x)=\ell$ et $\lim_{x\to x_0}g(x)=\ell'$ (avec $x_0 \in \mathbb{R}\cup\{\pm\infty\}$, $\ell,\ell'\in\mathbb{R}$), alors :
\begin{align*}
\lim (f+g) &= \ell+\ell', & \lim (f-g) &= \ell-\ell', \\
\lim (f\times g) &= \ell\ell', & \lim \frac{f}{g} &= \frac{\ell}{\ell'} \quad (\text{si } \ell'\neq 0).
\end{align*}
Si de plus $h$ est continue en $\ell$, $\lim h(f(x)) = h(\ell)$.
\end{propriete}

\begin{attention}[Formes indéterminées]
Les opérations ci-dessus \textbf{ne s'appliquent pas} dans les cas suivants (formes indéterminées) :
\[ \frac{0}{0},\quad \frac{\infty}{\infty},\quad 0\times\infty,\quad \infty-\infty,\quad 1^\infty,\quad 0^0,\quad \infty^0. \]
Il faut alors transformer l'expression (factorisation, quantité conjuguée, limites usuelles, encadrement) pour lever l'indétermination.
\end{attention}

\begin{methode}[Lever une forme indéterminée $0/0$ ou $\infty/\infty$ pour une fraction rationnelle]
\begin{enumerate}
\item Factoriser le numérateur et le dénominateur (identités remarquables, racines évidentes).
\item Simplifier les facteurs communs qui s'annulent en $x_0$ (ou le terme de plus haut degré pour $x\to\pm\infty$).
\item Calculer la limite de l'expression simplifiée par les opérations sur les limites.
\end{enumerate}
\end{methode}

\begin{methode}[Lever une forme indéterminée $\infty-\infty$ ou $0/0$ avec une racine carrée]
Utiliser la \textbf{quantité conjuguée} :
\begin{itemize}
\item Pour $\sqrt{A}-\sqrt{B}$ : multiplier par $\dfrac{\sqrt{A}+\sqrt{B}}{\sqrt{A}+\sqrt{B}}$.
\item Pour $\sqrt{A}-B$ : multiplier par $\dfrac{\sqrt{A}+B}{\sqrt{A}+B}$.
\end{itemize}
On utilise l'identité $(\sqrt{A}-\sqrt{B})(\sqrt{A}+\sqrt{B})=A-B$.
\end{methode}

\begin{definition}[Continuité en un point]
Une fonction $f$ est \cle{continue} en $x_0 \in D_f$ si et seulement si :
\[ \lim_{x\to x_0}f(x)=f(x_0). \]
Cela implique trois conditions : $f$ définie en $x_0$, limite existante (finie), égalité des deux.
\end{definition}

\begin{propriete}[Continuité sur un intervalle]
$f$ est continue sur un intervalle $I$ si elle est continue en tout point de $I$.
\begin{itemize}
\item Toute fonction polynomiale est continue sur $\mathbb{R}$.
\item Toute fonction rationnelle est continue sur son ensemble de définition.
\item Les fonctions $\sqrt{x}$, $\sin x$, $\cos x$, $\exp$, $\ln$ sont continues sur leur ensemble de définition.
\end{itemize}
\end{propriete}

\begin{definition}[Prolongement par continuité]
Si $f$ n'est pas définie en $x_0$ mais admet une limite finie $\ell$ en $x_0$, la fonction $\tilde{f}$ définie par
\[ \tilde{f}(x)=\begin{cases} f(x) & \text{si } x\neq x_0 \\ \ell & \text{si } x=x_0 \end{cases} \]
est continue en $x_0$. C'est le \textbf{prolongement par continuité} de $f$ en $x_0$.
\end{definition}

\begin{experience}[Découverte : Asymptotes et limites]
\begin{enumerate}
\item Soit $f(x)=\frac{1}{x}$. Tracer sa courbe. Que se passe-t-il près de $x=0$ ? Et pour $x$ très grand ?
\item On appelle \textbf{asymptote verticale} d'équation $x=a$ une droite que la courbe approche indéfiniment quand $x\to a$ (limite infinie).
\item On appelle \textbf{asymptote horizontale} d'équation $y=\ell$ une droite que la courbe approche quand $x\to\pm\infty$ (limite finie $\ell$).
\item On appelle \textbf{asymptote oblique} d'équation $y=ax+b$ ($a\neq 0$) une droite telle que $\lim_{x\to\pm\infty}(f(x)-(ax+b))=0$.
\end{enumerate}
\end{experience}

\begin{savaistu}[Le saviez-vous ? Le pont de la continuité]
En 1821, Cauchy formalise la définition $\varepsilon-\eta$ pour fonder le calcul infinitésimal. Avant lui, Euler manipulait les « infiniment petits » sans rigueur. Au Cameroun, la continuité modélise des phénomènes réels : le débit d'eau dans une canalisation (CAMWATER), la tension électrique (ENEO), ou le nombre d'abonnés MTN/Orange qui varie « sans saut » dans le temps... sauf lors d'une coupure (discontinuité) !
\end{savaistu}

\courssub{Je vérifie mes connaissances}

\begin{exercice}[QCM : une seule bonne réponse]
Pour chaque question, choisir la bonne réponse (aucune justification n'est demandée).
\begin{enumerate}
\item $\displaystyle\lim_{x\to 3}\frac{x^2-9}{x-3}$ est égale à :
\begin{enumerate}
\item[a)] $0$ \qquad b) $6$ \qquad c) $+\infty$ \qquad d) la limite n'existe pas.
\end{enumerate}
\item $\displaystyle\lim_{x\to +\infty}\frac{3x^2-x+1}{2x^2+5}$ est égale à :
\begin{enumerate}
\item[a)] $+\infty$ \qquad b) $0$ \qquad c) $\dfrac{3}{2}$ \qquad d) $\dfrac{1}{5}$.
\end{enumerate}
\item La fonction $f$ définie par $f(x)=\dfrac{x^2-4}{x-2}$ pour $x\neq 2$ :
\begin{enumerate}
\item[a)] est continue en $2$ \qquad b) admet une limite finie en $2$ égale à $4$ \qquad c) n'a pas de limite en $2$.
\end{enumerate}
\item Si $\displaystyle\lim_{x\to 1^-}f(x)=2$ et $\displaystyle\lim_{x\to 1^+}f(x)=5$, alors :
\begin{enumerate}
\item[a)] $f$ est continue en $1$ \qquad b) $\displaystyle\lim_{x\to 1}f(x)$ n'existe pas \qquad c) $f(1)=2$.
\end{enumerate}
\end{enumerate}
\end{exercice}

\begin{exercice}[Vrai ou faux, en justifiant]
Dire si chacune des affirmations suivantes est vraie ou fausse, en justifiant la réponse par un argument ou un contre-exemple.
\begin{enumerate}
\item « Si $\displaystyle\lim_{x\to 2}f(x)=3$, alors $f(2)=3$. »
\item « Si $f$ est continue en $2$, alors $\displaystyle\lim_{x\to 2}f(x)=f(2)$. »
\item « Toute fonction définie sur $\mathbb{R}$ est continue sur $\mathbb{R}$. »
\item « Si $\displaystyle\lim_{x\to +\infty}f(x)=+\infty$ et $\displaystyle\lim_{x\to +\infty}g(x)=+\infty$, alors $\displaystyle\lim_{x\to +\infty}\big(f(x)-g(x)\big)=0$. »
\end{enumerate}
\end{exercice}

\begin{exercice}[Définitions et vocabulaire]
\begin{enumerate}
\item Compléter la définition : « On dit que $f$ est \cle{continue} en $x_0$ lorsque $f$ est définie en $x_0$ et que \ldots »
\item Donner la définition de la limite à droite de $f$ en $x_0$, notée $\displaystyle\lim_{x\to x_0^+}f(x)$.
\item Citer les quatre formes indéterminées principales rencontrées dans le calcul de limites algébriques.
\item Donner les limites usuelles suivantes : $\displaystyle\lim_{x\to 0}\frac{\sqrt{x+1}-1}{x}$ et $\displaystyle\lim_{x\to +\infty}\frac{P(x)}{Q(x)}$ où $P$ et $Q$ sont deux polynômes de même degré (on précisera la règle utilisée).
\end{enumerate}
\end{exercice}

\begin{exercice}[Calculs directs]
Calculer les limites suivantes en utilisant les opérations sur les limites et les limites usuelles.
\begin{enumerate}
\item $\displaystyle\lim_{x\to +\infty}\left(2x^2-3x+7\right)$ ;
\item $\displaystyle\lim_{x\to -\infty}\frac{4x-1}{x+3}$ ;
\item $\displaystyle\lim_{x\to 5}\frac{x^2+1}{x-3}$ ;
\item $\displaystyle\lim_{x\to +\infty}\left(\frac{1}{x}+3\right)\left(x^2-2\right)$.
\end{enumerate}
\end{exercice}

\courssub{Je m'entraîne}

\begin{exercice}[Identifier la forme puis calculer]
Pour chacune des limites suivantes, identifier d'abord la forme (déterminée ou indéterminée), puis calculer la limite.
\begin{enumerate}
\item $\displaystyle\lim_{x\to 3}\frac{x^2-9}{x-3}$ ;
\item $\displaystyle\lim_{x\to 0}\frac{\sqrt{x+4}-2}{x}$ ;
\item $\displaystyle\lim_{x\to +\infty}\frac{3x^2-x+1}{2x^2+5}$ ;
\item $\displaystyle\lim_{x\to +\infty}\left(\sqrt{x^2+2x}-x\right)$.
\end{enumerate}
\end{exercice}

\begin{exercice}[Prolongement par continuité]
Soit $f$ la fonction définie sur $\mathbb{R}\setminus\{2\}$ par $f(x)=\dfrac{x^2-4}{x-2}$.
\begin{enumerate}
\item Déterminer l'ensemble de définition de $f$.
\item Calculer la limite de $f$ en $2$.
\item En déduire le prolongement par continuité de $f$ en $2$ : on précisera l'expression de la fonction prolongée $\tilde{f}$ et son ensemble de définition.
\end{enumerate}
\end{exercice}

\begin{exercice}[Continuité d'une fonction définie par morceaux]
Soit $a$ un réel et $f$ la fonction définie sur $\mathbb{R}$ par :
\[ f(x)=\begin{cases} 2x+a & \text{si } x<1 \\ x^2+2 & \text{si } x\geq 1 \end{cases} \]
\begin{enumerate}
\item Calculer $f(1)$.
\item Calculer $\displaystyle\lim_{x\to 1^+}f(x)$ et $\displaystyle\lim_{x\to 1^-}f(x)$ en fonction de $a$.
\item Étudier la continuité de $f$ en $x_0=1$ lorsque $a=0$.
\item Déterminer la valeur du réel $a$ pour laquelle $f$ est continue en $1$.
\end{enumerate}
\end{exercice}

\begin{exercice}[Quantité conjuguée]
Soit $h$ la fonction définie pour $x\geq 1$ par $h(x)=x-\sqrt{x^2-1}$.
\begin{enumerate}
\item Vérifier que le calcul direct de la limite de $h$ en $+\infty$ conduit à une forme indéterminée.
\item En multipliant par la quantité conjuguée, montrer que pour tout $x\geq 1$ :
\[ h(x)=\frac{1}{x+\sqrt{x^2-1}}. \]
\item En déduire que $h$ admet une limite finie en $+\infty$ et la calculer.
\item Interpréter graphiquement ce résultat pour la courbe représentative de $h$.
\end{enumerate}
\end{exercice}

\courssub{Je me prépare au Bac}

\begin{exercice}[Type Bac : étude complète d'une fonction rationnelle]
Une entreprise de transformation de manioc à Bafoussam modélise son coût moyen de production, en milliers de francs CFA par tonne, par la fonction
\[ f(x)=\frac{2x^2-3x+1}{x^2-4} \]
où $x$ désigne le nombre de tonnes produites par jour.
\begin{enumerate}
\item Déterminer l'ensemble de définition $D_f$ de $f$.
\item Calculer les limites de $f$ en $-\infty$ et en $+\infty$. Interpréter graphiquement le résultat.
\item Calculer les limites de $f$ à gauche et à droite en $-2$ et en $2$.
\item En déduire toutes les asymptotes à la courbe représentative $(\mathcal{C})$ de $f$, en précisant leur nature et leur équation.
\item La fonction $f$ est-elle continue sur $D_f$ ? Justifier. Est-elle continue sur $\mathbb{R}$ ?
\item Dans le contexte de l'entreprise, expliquer pourquoi une production journalière voisine de $2$ tonnes pose problème.
\end{enumerate}
\end{exercice}

\begin{popfigure}
\begin{tikzpicture}[scale=0.8, domain=-5:5]
\draw[popline, thick, ->] (-5.5,0) -- (5.5,0) node[right] {$x$};
\draw[popline, thick, ->] (0,-5.5) -- (0,5.5) node[above] {$y$};
% Asymptotes
\draw[popmuted, dashed] (-2,-5) -- (-2,5) node[above left] {$x=-2$};
\draw[popmuted, dashed] (2,-5) -- (2,5) node[above right] {$x=2$};
\draw[popmuted, dashed] (-5,2) -- (5,2) node[right] {$y=2$};
% Branche gauche x < -2
\draw[popblue, thick, samples=100] plot (\x, { (2*\x*\x - 3*\x + 1)/(\x*\x - 4) }) node[above left] {$\mathcal{C}_f$};
% Branche milieu -2 < x < 2
\draw[popblue, thick, samples=100] plot (\x, { (2*\x*\x - 3*\x + 1)/(\x*\x - 4) });
% Branche droite x > 2
\draw[popblue, thick, samples=100] plot (\x, { (2*\x*\x - 3*\x + 1)/(\x*\x - 4) });
\foreach \x in {-2,2} \draw (\x,2pt) -- (\x,-2pt) node[below] {\x};
\draw (0,2pt) -- (0,-2pt) node[below left] {0};
\draw (2pt,-0.25) -- (-2pt,-0.25) node[left] {$-0,25$}; % f(0)
\end{tikzpicture}
\legende{Courbe de $f(x)=\frac{2x^2-3x+1}{x^2-4}$ avec ses asymptotes (pointillés).}
\end{popfigure}

\begin{exercice}[Type Bac : lecture graphique]
Le réseau d'une agence Orange de Douala enregistre le débit de connexion $g(t)$ (en mégabits par seconde) en fonction du temps $t$ (en heures). La courbe représentative de $g$ admet la droite d'équation $x=-1$ comme asymptote verticale, la droite d'équation $y=2$ comme asymptote horizontale en $+\infty$, et présente un point de discontinuité en $x=3$ (coupure du réseau) : la courbe « saute » d'une valeur à une autre en $x=3$, avec $g(3)=6$.
\begin{enumerate}
\item Lire graphiquement les limites suivantes : $\displaystyle\lim_{x\to -1^-}g(x)$, $\displaystyle\lim_{x\to -1^+}g(x)$, $\displaystyle\lim_{x\to +\infty}g(x)$ et $\displaystyle\lim_{x\to -\infty}g(x)$ (on admettra que la courbe admet aussi une asymptote horizontale d'équation $y=0$ en $-\infty$).
\item La fonction $g$ admet-elle une limite en $-1$ ? Justifier.
\item On donne $\displaystyle\lim_{x\to 3^-}g(x)=4$ et $\displaystyle\lim_{x\to 3^+}g(x)=1$.
\begin{enumerate}
\item La fonction $g$ admet-elle une limite en $3$ ?
\item La fonction $g$ est-elle continue en $3$ ? Justifier en utilisant la définition de la continuité.
\end{enumerate}
\item Dresser le tableau récapitulatif des asymptotes de la courbe de $g$.
\item Expliquer, en langage courant, ce que représente la discontinuité en $t=3$ pour les abonnés de l'agence.
\end{enumerate}
\end{exercice}

\begin{popfigure}
\begin{tikzpicture}[scale=0.7, domain=-3:6]
\draw[popline, thick, ->] (-3.5,0) -- (6.5,0) node[right] {$t$ (heures)};
\draw[popline, thick, ->] (0,-0.5) -- (0,7) node[above] {$g(t)$ (Mbps)};
% Asymptotes
\draw[popmuted, dashed] (-1,-0.5) -- (-1,7) node[above] {$x=-1$};
\draw[popmuted, dashed] (-3,2) -- (6,2) node[right] {$y=2$};
\draw[popmuted, dashed] (-3,0) -- (6,0) node[right] {$y=0$};
% Courbe gauche (x < -1) : tend vers +inf en -1-, vers 0 en -inf
\draw[popblue, thick, samples=100] plot (\x, { 2/(1+exp(-\x)) + 3/(1+exp(\x+1)) }); % approximation shape
% Point discontinuité x=3
\draw[popblue, thick] (-1, 0.5) .. controls (0,1) and (2,3.5) .. (3,4);
\fill[popblue] (3,4) circle (2pt) node[above left] {$\lim_{3^-}=4$};
\fill[poporange] (3,6) circle (2pt) node[above right] {$g(3)=6$};
\fill[popgreen] (3,1) circle (2pt) node[below right] {$\lim_{3^+}=1$};
% Courbe droite (x > 3) : tend vers 2 en +inf
\draw[popblue, thick] (3,1) .. controls (4,1.5) and (5,1.8) .. (6,1.95);
\foreach \x in {-1,3} \draw (\x,2pt) -- (\x,-2pt) node[below] {\x};
\draw (2pt,2) -- (-2pt,2) node[left] {2};
\draw (2pt,1) -- (-2pt,1) node[left] {1};
\draw (2pt,4) -- (-2pt,4) node[left] {4};
\draw (2pt,6) -- (-2pt,6) node[left] {6};
\end{tikzpicture}
\legende{Schéma qualitatif de la courbe de $g$ (débit Orange).}
\end{popfigure}

\begin{exercice}[Type Bac : limites, continuité et paramètre]
Une société de transport par bus entre Yaoundé et Bafoussam facture le prix d'un billet (en milliers de francs CFA) selon la distance parcourue $x$ (en dizaines de kilomètres) par la fonction :
\[ f(x)=\begin{cases} \dfrac{x^2-1}{x-1} & \text{si } x>1 \\[6pt] ax+1 & \text{si } x\leq 1 \end{cases} \]
où $a$ est un réel fixé par la direction.
\begin{enumerate}
\item Déterminer l'ensemble de définition de $f$.
\item Montrer que pour tout $x>1$, $f(x)=x+1$, puis calculer $\displaystyle\lim_{x\to 1^+}f(x)$.
\item Calculer $f(1)$ et $\displaystyle\lim_{x\to 1^-}f(x)$ en fonction de $a$.
\item Déterminer la valeur de $a$ pour laquelle $f$ est continue en $1$. Que représente cette continuité pour le client dont le trajet franchit la distance de $10$ km ?
\item On suppose désormais $a=1$.
\begin{enumerate}
\item Calculer $\displaystyle\lim_{x\to +\infty}\big(f(x)-x\big)$ et interpréter graphiquement le résultat.
\item Calculer $\displaystyle\lim_{x\to +\infty}\frac{f(x)}{x}$.
\end{enumerate}
\item La fonction $f$ (avec $a=1$) est-elle continue sur $\mathbb{R}$ tout entier ? Justifier.
\end{enumerate}
\end{exercice}



\newpage

\coursec{Corrigés des exercices}

\coursec{Exercices corrigés : Limites et continuité}

\begin{corrige}[Exercice 1 — QCM]
\begin{enumerate}
\item \textbf{Réponse b) :} $6$. Pour $x\neq 3$, $\dfrac{x^2-9}{x-3}=\dfrac{(x-3)(x+3)}{x-3}=x+3$, donc $\lim_{x\to 3}\dfrac{x^2-9}{x-3}=\lim_{x\to 3}(x+3)=6$.
\item \textbf{Réponse c) :} $\dfrac{3}{2}$. Numérateur et dénominateur sont de même degré $2$ : la limite en $+\infty$ est le quotient des coefficients directeurs, soit $\dfrac{3}{2}$.
\item \textbf{Réponse b) :} pour $x\neq 2$, $f(x)=\dfrac{(x-2)(x+2)}{x-2}=x+2$, donc $\lim_{x\to 2}f(x)=4$. Mais $f$ n'est pas définie en $2$ : elle n'est donc pas continue en $2$ (réponse a) fausse), et la limite existe (réponse c) fausse).
\item \textbf{Réponse b) :} les limites à gauche ($2$) et à droite ($5$) sont différentes, donc $\lim_{x\to 1}f(x)$ n'existe pas, d'après le théorème de la limite bilatère.
\end{enumerate}
\textbf{Point de vigilance :} dans la question 3, l'existence d'une limite finie ne suffit pas pour la continuité : il faut en plus que $f$ soit définie au point et que la limite égale la valeur.
\end{corrige}

\begin{corrige}[Exercice 2 — Vrai ou faux]
\begin{enumerate}
\item \textbf{Faux.} La limite en $x_0$ ne dépend jamais de la valeur $f(x_0)$, qui peut être différente ou même ne pas exister. Contre-exemple : $f(x)=3$ si $x\neq 2$ et $f(2)=0$. Alors $\lim_{x\to 2}f(x)=3$ mais $f(2)=0\neq 3$.
\item \textbf{Vrai.} C'est exactement la définition de la continuité en un point : $f$ continue en $2 \iff \lim_{x\to 2}f(x)=f(2)$.
\item \textbf{Faux.} Être définie partout ne garantit pas la continuité. Contre-exemple : la fonction partie entière $E(x)$ est définie sur $\mathbb{R}$ mais discontinue en tout entier (saut de $1$). Autre exemple : $f(x)=\dfrac{1}{x}$ si $x\neq 0$ et $f(0)=0$ est définie sur $\mathbb{R}$ mais discontinue en $0$.
\item \textbf{Faux.} C'est une forme indéterminée $\infty-\infty$ : tout peut arriver. Contre-exemple : $f(x)=x+1$ et $g(x)=x$ tendent toutes deux vers $+\infty$, mais $f(x)-g(x)=1\xrightarrow[x\to+\infty]{}1\neq 0$. (Avec $f(x)=2x$, $g(x)=x$, la différence tend vers $+\infty$ ; avec $f(x)=x$, $g(x)=x$, elle vaut $0$.)
\end{enumerate}
\textbf{Point de vigilance :} pour réfuter une affirmation, un seul contre-exemple suffit, mais il faut le vérifier explicitement (calculer les limites annoncées).
\end{corrige}

\begin{corrige}[Exercice 3 — Définitions et vocabulaire]
\begin{enumerate}
\item « On dit que $f$ est \cle{continue} en $x_0$ lorsque $f$ est définie en $x_0$ et que $\displaystyle\lim_{x\to x_0}f(x)=f(x_0)$. »
\item On dit que $f$ admet la limite $\ell$ à droite en $x_0$, et on note $\lim_{x\to x_0^+}f(x)=\ell$, si :
\[ \forall \varepsilon>0,\ \exists \eta>0,\ \forall x\in D_f,\ x_0<x<x_0+\eta \implies |f(x)-\ell|<\varepsilon. \]
\item Les quatre formes indéterminées algébriques principales sont : $\dfrac{0}{0}$, $\dfrac{\infty}{\infty}$, $0\times\infty$ et $\infty-\infty$.
\item $\displaystyle\lim_{x\to 0}\frac{\sqrt{x+1}-1}{x}=\frac{1}{2}$ (limite usuelle admise). Si $P$ et $Q$ sont deux polynômes de même degré, de coefficients directeurs respectifs $a_n$ et $b_n$, alors $\displaystyle\lim_{x\to +\infty}\frac{P(x)}{Q(x)}=\frac{a_n}{b_n}$ : en l'infini, une fraction rationnelle se comporte comme le quotient de ses termes de plus haut degré.
\end{enumerate}
\end{corrige}

\begin{corrige}[Exercice 4 — Calculs directs]
\begin{enumerate}
\item $2x^2-3x+7$ est un polynôme : en $+\infty$, il se comporte comme son terme de plus haut degré $2x^2$. Or $\lim_{x\to +\infty}2x^2=+\infty$, donc $\boxed{\lim_{x\to +\infty}(2x^2-3x+7)=+\infty}$.
\item Fraction rationnelle de degrés égaux ($1$ et $1$) : $\lim_{x\to -\infty}\dfrac{4x-1}{x+3}=\dfrac{4}{1}=\boxed{4}$.
\item Le dénominateur ne s'annule pas en $5$ ($5-3=2\neq 0$) : les opérations s'appliquent directement. $\lim_{x\to 5}\dfrac{x^2+1}{x-3}=\dfrac{25+1}{5-3}=\dfrac{26}{2}=\boxed{13}$.
\item $\lim_{x\to +\infty}\left(\dfrac{1}{x}+3\right)=0+3=3$ et $\lim_{x\to +\infty}(x^2-2)=+\infty$. Par produit : $3\times(+\infty)=\boxed{+\infty}$ (forme déterminée : le premier facteur a une limite finie non nulle).
\end{enumerate}
\end{corrige}

\begin{corrige}[Exercice 5 — Identifier la forme puis calculer]
\begin{enumerate}
\item \textbf{Forme $\dfrac{0}{0}$} (numérateur $9-9=0$, dénominateur $0$). On factorise :
\[ \frac{x^2-9}{x-3}=\frac{(x-3)(x+3)}{x-3}=x+3 \quad (x\neq 3), \]
donc $\displaystyle\lim_{x\to 3}\frac{x^2-9}{x-3}=\boxed{6}$.
\item \textbf{Forme $\dfrac{0}{0}$} ($\sqrt{4}-2=0$). On multiplie par la quantité conjuguée :
\[ \frac{\sqrt{x+4}-2}{x}=\frac{(\sqrt{x+4}-2)(\sqrt{x+4}+2)}{x(\sqrt{x+4}+2)}=\frac{(x+4)-4}{x(\sqrt{x+4}+2)}=\frac{1}{\sqrt{x+4}+2}, \]
d'où $\displaystyle\lim_{x\to 0}\frac{\sqrt{x+4}-2}{x}=\frac{1}{\sqrt{4}+2}=\boxed{\frac{1}{4}}$.
\item \textbf{Forme $\dfrac{\infty}{\infty}$}. On divise numérateur et dénominateur par $x^2$ (terme dominant) :
\[ \frac{3x^2-x+1}{2x^2+5}=\frac{3-\frac{1}{x}+\frac{1}{x^2}}{2+\frac{5}{x^2}}\xrightarrow[x\to +\infty]{}\boxed{\frac{3}{2}}, \]
car $\frac{1}{x}\to 0$ et $\frac{1}{x^2}\to 0$.
\item \textbf{Forme $\infty-\infty$}. Quantité conjuguée :
\[ \sqrt{x^2+2x}-x=\frac{(\sqrt{x^2+2x}-x)(\sqrt{x^2+2x}+x)}{\sqrt{x^2+2x}+x}=\frac{2x}{\sqrt{x^2+2x}+x}. \]
Pour $x>0$, $\sqrt{x^2+2x}=x\sqrt{1+\frac{2}{x}}$, donc l'expression vaut $\dfrac{2x}{x\sqrt{1+\frac{2}{x}}+x}=\dfrac{2}{\sqrt{1+\frac{2}{x}}+1}\xrightarrow[x\to +\infty]{}\dfrac{2}{2}=\boxed{1}$.
\end{enumerate}
\textbf{Point de vigilance :} dans la question 4, pour $x>0$ on a $\sqrt{x^2}=x$ (et non $-x$) : c'est ce qui permet de factoriser $x$ sous la racine. Pour $x\to -\infty$, le calcul serait différent.
\end{corrige}

\begin{corrige}[Exercice 6 — Prolongement par continuité]
\begin{enumerate}
\item $f$ est définie partout sauf là où le dénominateur s'annule : $x-2=0\iff x=2$. Donc $D_f=\mathbb{R}\setminus\{2\}$.
\item Pour $x\neq 2$ : $f(x)=\dfrac{x^2-4}{x-2}=\dfrac{(x-2)(x+2)}{x-2}=x+2$. Par conséquent :
\[ \lim_{x\to 2}f(x)=\lim_{x\to 2}(x+2)=\boxed{4}. \]
\item La limite en $2$ existe et est finie : $f$ est donc prolongeable par continuité en $2$. Le prolongement est la fonction $\tilde{f}$ définie sur $\mathbb{R}$ par :
\[ \tilde{f}(x)=\begin{cases} \dfrac{x^2-4}{x-2} & \text{si } x\neq 2 \\[4pt] 4 & \text{si } x=2 \end{cases} \]
c'est-à-dire simplement $\tilde{f}(x)=x+2$ pour tout $x\in\mathbb{R}$. C'est une fonction affine, donc continue sur $\mathbb{R}$ tout entier.
\end{enumerate}
\textbf{Point de vigilance :} dire que $f$ est prolongeable par continuité exige de vérifier que la limite est \emph{finie}. Une limite infinie (asymptote verticale) interdit tout prolongement continu.
\end{corrige}

\begin{corrige}[Exercice 7 — Continuité d'une fonction définie par morceaux]
\begin{enumerate}
\item Comme $1\geq 1$, on utilise la deuxième branche : $f(1)=1^2+2=\boxed{3}$.
\item \textbf{À droite :} pour $x>1$, $f(x)=x^2+2$, donc $\displaystyle\lim_{x\to 1^+}f(x)=1+2=3$.
\textbf{À gauche :} pour $x<1$, $f(x)=2x+a$, donc $\displaystyle\lim_{x\to 1^-}f(x)=2+a$.
\item Si $a=0$ : $\lim_{x\to 1^-}f(x)=2$ et $\lim_{x\to 1^+}f(x)=3$. Les limites latérales sont différentes : d'après le théorème de la limite bilatère, $\lim_{x\to 1}f(x)$ n'existe pas. Donc \textbf{$f$ n'est pas continue en $1$} (discontinuité de première espèce, ou « saut », d'amplitude $1$).
\item $f$ est continue en $1$ si et seulement si les trois conditions sont réunies : $f$ définie en $1$ (oui, $f(1)=3$), limite existante, et égalité. Il faut et il suffit :
\[ \lim_{x\to 1^-}f(x)=\lim_{x\to 1^+}f(x)=f(1)\iff 2+a=3\iff \boxed{a=1}. \]
\end{enumerate}
\textbf{Point de vigilance :} la continuité exige la triple égalité limite à gauche = limite à droite = valeur au point. Vérifier seulement l'égalité des deux limites latérales est insuffisant en toute rigueur (même si ici $f(1)=3$ coïncide avec la limite à droite).
\end{corrige}

\begin{corrige}[Exercice 8 — Quantité conjuguée]
\begin{enumerate}
\item $\lim_{x\to +\infty}x=+\infty$ et $\lim_{x\to +\infty}\sqrt{x^2-1}=+\infty$ : le calcul direct donne la forme indéterminée $\infty-\infty$, dont on ne peut rien conclure sans transformation.
\item Pour tout $x\geq 1$, $x+\sqrt{x^2-1}>0$, on peut donc multiplier par la quantité conjuguée :
\begin{align*}
h(x)&=(x-\sqrt{x^2-1})\times\frac{x+\sqrt{x^2-1}}{x+\sqrt{x^2-1}}\\
&=\frac{x^2-\left(\sqrt{x^2-1}\right)^2}{x+\sqrt{x^2-1}}=\frac{x^2-(x^2-1)}{x+\sqrt{x^2-1}}=\frac{1}{x+\sqrt{x^2-1}}.
\end{align*}
\item Quand $x\to +\infty$ : $x\to +\infty$ et $\sqrt{x^2-1}\to +\infty$, donc par somme $x+\sqrt{x^2-1}\to +\infty$, puis par quotient :
\[ \lim_{x\to +\infty}h(x)=\lim_{x\to +\infty}\frac{1}{x+\sqrt{x^2-1}}=\boxed{0}. \]
La limite est donc bien finie.
\item \textbf{Interprétation graphique :} la courbe représentative de $h$ admet la droite d'équation $y=0$ (l'axe des abscisses) comme \textbf{asymptote horizontale} en $+\infty$ : pour les grandes valeurs de $x$, la courbe se rapproche indéfiniment de l'axe des abscisses tout en restant au-dessus (car $h(x)>0$).
\end{enumerate}
\textbf{Point de vigilance :} l'identité utilisée est $(a-b)(a+b)=a^2-b^2$ avec $\left(\sqrt{x^2-1}\right)^2=x^2-1$ : ne pas oublier les parenthèses, sinon on obtient $-x^2-1$ au lieu de $-(x^2-1)$.
\end{corrige}

\begin{corrige}[Exercice 9 — Type Bac : fonction rationnelle (coût du manioc)]
\begin{enumerate}
\item $f$ est définie lorsque $x^2-4\neq 0$, c'est-à-dire $x\neq -2$ et $x\neq 2$. Donc $D_f=\mathbb{R}\setminus\{-2\,;\,2\}$.
\item Numérateur et dénominateur sont de même degré $2$, de coefficients directeurs $2$ et $1$ :
\[ \lim_{x\to +\infty}f(x)=\lim_{x\to +\infty}\frac{2x^2}{x^2}=\boxed{2}\quad\text{et de même}\quad\lim_{x\to -\infty}f(x)=\boxed{2}. \]
\textbf{Interprétation graphique :} la courbe $(\mathcal{C})$ admet la droite d'équation $y=2$ comme asymptote horizontale en $+\infty$ et en $-\infty$. Concrètement, pour de très grosses productions, le coût moyen se stabilise autour de $2$ milliers de francs CFA par tonne.
\item \textbf{Étude en $-2$.} Valeur du numérateur en $-2$ : $2(4)-3(-2)+1=8+6+1=15>0$.
Signe du dénominateur : $x^2-4>0$ si $x<-2$ et $x^2-4<0$ si $-2<x<2$.
\[ \lim_{x\to -2^-}f(x)=\frac{15}{0^+}=\boxed{+\infty}\ ;\qquad \lim_{x\to -2^+}f(x)=\frac{15}{0^-}=\boxed{-\infty}. \]
\textbf{Étude en $2$.} Valeur du numérateur en $2$ : $2(4)-3(2)+1=8-6+1=3>0$.
\[ \lim_{x\to 2^-}f(x)=\frac{3}{0^-}=\boxed{-\infty}\ ;\qquad \lim_{x\to 2^+}f(x)=\frac{3}{0^+}=\boxed{+\infty}. \]
\item \textbf{Asymptotes verticales :} les droites $x=-2$ et $x=2$ (limites infinies en ces points).
\textbf{Asymptote horizontale :} la droite $y=2$ (limites finies égales à $2$ en $\pm\infty$).
Il n'y a pas d'asymptote oblique (la présence d'une asymptote horizontale l'exclut).
\item $f$ est une fonction rationnelle : elle est donc \textbf{continue sur son ensemble de définition} $D_f=\mathbb{R}\setminus\{-2\,;\,2\}$. En revanche, $f$ n'est pas définie en $-2$ ni en $2$ (et les limites y sont infinies, donc aucun prolongement possible) : \textbf{$f$ n'est pas continue sur $\mathbb{R}$}.
\item Lorsque la production $x$ s'approche de $2$ tonnes par valeur supérieure, le coût moyen $f(x)$ tend vers $+\infty$ : le coût par tonne \emph{explose}. Le dénominateur $x^2-4$ devient quasi nul alors que le numérateur vaut environ $3$ : le modèle traduit un seuil critique de l'usine (saturation des capacités, pannes, heures supplémentaires...). L'entreprise doit donc éviter une production journalière voisine de $2$ tonnes et se placer nettement au-dessus de ce seuil, où le coût moyen redevient raisonnable et proche de $2$ milliers de francs.
\end{enumerate}
\textbf{Barème indicatif (sur 6 pts) :} Q1 : $0,5$ ; Q2 : $1$ ($0,5$ par limite) ; Q3 : $1,5$ ; Q4 : $1$ ; Q5 : $1$ ; Q6 : $1$.
\textbf{Points de vigilance Bac :} (i) pour les limites en $-2$ et $2$, il est \emph{indispensable} de justifier le signe du dénominateur de chaque côté ($0^+$ ou $0^-$) : écrire directement « $\pm\infty$ » sans étude de signe coûte des points ; (ii) citer la propriété « toute fonction rationnelle est continue sur son ensemble de définition » ; (iii) ne pas confondre « continue sur $D_f$ » et « continue sur $\mathbb{R}$ ».
\end{corrige}

\begin{popfigure}
\begin{tikzpicture}[scale=0.65, domain=-5:5]
\draw[popdark, ->] (-5.5,0) -- (5.5,0) node[right] {$x$};
\draw[popdark, ->] (0,-5.5) -- (0,5.5) node[above] {$y$};
% Asymptotes
\draw[popblue, dashed, thick] (-2,-5.5) -- (-2,5.5) node[above] {$x=-2$};
\draw[popblue, dashed, thick] (2,-5.5) -- (2,5.5) node[above] {$x=2$};
\draw[popblue, dashed, thick] (-5.5,2) -- (5.5,2) node[right] {$y=2$};
% Curve branches (plot of f(x) = (2x^2-3x+1)/(x^2-4))
\draw[poporange, very thick, samples=200, domain=-5:-2.15] plot (\x, {(2*\x*\x - 3*\x + 1)/(\x*\x - 4)});
\draw[poporange, very thick, samples=200, domain=-1.9:1.9] plot (\x, {(2*\x*\x - 3*\x + 1)/(\x*\x - 4)});
\draw[poporange, very thick, samples=200, domain=2.15:5] plot (\x, {(2*\x*\x - 3*\x + 1)/(\x*\x - 4)});
% Points of interest
\fill[poporange] (-1,0) circle (2pt) node[below left] {$(-1,0)$};
\fill[poporange] (0.5,0) circle (2pt) node[above right] {$(\frac{1}{2},0)$};
\fill[poporange] (0,-0.25) circle (2pt) node[left] {$(0,-\frac{1}{4})$};
\end{tikzpicture}
\legende{Courbe de $f(x)=\frac{2x^2-3x+1}{x^2-4}$ et ses asymptotes (verticales $x=\pm2$, horizontale $y=2$).}
\end{popfigure}

\begin{corrige}[Exercice 10 — Type Bac : lecture graphique (débit Orange)]
\begin{enumerate}
\item Lecture directe sur la courbe et ses asymptotes :
\begin{itemize}
\item Asymptote verticale $x=-1$ : $\displaystyle\lim_{x\to -1^-}g(x)=+\infty$ et $\displaystyle\lim_{x\to -1^+}g(x)=-\infty$ (la courbe s'éloigne indéfiniment de part et d'autre de la droite) ;
\item Asymptote horizontale $y=2$ en $+\infty$ : $\displaystyle\lim_{x\to +\infty}g(x)=2$ ;
\item Asymptote horizontale $y=0$ en $-\infty$ : $\displaystyle\lim_{x\to -\infty}g(x)=0$.
\end{itemize}
\item \textbf{Non.} Les limites à gauche et à droite en $-1$ sont infinies et de signes opposés : elles ne sont ni égales ni finies. D'après le théorème de la limite bilatère, $g$ n'admet pas de limite en $-1$.
\item \begin{enumerate}
\item \textbf{Non.} $\displaystyle\lim_{x\to 3^-}g(x)=4\neq 1=\displaystyle\lim_{x\to 3^+}g(x)$ : les limites latérales étant différentes, la limite de $g$ en $3$ n'existe pas.
\item \textbf{Non.} La continuité en $3$ exigerait $\lim_{x\to 3}g(x)=g(3)$. Or la limite en $3$ n'existe pas (question précédente) : la condition de la définition n'est pas satisfaite. De plus, même chaque limite latérale ($4$ et $1$) diffère de $g(3)=6$. Il s'agit d'une \textbf{discontinuité de première espèce} (un « saut »).
\end{enumerate}
\item Tableau récapitulatif :
\begin{center}
\begin{tabularx}{\linewidth}{|l|X|X|}\hline
\rowcolor{popblueL}
\textbf{Asymptote} & \textbf{Équation} & \textbf{Justification} \\ \hline
Verticale & $x=-1$ & Limites infinies en $-1$ à gauche et à droite \\ \hline
Horizontale en $+\infty$ & $y=2$ & $\lim_{x\to +\infty}g(x)=2$ \\ \hline
Horizontale en $-\infty$ & $y=0$ & $\lim_{x\to -\infty}g(x)=0$ \\ \hline
\end{tabularx}
\end{center}
\item En langage courant : à l'instant $t=3$ h, le débit chute brutalement d'environ $4$ Mbps à environ $1$ Mbps, sans transition progressive. Pour les abonnés, cela se traduit par une \textbf{coupure soudaine du service} : appels vidéo interrompus, téléchargements bloqués, puis une reprise avec un débit fortement dégradé. La valeur $g(3)=6$ Mbps enregistrée à l'instant exact est un pic aberrant de la mesure au moment de l'incident, sans signification pour la qualité réelle de la connexion.
\end{enumerate}
\textbf{Barème indicatif (sur 5 pts) :} Q1 : $1,5$ ; Q2 : $0,5$ ; Q3 : $1,5$ ; Q4 : $0,5$ ; Q5 : $1$.
\textbf{Points de vigilance Bac :} (i) en lecture graphique, toujours relier chaque limite à l'asymptote correspondante (c'est la justification attendue) ; (ii) pour la continuité, citer explicitement la définition $\lim_{x\to 3}g(x)=g(3)$ et montrer qu'elle échoue ; (iii) ne pas confondre « pas de limite » et « limite infinie » : une limite infinie est une limite (non finie), ici les limites latérales finies et distinctes empêchent toute limite.
\end{corrige}

\begin{popfigure}
\begin{tikzpicture}[scale=0.7, domain=-4:5]
\draw[popdark, ->] (-4.5,0) -- (5.5,0) node[right] {$x$ (h)};
\draw[popdark, ->] (0,-1.5) -- (0,6.5) node[above] {$y$ (Mbps)};
% Asymptotes
\draw[popblue, dashed, thick] (-1,-1.5) -- (-1,6.5) node[above] {$x=-1$};
\draw[popblue, dashed, thick] (-4.5,2) -- (5.5,2) node[right] {$y=2$};
\draw[popblue, dashed, thick] (-4.5,0) -- (5.5,0) node[right] {$y=0$};
% Curve sketch (qualitative using smooth coordinates)
% Left branch: from y~0 at -inf, up to +inf at -1-
\draw[poporange, very thick] plot[smooth, tension=0.8] coordinates {(-4,0.5) (-3,1) (-2,2.5) (-1.5,5)};
% Right branch: from -inf at -1+, up to y=2 at +inf, with jump at 3
\draw[poporange, very thick] plot[smooth, tension=0.8] coordinates {(-0.5,-2) (0,0.5) (1,1.5) (2.5,3.5)};
\draw[poporange, very thick] plot[smooth, tension=0.8] coordinates {(3.5,1.2) (4,1.5) (5,1.8)};
% Jump at x=3
\fill[poporange] (3,4) circle (2pt); 
\fill[poporange] (3,1) circle (2pt); 
\fill[poppink] (3,6) circle (2pt) node[above right] {$g(3)=6$};
\draw[popmuted, dashed, thin] (3,0) -- (3,6.5);
% Labels for limits
\node[poporange, anchor=south east] at (2.9,4) {$\lim_{x\to 3^-}=4$};
\node[poporange, anchor=north east] at (2.9,1) {$\lim_{x\to 3^+}=1$};
\end{tikzpicture}
\legende{Représentation graphique schématique de la fonction $g$ (débit Orange) : asymptotes $x=-1$, $y=0$, $y=2$ ; discontinuité de première espèce en $x=3$.}
\end{popfigure}

\begin{corrige}[Exercice 11 — Type Bac : limites, continuité et paramètre (transport)]
\begin{enumerate}
\item Pour $x>1$, le dénominateur $x-1$ est non nul ; pour $x\leq 1$, l'expression $ax+1$ est définie pour tout réel. Donc $D_f=\boxed{\mathbb{R}}$.
\item Pour tout $x>1$ : $x^2-1=(x-1)(x+1)$, donc
\[ f(x)=\frac{(x-1)(x+1)}{x-1}=x+1. \]
Par suite : $\displaystyle\lim_{x\to 1^+}f(x)=\lim_{x\to 1^+}(x+1)=\boxed{2}$.
\item Comme $1\leq 1$, on utilise la deuxième branche : $f(1)=a\times 1+1=\boxed{a+1}$.
Pour $x<1$, $f(x)=ax+1$ (fonction affine, continue) : $\displaystyle\lim_{x\to 1^-}f(x)=\boxed{a+1}$.
\item $f$ est continue en $1$ si et seulement si $\lim_{x\to 1^-}f(x)=\lim_{x\to 1^+}f(x)=f(1)$, c'est-à-dire :
\[ a+1=2\iff \boxed{a=1}. \]
\textbf{Interprétation :} avec $a=1$, le tarif ne subit aucun saut au franchissement du seuil de $10$ km ($x=1$) : le prix passe continûment de $2$ milliers de francs (limite commune) à la valeur facturée. Le client dont le trajet franchit exactement $10$ km paie un prix cohérent avec ceux des trajets immédiatement plus courts ou plus longs : la tarification est équitable au seuil.
\item On suppose $a=1$. Alors $f(x)=x+1$ à la fois pour $x\leq 1$ et pour $x>1$ : $f(x)=x+1$ sur $\mathbb{R}$ entier.
\begin{enumerate}
\item $f(x)-x=1$ pour tout $x$, donc $\displaystyle\lim_{x\to +\infty}\big(f(x)-x\big)=\boxed{1}$.
\textbf{Interprétation graphique :} la distance verticale entre la courbe de $f$ et la droite d'équation $y=x$ reste constante, égale à $1$ : la courbe est la droite $y=x+1$, parallèle à $y=x$ (la droite $y=x+1$ est sa propre asymptote oblique).
\item $\displaystyle\lim_{x\to +\infty}\frac{f(x)}{x}=\lim_{x\to +\infty}\frac{x+1}{x}=\lim_{x\to +\infty}\left(1+\frac{1}{x}\right)=\boxed{1}$ : c'est le coefficient directeur de la direction asymptotique.
\end{enumerate}
\item \textbf{Oui.} Avec $a=1$, $f$ coïncide sur $\mathbb{R}$ avec la fonction affine $x\mapsto x+1$, qui est une fonction polynomiale, donc continue sur $\mathbb{R}$. (En particulier, la question 4 garantit la continuité au point de raccordement $x=1$, et chaque branche est continue de son côté.)
\end{enumerate}
\textbf{Barème indicatif (sur 6 pts) :} Q1 : $0,5$ ; Q2 : $1$ ; Q3 : $1$ ; Q4 : $1,5$ (dont $0,5$ pour l'interprétation) ; Q5 : $1$ ; Q6 : $1$.
\textbf{Points de vigilance Bac :} (i) la simplification $\frac{x^2-1}{x-1}=x+1$ n'est valable que pour $x\neq 1$ : le préciser ; (ii) pour $f(1)$, bien choisir la branche $x\leq 1$ (le point $1$ appartient à la deuxième branche, pas à la première) ; (iii) la continuité au point de raccordement se conclut par la triple égalité limite à gauche = limite à droite = $f(1)$, en citant le théorème de la limite bilatère.
\end{corrige}

\begin{popfigure}
\begin{tikzpicture}[scale=0.7, domain=-1:5]
\draw[popdark, ->] (-0.5,0) -- (5,0) node[right] {$x$ (dizaines de km)};
\draw[popdark, ->] (0,-0.5) -- (0,6) node[above] {$y$ (kFCFA)};
% Line y = x+1
\draw[poporange, very thick] (-0.5,0.5) -- (5,6) node[above right] {$y=x+1$};
% Asymptote y = x
\draw[popblue, dashed, thick] (0,0) -- (5,5) node[above] {$y=x$ (asymptote oblique)};
% Point at x=1
\fill[poporange] (1,2) circle (2.5pt) node[below left] {$(1,2)$};
% Distance representation
\draw[popmuted, <->] (2,3) -- (2,2) node[midway, right] {$1$};
\node[popmuted, anchor=west] at (2.2,2.5) {$f(x)-x=1$};
\end{tikzpicture}
\legende{Courbe de $f$ pour $a=1$ : droite $y=x+1$, asymptote oblique $y=x$, distance verticale constante $=1$.}
\end{popfigure}

\newpage

\coursec{Fiche bilan}

\coursec{Limites et continuité}

\courssub{Introduction et intuition}
\begin{savaistu}[D'où vient l'idée de limite ?]
Historiquement, le concept de limite est né de la nécessité de calculer des aires, des volumes et des vitesses instantanées (calcul différentiel). Newton et Leibniz l'ont manipulé intuitivement (\og quantités évanouissantes \fg). La définition rigoureuse ($\varepsilon-\delta$) est due à Cauchy et Weierstrass au XIX\textsuperscript{e} siècle. Au Cameroun, on l'utilise par exemple pour modéliser la charge maximale d'un pont (limite de résistance), la concentration limite d'un polluant dans la rivière Wouri, ou l'audience maximale d'une antenne MTN/Orange avant saturation.
\end{savaistu}

\begin{experience}[Approche intuitive : la flèche et la cible]
Imagine un archer (toi) qui tire une flèche vers une cible placée en $x_0$. La position de la flèche est $f(x)$.
\begin{itemize}
  \item Si la flèche se plante pile au centre ($\ell$) quand tu tires de plus en plus près ($x\to x_0$), on dit que $\lim_{x\to x_0}f(x)=\ell$.
  \item Si la flèche part vers le ciel ($+\infty$) ou s'enfonce dans le sol ($-\infty$), la limite est infinie.
  \item Si la flèche oscille sans se stabiliser, la limite \textbf{n'existe pas}.
\end{itemize}
Cette image guide toute la leçon : \textbf{la limite décrit le comportement de $f(x)$ quand $x$ \emph{s'approche} d'une valeur, sans nécessairement l'atteindre.}
\end{experience}

\coursec{Limite d'une fonction en un point}

\courssub{Définition et vocabulaire}
\begin{definition}[Limite finie en un point]
Soit $f$ une fonction définie sur un intervalle $I$ (sauf peut-être en $x_0$) et $x_0$ une valeur d'adhérence de $I$. Soit $\ell\in\mathbb{R}$.
On dit que \textbf{$f$ admet la limite $\ell$ en $x_0$} si, quand $x$ se rapproche arbitrairement de $x_0$ (sans y être égal), $f(x)$ se rapproche arbitrairement de $\ell$.
\par Notation : $\displaystyle\lim_{x\to x_0}f(x)=\ell$.
\end{definition}

\begin{definition}[Limites à droite et à gauche]
\begin{itemize}
  \item \cle{Limite à droite} en $x_0$ : $x$ tend vers $x_0$ par des valeurs \textbf{strictement supérieures} à $x_0$. Notation : $\displaystyle\lim_{x\to x_0^+}f(x)=\ell_+$.
  \item \cle{Limite à gauche} en $x_0$ : $x$ tend vers $x_0$ par des valeurs \textbf{strictement inférieures} à $x_0$. Notation : $\displaystyle\lim_{x\to x_0^-}f(x)=\ell_-$.
\end{itemize}
\end{definition}

\begin{propriete}[Théorème de l'unicité et liaison bilatérale]
\begin{enumerate}
  \item Si la limite existe, elle est \textbf{unique}.
  \item $\displaystyle\lim_{x\to x_0}f(x)=\ell \quad\Longleftrightarrow\quad \lim_{x\to x_0^-}f(x)=\lim_{x\to x_0^+}f(x)=\ell$.
\end{enumerate}
\emph{Au Bac : on doit savoir énoncer cette équivalence et l'utiliser pour montrer qu'une limite n'existe pas (limites latérales différentes).}
\end{propriete}

\courssub{Limites infinies en un point (asymptotes verticales)}
\begin{definition}
On note $\displaystyle\lim_{x\to x_0}f(x)=+\infty$ si $f(x)$ devient arbitrairement grand (positif) quand $x\to x_0$.
De même pour $-\infty$. Si les limites latérales sont infinies de signes contraires, on parle de \cle{pôle} (ex: $1/x$ en $0$).
\end{definition}

\begin{exemplebox}[Exemple : $f(x)=\dfrac{1}{x-2}$ en $x_0=2$]
$f$ n'est pas définie en $2$.
\begin{align*}
\lim_{x\to 2^+}\frac{1}{x-2} &= +\infty \quad (\text{dénominateur } \to 0^+) \\
\lim_{x\to 2^-}\frac{1}{x-2} &= -\infty \quad (\text{dénominateur } \to 0^-) \\
\Rightarrow \lim_{x\to 2}f(x) &\text{ n'existe pas (pas de limite finie, ni infinie unique).}
\end{align*}
La droite $x=2$ est une \cle{asymptote verticale}.
\end{exemplebox}

\begin{attention}[Piège classique : $x\to x_0$ vs $x=x_0$]
La limite en $x_0$ \textbf{ne dépend pas} de la valeur $f(x_0)$ (qui peut ne pas exister).
\par Exemple : $f(x)=\dfrac{x^2-4}{x-2}$. $f(2)$ n'existe pas, mais $\lim_{x\to 2}f(x)=\lim_{x\to 2}(x+2)=4$.
\par Ne confonds pas \og $f(x_0)$ \fg et \og $\lim_{x\to x_0}f(x)$ \fg.
\end{attention}

\coursec{Limite d'une fonction en l'infini}

\courssub{Définition et asymptotes horizontales}
\begin{definition}[Limite en $+\infty$ et $-\infty$]
\begin{itemize}
  \item $\displaystyle\lim_{x\to +\infty}f(x)=\ell$ : quand $x$ devient arbitrairement grand, $f(x)$ se rapproche de $\ell$.
  \item $\displaystyle\lim_{x\to -\infty}f(x)=\ell$ : quand $x$ devient arbitrairement grand en valeur absolue et négatif, $f(x)$ se rapproche de $\ell$.
\end{itemize}
Si $\ell\in\mathbb{R}$, la droite $y=\ell$ est une \cle{asymptote horizontale}.
\end{definition}

\begin{propriete}[Limites usuelles à connaître par cœur]
\begin{center}
\begin{tabularx}{\linewidth}{|l|X|X|}
\hline
\rowcolor{popblueL} \textbf{Fonction $f(x)$} & \textbf{$\boldsymbol{\lim_{x\to +\infty}}$} & \textbf{$\boldsymbol{\lim_{x\to -\infty}}$} \\
\hline
$x^n\ (n\in\mathbb{N}^*)$ & $+\infty$ & $+\infty$ si $n$ pair, $-\infty$ si $n$ impair \\
\hline
$\dfrac{1}{x}$ & $0^+$ & $0^-$ \\
\hline
$\sqrt{x}$ & $+\infty$ & \textit{Non défini} (domaine $[0,+\infty[$) \\
\hline
$|x|$ & $+\infty$ & $+\infty$ \\
\hline
\end{tabularx}
\end{center}
\par\medskip
\begin{tabularx}{\linewidth}{|l|X|X|}
\hline
\rowcolor{popblueL} \textbf{Fonction $f(x)$} & \textbf{$\boldsymbol{\lim_{x\to 0^+}}$} & \textbf{$\boldsymbol{\lim_{x\to 0^-}}$} \\
\hline
$\dfrac{1}{x}$ & $+\infty$ & $-\infty$ \\
\hline
$\sqrt{x}$ & $0$ & \textit{Non défini} \\
\hline
$|x|$ & $0$ & $0$ \\
\hline
$x^n\ (n\ge 1)$ & $0$ & $0$ \\
\hline
\end{tabularx}
\end{propriete}

\begin{savaistu}[Contexte camerounais : Asymptote et saturation]
Modélisons le débit internet (en Mbps) d'une box Orange/Mtn selon le nombre $x$ d'utilisateurs connectés simultanément : $D(x)=\dfrac{100x}{x+10}$.
$\lim_{x\to +\infty}D(x)=100$. La bande passante maximale (asymptote horizontale $y=100$) est de 100 Mbps. Quel que soit le nombre d'abonnés, le débit ne dépassera jamais 100 Mbps.
\end{savaistu}

\coursec{Opérations sur les limites}

\courssub{Règles de calcul (limites finies)}
\begin{propriete}[Opérations sur les limites finies]
Soit $f,g$ deux fonctions admettant des limites finies $\ell$ et $\ell'$ en $x_0$ (ou en $\pm\infty$).
\begin{align*}
\lim (f+g) &= \ell+\ell' \\
\lim (f-g) &= \ell-\ell' \\
\lim (f\times g) &= \ell\times\ell' \\
\lim \left(\frac{f}{g}\right) &= \frac{\ell}{\ell'} \quad\text{si } \ell'\neq 0
\end{align*}
\end{propriete}

\courssub{Règles avec l'infini (calcul formel)}
\begin{propriete}[Opérations avec $\pm\infty$ — Règles de signes]
On utilise les conventions suivantes (valables pour $x\to x_0$ ou $x\to\pm\infty$) :
\begin{center}
\begin{tabular}{|c|c|c|c|}
\hline
\rowcolor{popblueL} \textbf{Opération} & \textbf{Résultat} & \textbf{Condition} & \textbf{Indétermination} \\
\hline
$+\infty + (+\infty)$ & $+\infty$ & & \\
$-\infty + (-\infty)$ & $-\infty$ & & \\
$+\infty + (-\infty)$ & \textbf{?} & & \textbf{OUI ($+\infty-\infty$)} \\
\hline
$(\pm\infty)\times(\pm\infty)$ & $+\infty$ & Même signe & \\
$(\pm\infty)\times(\mp\infty)$ & $-\infty$ & Signes contraires & \\
$0 \times (\pm\infty)$ & \textbf{?} & & \textbf{OUI ($0\times\infty$)} \\
\hline
$\dfrac{\ell}{\pm\infty}\ (\ell\in\mathbb{R})$ & $0$ & & \\
$\dfrac{\pm\infty}{\ell}\ (\ell\in\mathbb{R}^*)$ & $\pm\infty$ & Règle des signes & \\
$\dfrac{0}{0}$ ou $\dfrac{\pm\infty}{\pm\infty}$ & \textbf{?} & & \textbf{OUI} \\
\hline
\end{tabular}
\end{center}
\end{propriete}

\begin{attention}[Les 4 formes indéterminées majeures]
Ne jamais écrire \og $=\ ?$ \fg ou \og $=$ indéterminé \fg comme résultat final. Une forme indéterminée signifie : \og \textbf{la règle ne s'applique pas, il faut transformer l'expression} \fg.
\begin{enumerate}
  \item $+\infty - \infty$ (somme)
  \item $0 \times \infty$ (produit)
  \item $\dfrac{0}{0}$ (quotient)
  \item $\dfrac{\infty}{\infty}$ (quotient)
\end{enumerate}
\end{attention}

\courssub{Limite d'une fonction composée}
\begin{propriete}[Théorème de la limite d'une composée]
Soit $f$ et $g$ deux fonctions. Si :
\begin{enumerate}
  \item $\displaystyle\lim_{x\to a}f(x)=b$ (fini ou infini),
  \item $\displaystyle\lim_{y\to b}g(y)=c$ (fini ou infini),
  \item \textbf{Condition cruciale :} si $b$ est fini, $g$ est \textbf{continue en $b$} (ou $f(x)\neq b$ dans un voisinage de $a$).
\end{enumerate}
Alors $\displaystyle\lim_{x\to a}g(f(x))=c$.
\par\medskip
\emph{Au Bac : on cite souvent \og limite de la composée = composée des limites \fg en vérifiant la continuité de la fonction extérieure au point limite de la fonction intérieure.}
\end{propriete}

\begin{exemplebox}[Composée : $\lim_{x\to +\infty}\exp(-x^2)$]
$f(x)=-x^2$, $g(y)=e^y$.
$\lim_{x\to +\infty}f(x)=-\infty$.
$\lim_{y\to -\infty}g(y)=0$.
$g$ est continue sur $\mathbb{R}$.
$\Rightarrow \lim_{x\to +\infty}e^{-x^2}=0$.
\end{exemplebox}

\coursec{Levier une forme indéterminée}

\courssub{Méthode 1 : Factorisation par le terme de plus haut degré (en $\pm\infty$)}
\begin{methode}[Limite de fraction rationnelle en $\pm\infty$]
Pour $f(x)=\dfrac{P(x)}{Q(x)}$ où $P,Q$ sont des polynômes.
\begin{enumerate}
  \item Identifier le terme de plus haut degré au numérateur ($a x^n$) et au dénominateur ($b x^m$).
  \item Factoriser $x^n$ au numérateur et $x^m$ au dénominateur.
  \item Simplifier par $x^{\min(n,m)}$.
  \item Calculer la limite des facteurs restants (qui tendent vers des constantes).
\end{enumerate}
\par\medskip
\textbf{Règle rapide (termes dominants) :} $\dfrac{P(x)}{Q(x)} \sim_{+\infty} \dfrac{a x^n}{b x^m} = \dfrac{a}{b}x^{n-m}$.
\begin{itemize}
  \item $n<m \Rightarrow$ limite $0$.
  \item $n=m \Rightarrow$ limite $\dfrac{a}{b}$.
  \item $n>m \Rightarrow$ limite $\pm\infty$ (signe selon parité de $n-m$ et signes de $a,b$).
\end{itemize}
\end{methode}

\begin{exoresolu}[Factorisation en $+\infty$]
Étudier $\lim_{x\to +\infty}\dfrac{3x^2-5x+2}{2x^2+7x-1}$.
\par\medskip
\textbf{Solution :}
Forme $\dfrac{+\infty}{+\infty}$.
Factorisons par $x^2$ (plus haut degré) :
\[
\frac{3x^2-5x+2}{2x^2+7x-1} = \frac{x^2(3-\frac{5}{x}+\frac{2}{x^2})}{x^2(2+\frac{7}{x}-\frac{1}{x^2})} = \frac{3-\frac{5}{x}+\frac{2}{x^2}}{2+\frac{7}{x}-\frac{1}{x^2}}
\]
$\lim_{x\to +\infty}\frac{5}{x}=\lim_{x\to +\infty}\frac{2}{x^2}=\lim_{x\to +\infty}\frac{7}{x}=\lim_{x\to +\infty}\frac{1}{x^2}=0$.
Donc la limite vaut $\dfrac{3-0+0}{2+0-0}=\dfrac{3}{2}$.
\end{exoresolu}

\courssub{Méthode 2 : Factorisation par $(x-x_0)$ (forme $\frac{0}{0}$ en $x_0$)}
\begin{methode}[Limite en $x_0$ donnant $\frac{0}{0}$]
Si $\lim_{x\to x_0}P(x)=\lim_{x\to x_0}Q(x)=0$ pour des polynômes $P,Q$.
\begin{enumerate}
  \item $x_0$ est racine de $P$ et $Q$ : $(x-x_0)$ divise $P$ et $Q$.
  \item Factoriser $P(x)=(x-x_0)P_1(x)$ et $Q(x)=(x-x_0)Q_1(x)$.
  \item Simplifier par $(x-x_0)$ (valide car $x\neq x_0$ pour la limite).
  \item Calculer $\lim_{x\to x_0}\frac{P_1(x)}{Q_1(x)}$.
\end{enumerate}
\end{methode}

\begin{exoresolu}[Factorisation en $x_0$]
$\lim_{x\to 3}\dfrac{x^2-9}{x^2-4x+3}$.
\par\medskip
\textbf{Solution :}
$P(3)=0, Q(3)=0 \Rightarrow$ forme $\frac{0}{0}$.
$P(x)=(x-3)(x+3)$.
$Q(x)=(x-3)(x-1)$.
Pour $x\neq 3$ : $\dfrac{(x-3)(x+3)}{(x-3)(x-1)} = \dfrac{x+3}{x-1}$.
Limite $= \dfrac{3+3}{3-1} = \dfrac{6}{2}=3$.
\end{exoresolu}

\courssub{Méthode 3 : Quantité conjuguée (racines carrées)}
\begin{methode}[Expression avec $\sqrt{\cdot}$ donnant $0/0$ ou $\infty-\infty$]
Utiliser l'identité : $\sqrt{A}-\sqrt{B} = \dfrac{A-B}{\sqrt{A}+\sqrt{B}}$.
\begin{enumerate}
  \item Identifier la différence de racines.
  \item Multiplier numérateur et dénominateur par la quantité conjuguée ($\sqrt{A}+\sqrt{B}$).
  \item Simplifier l'expression obtenue (souvent apparition d'un facteur $(x-x_0)$).
  \item Calculer la limite.
\end{enumerate}
\end{methode}

\begin{exoresolu}[Quantité conjuguée]
$\lim_{x\to 4}\dfrac{\sqrt{x}-2}{x-4}$.
\par\medskip
\textbf{Solution :}
Forme $\frac{0}{0}$.
Multiplier par $\sqrt{x}+2$ :
\[
\frac{\sqrt{x}-2}{x-4} \times \frac{\sqrt{x}+2}{\sqrt{x}+2} = \frac{x-4}{(x-4)(\sqrt{x}+2)} = \frac{1}{\sqrt{x}+2} \quad (x\neq 4)
\]
Limite $= \dfrac{1}{\sqrt{4}+2} = \dfrac{1}{4}$.
\end{exoresolu}

\begin{attention}[Erreurs fréquentes sur les conjuguées]
\begin{itemize}
  \item Oublier de multiplier \textbf{numérateur ET dénominateur}.
  \item Ne pas simplifier le facteur $(x-x_0)$ après développement.
  \item Confondre $\sqrt{A}+\sqrt{B}$ (conjuguée de $\sqrt{A}-\sqrt{B}$) et $\sqrt{A}-\sqrt{B}$.
\end{itemize}
\end{attention}

\coursec{Continuité en un point et sur un intervalle}

\courssub{Définition de la continuité en un point}
\begin{definition}[Continuité en $x_0$]
Une fonction $f$ est \cle{continue en $x_0$} si et seulement si les trois conditions suivantes sont réunies :
\begin{enumerate}
  \item $f$ est \textbf{définie} en $x_0$ ($x_0\in D_f$).
  \item $f$ admet une \textbf{limite finie} en $x_0$ : $\lim_{x\to x_0}f(x)=\ell\in\mathbb{R}$.
  \item \textbf{Égalité} : $\lim_{x\to x_0}f(x)=f(x_0)$.
\end{enumerate}
Si l'une des conditions manque, $f$ est \cle{discontinue} en $x_0$.
\end{definition}

\begin{propriete}[Types de discontinuité en $x_0$]
\begin{itemize}
  \item \textbf{Discontinuité de première espèce (saut) :} Limites latérales finies mais différentes ($\ell_-\neq\ell_+$). Ex: partie entière $[x]$ en un entier.
  \item \textbf{Discontinuité de deuxième espèce (infinie) :} Au moins une limite latérale est infinie (asymptote verticale). Ex: $1/x$ en $0$.
  \item \textbf{Discontinuité évitable (trou) :} Limite finie $\ell$ existe mais $f(x_0)$ n'existe pas ou $f(x_0)\neq\ell$. On peut \og boucher le trou \fg en redéfinissant $f(x_0)=\ell$.
\end{itemize}
\end{propriete}

\courssub{Continuité sur un intervalle}
\begin{definition}
$f$ est \cle{continue sur un intervalle $I$} si elle est continue en \textbf{tout point} de $I$.
\par Pour un intervalle fermé $[a,b]$, on demande en plus la continuité à droite en $a$ et à gauche en $b$.
\end{definition}

\begin{propriete}[Fonctions usuelles continues sur leur domaine]
Les fonctions suivantes sont continues sur tout intervalle inclus dans leur ensemble de définition :
\begin{itemize}
  \item Fonctions \textbf{polynômes} (sur $\mathbb{R}$).
  \item Fonctions \textbf{rationnelles} (sur $\mathbb{R}\setminus\{\text{zéros du dénominateur}\}$).
  \item $x\mapsto \sqrt{x}$ (sur $[0,+\infty[$).
  \item $x\mapsto |x|$ (sur $\mathbb{R}$).
  \item $x\mapsto \exp(x), \ln(x)$ (sur $\mathbb{R}$ et $]0,+\infty[$).
  \item $x\mapsto \sin x, \cos x$ (sur $\mathbb{R}$).
\end{itemize}
\end{propriete}

\begin{propriete}[Stabilité de la continuité par opérations]
Si $f$ et $g$ sont continues en $x_0$ (ou sur $I$) :
\begin{itemize}
  \item $f+g$, $f-g$, $f\times g$ sont continues en $x_0$ (sur $I$).
  \item $\dfrac{f}{g}$ est continue en $x_0$ \textbf{si $g(x_0)\neq 0$} (sur $I$ où $g\neq 0$).
  \item Si $f$ continue en $x_0$ et $g$ continue en $f(x_0)$, alors $g\circ f$ est continue en $x_0$.
\end{itemize}
\end{propriete}

\begin{exemplebox}[Étude de continuité : $f(x)=\dfrac{x^2-1}{x-1}$ sur $\mathbb{R}$]
$D_f=\mathbb{R}\setminus\{1\}$.
\begin{itemize}
  \item Sur $]-\infty,1[$ et $]1,+\infty[$ : $f$ est rationnelle, dénominateur non nul $\Rightarrow$ continue.
  \item En $1$ : $f$ n'est pas définie $\Rightarrow$ \textbf{discontinue} (discontinuité évitable car $\lim_{x\to 1}f(x)=2$).
\end{itemize}
\end{exemplebox}

\courssub{Théorème des valeurs intermédiaires (TVI) — Exigible au Bac}
\begin{propriete}[Théorème des valeurs intermédiaires]
Soit $f$ continue sur un intervalle $I=[a,b]$ ($a<b$).
Pour tout réel $k$ compris entre $f(a)$ et $f(b)$, il existe \textbf{au moins un} $c\in[a,b]$ tel que $f(c)=k$.
\par\medskip
\emph{Corollaire (Recherche de zéro) :} Si $f(a)\times f(b)<0$, il existe au moins un $c\in[a,b]$ tel que $f(c)=0$.
\par Si de plus $f$ est \textbf{strictement monotone} sur $[a,b]$, ce $c$ est \textbf{unique}.
\end{propriete}

\begin{methode}[Résolution d'équation $f(x)=k$ par TVI]
\begin{enumerate}
  \item Vérifier que $f$ est continue sur $[a,b]$.
  \item Calculer $f(a)$ et $f(b)$.
  \item Vérifier que $k\in[f(a),f(b)]$ (ou $f(a)f(b)<0$ pour $k=0$).
  \item Conclure : \og Il existe au moins une solution $c\in[a,b]$ \fg.
  \item (Si monotonie prouvée) \og La solution est unique \fg.
\end{enumerate}
\end{methode}

\coursec{Interprétations graphiques et synthèse}

\courssub{Lecture graphique des limites et continuité}
\begin{popfigure}
\begin{tikzpicture}[scale=0.8, domain=-3:3]
\draw[->, thick] (-3.5,0) -- (3.5,0) node[right] {$x$};
\draw[->, thick] (0,-2.5) -- (0,3.5) node[above] {$y$};
% Asymptote verticale x=1
\draw[dashed, popred, thick] (1,-2.5) -- (1,3.5) node[above right] {$x=1$};
% Asymptote horizontale y=2
\draw[dashed, popblue, thick] (-3.5,2) -- (3.5,2) node[right] {$y=2$};
% Courbe gauche x<1 : y = 2 + 1/(x-1) -> -inf at 1-
\draw[popcurve, thick, domain=-3:0.95, samples=100] plot (\x, {2 + 1/(\x-1)});
% Courbe droite x>1 : y = 2 + 1/(x-1) -> +inf at 1+
\draw[popcurve, thick, domain=1.05:3, samples=100] plot (\x, {2 + 1/(\x-1)});
% Point trou à (0,1) -> f(0)=1? Let's put a hole at (0,1) and point at (0,3)
\fill[popgreen] (0,3) circle (2pt) node[above left] {$f(0)=3$};
\draw[popgreen, thick] (0,1) circle (4pt) node[below right] {Trou $\lim=1$};
% Point discontinuité saut à x=-2
\fill[poporange] (-2,1) circle (2pt);
\draw[poporange, thick] (-2,-1) circle (4pt) node[below] {Saut};
\end{tikzpicture}
\legende{Graphique récapitulatif : asymptotes verticales/horizontales, trou (discontinuité évitable), saut (discontinuité 1\textsuperscript{ère} espèce).}
\end{popfigure}

\begin{aretenir}[Synthèse visuelle : Ce qu'il faut voir sur un graphique]
\begin{itemize}
  \item \textbf{Limite finie $\ell$ en $x_0$} : La courbe se rapproche du point $(x_0,\ell)$ des deux côtés.
  \item \textbf{Limite infinie} : La courbe s'éloigne verticalement vers le haut/bas près de $x_0$ (asymptote verticale).
  \item \textbf{Limite en $\pm\infty$} : La courbe se rapproche d'une horizontale $y=\ell$ (asymptote horizontale) ou part vers $\pm\infty$.
  \item \textbf{Continuité en $x_0$} : Le point $(x_0,f(x_0))$ est \textbf{sur la courbe} et la courbe ne \og casse \fg pas (on trace sans lever le crayon).
  \item \textbf{Discontinuité évitable} : Un \textbf{trou} (cercle vide) à $(x_0,\ell)$ ; éventuellement un point isolé ailleurs.
  \item \textbf{Discontinuité de saut} : Deux \textbf{extrémités distinctes} de courbe au-dessus de $x_0$ (cercles pleins ou vides selon définition).
  \item \textbf{Discontinuité infinie} : Les deux branches partent en sens opposés vers $\pm\infty$ de part et d'autre de l'asymptote verticale.
\end{itemize}
\end{aretenir}

\courssub{Exercices types Bac (avec corrections)}

\begin{exercice}[Exercice 1 : Limites en un point et en l'infini]
Soit $f(x)=\dfrac{2x^2-3x+1}{x-2}$.
\begin{enumerate}
  \item Déterminer $D_f$.
  \item Calculer $\lim_{x\to 2^-}f(x)$ et $\lim_{x\to 2^+}f(x)$. En déduire l'existence de la limite en $2$.
  \item Calculer $\lim_{x\to +\infty}f(x)$ et $\lim_{x\to -\infty}f(x)$. Y a-t-il une asymptote horizontale ? Oblique ?
  \item Étudier la continuité de $f$ sur $D_f$.
\end{enumerate}
\end{exercice}

\begin{corrige}[Exercice 1]
\begin{enumerate}
  \item $D_f = \mathbb{R}\setminus\{2\}$.
  \item Numérateur en $2$ : $2(4)-6+1=3\neq 0$. Dénominateur $\to 0$.
  $x\to 2^- \Rightarrow x-2\to 0^- \Rightarrow f(x)\to \frac{3}{0^-}=-\infty$.
  $x\to 2^+ \Rightarrow x-2\to 0^+ \Rightarrow f(x)\to \frac{3}{0^+}=+\infty$.
  Limites latérales infinies de signes contraires $\Rightarrow$ \textbf{pas de limite en $2$}. Asymptote verticale $x=2$.
  \item Degrés : num 2, dén 1 $\Rightarrow$ limite $\pm\infty$ selon signe coefficient directeur.
  $\dfrac{2x^2}{x}=2x$. $\lim_{+\infty}=+\infty$, $\lim_{-\infty}=-\infty$. Pas d'asymptote horizontale.
  Division euclidienne : $(2x^2-3x+1) = (x-2)(2x+1) + 3$.
  $f(x)=2x+1+\frac{3}{x-2}$. Asymptote oblique : $y=2x+1$.
  \item $f$ est rationnelle, continue sur $D_f$ (dénominateur non nul sur $D_f$).
\end{enumerate}
\end{corrige}

\begin{exercice}[Exercice 2 : Formes indéterminées et quantité conjuguée]
Calculer les limites suivantes :
\begin{enumerate}
  \item $\lim_{x\to 1}\dfrac{x^3-1}{x^2-1}$
  \item $\lim_{x\to +\infty}\left(\sqrt{x^2+3x}-x\right)$
  \item $\lim_{x\to 0}\dfrac{\sqrt{x+4}-2}{x}$
\end{enumerate}
\end{exercice}

\begin{corrige}[Exercice 2]
\begin{enumerate}
  \item Forme $\frac{0}{0}$ en $1$.
  $x^3-1=(x-1)(x^2+x+1)$, $x^2-1=(x-1)(x+1)$.
  $\lim_{x\to 1}\dfrac{x^2+x+1}{x+1}=\dfrac{3}{2}$.
  \item Forme $+\infty-\infty$.
  $\sqrt{x^2+3x}-x = \dfrac{(x^2+3x)-x^2}{\sqrt{x^2+3x}+x} = \dfrac{3x}{x(\sqrt{1+\frac{3}{x}}+1)} = \dfrac{3}{\sqrt{1+\frac{3}{x}}+1}$.
  Limite $= \dfrac{3}{1+1}=\dfrac{3}{2}$.
  \item Forme $\frac{0}{0}$.
  Multiplier par $\sqrt{x+4}+2$ :
  $\dfrac{x+4-4}{x(\sqrt{x+4}+2)} = \dfrac{x}{x(\sqrt{x+4}+2)} = \dfrac{1}{\sqrt{x+4}+2}$.
  Limite $= \dfrac{1}{2+2}=\dfrac{1}{4}$.
\end{enumerate}
\end{corrige}

\begin{exercice}[Exercice 3 : Continuité et TVI]
Soit $f(x)=x^3-3x+1$ sur $[-2,2]$.
\begin{enumerate}
  \item Montrer que $f$ est continue sur $[-2,2]$.
  \item Calculer $f(-2)$ et $f(2)$. En déduire, d'après le TVI, que l'équation $f(x)=0$ admet au moins une solution dans $[-2,2]$.
  \item Étudier les variations de $f$ sur $[-2,2]$ pour montrer que cette solution est unique.
\end{enumerate}
\end{exercice}

\begin{corrige}[Exercice 3]
\begin{enumerate}
  \item $f$ est polynôme $\Rightarrow$ continue sur $\mathbb{R}$, donc sur $[-2,2]$.
  \item $f(-2)=-8+6+1=-1$. $f(2)=8-6+1=3$.
  $f(-2)\times f(2) = -3 < 0$.
  Par TVI (corollaire), $\exists c\in[-2,2]$ tel que $f(c)=0$.
  \item $f'(x)=3x^2-3=3(x-1)(x+1)$.
  $f'$ s'annule en $-1$ et $1$.
  Tableau de variations :
  \begin{center}
  \begin{tabular}{|c|ccccccc|}
  \hline
  $x$ & $-2$ & & $-1$ & & $1$ & & $2$ \\
  \hline
  $f'$ & & $+$ & $0$ & $-$ & $0$ & $+$ & \\
  \hline
  $f$ & $-1$ & $\nearrow$ & $3$ & $\searrow$ & $-1$ & $\nearrow$ & $3$ \\
  \hline
  \end{tabular}
  \end{center}
  $f$ est strictement croissante sur $[-2,-1]$, décroissante sur $[-1,1]$, croissante sur $[1,2]$.
  Sur $[-2,-1]$ : $f(-2)=-1<0$, $f(-1)=3>0$ $\Rightarrow$ une unique solution $c_1\in[-2,-1]$.
  Sur $[-1,1]$ : $f(-1)=3>0$, $f(1)=-1<0$ $\Rightarrow$ une unique solution $c_2\in[-1,1]$.
  Sur $[1,2]$ : $f(1)=-1<0$, $f(2)=3>0$ $\Rightarrow$ une unique solution $c_3\in[1,2]$.
  \textbf{Il y a trois solutions distinctes.} (L'énoncé disait \og une solution \fg, c'était un piège : TVI donne \og au moins une \fg, la monotonie par morceaux donne trois).
\end{enumerate}
\end{corrige}

\coursec{Bilan final et checklist Bac}

\begin{popfigure}
\begin{center}\tcbox[colback=popblue,colframe=popblue,coltext=white,arc=9pt,boxsep=4pt,left=6pt,right=6pt]{\bfseries\small Limites et continuité}\end{center}
\vspace{-2pt}
\begin{tcbraster}[raster columns=3,raster equal height=rows,raster column skip=6pt,raster row skip=6pt,enhanced,blankest]
\begin{tcolorbox}[enhanced,colback=poporange,colframe=poporange,coltext=white,boxrule=0.9pt,arc=3pt,left=4pt,right=4pt,top=3pt,bottom=3pt,boxsep=1pt,halign=left]\bfseries\scriptsize Limite en $x_0$ : finie, $\pm\infty$, latérales $\ell_-,\ell_+$\end{tcolorbox}
\begin{tcolorbox}[enhanced,colback=popgreen,colframe=popgreen,coltext=white,boxrule=0.9pt,arc=3pt,left=4pt,right=4pt,top=3pt,bottom=3pt,boxsep=1pt,halign=left]\bfseries\scriptsize Limite en $\pm\infty$ : asymptotes H/O\end{tcolorbox}
\begin{tcolorbox}[enhanced,colback=poppurple,colframe=poppurple,coltext=white,boxrule=0.9pt,arc=3pt,left=4pt,right=4pt,top=3pt,bottom=3pt,boxsep=1pt,halign=left]\bfseries\scriptsize Limites usuelles : $x^n, 1/x, \sqrt{x}, |x|, \exp, \ln$\end{tcolorbox}
\begin{tcolorbox}[enhanced,colback=poppink,colframe=poppink,coltext=white,boxrule=0.9pt,arc=3pt,left=4pt,right=4pt,top=3pt,bottom=3pt,boxsep=1pt,halign=left]\bfseries\scriptsize Opérations : $\Sigma,\Pi,/,\circ$ + 4 indéterminées\end{tcolorbox}
\begin{tcolorbox}[enhanced,colback=popgold,colframe=popgold,coltext=white,boxrule=0.9pt,arc=3pt,left=4pt,right=4pt,top=3pt,bottom=3pt,boxsep=1pt,halign=left]\bfseries\scriptsize Méthodes : Facto deg max, Facto $(x-x_0)$, Conjuguée\end{tcolorbox}
\begin{tcolorbox}[enhanced,colback=popdark,colframe=popdark,coltext=white,boxrule=0.9pt,arc=3pt,left=4pt,right=4pt,top=3pt,bottom=3pt,boxsep=1pt,halign=left]\bfseries\scriptsize Continuité : 3 cond., TVI, fcts usuelles\end{tcolorbox}
\end{tcbraster}

\legende{Carte mentale récapitulative : les 6 piliers de la leçon.}
\end{popfigure}

\begin{aretenir}[L'essentiel à retenir pour l'épreuve]
\textbf{1. Définitions incontournables}
\begin{itemize}
  \item $\lim_{x\to x_0}f(x)=\ell \iff \forall \varepsilon>0, \exists \eta>0, |x-x_0|<\eta \Rightarrow |f(x)-\ell|<\varepsilon$ (savoir l'écrire en français).
  \item $\lim_{x\to x_0}f(x)$ existe $\iff \lim_{x\to x_0^-}f(x)=\lim_{x\to x_0^+}f(x)$.
  \item $f$ continue en $x_0 \iff \begin{cases} x_0\in D_f \\ \lim_{x\to x_0}f(x) \text{ existe et est finie} \\ \lim_{x\to x_0}f(x)=f(x_0) \end{cases}$
\end{itemize}

\textbf{2. Limites usuelles (tableau de référence)}
$+\infty$ l'emporte sur tout polynôme. $\exp$ l'emporte sur toute puissance. $\ln$ est négligeable devant toute puissance.
\begin{align*}
\lim_{x\to 0}\frac{\sin x}{x}=1,\quad \lim_{x\to 0}\frac{\tan x}{x}=1,\quad \lim_{x\to 0}\frac{e^x-1}{x}=1,\quad \lim_{x\to 0}\frac{\ln(1+x)}{x}=1
\end{align*}
(Souvent utiles pour lever des $0/0$ sans développement limité formel).

\textbf{3. Opérations : Tableau de vérité pour l'infini}
\begin{center}
\begin{tabular}{|c|c|}
\hline
\rowcolor{popblueL} \textbf{Forme} & \textbf{Résultat} \\
\hline
$+\infty+\infty$, $-\infty-\infty$ & $\pm\infty$ \\
$+\infty-\infty$ & \textbf{Indéterminée} \\
$(\pm\infty)\times(\pm\infty)$ & $+\infty$ \\
$(\pm\infty)\times(\mp\infty)$ & $-\infty$ \\
$0\times(\pm\infty)$ & \textbf{Indéterminée} \\
$c/(\pm\infty)\ (c\in\mathbb{R})$ & $0$ \\
$(\pm\infty)/c\ (c\neq 0)$ & $\pm\infty$ \\
$0/0$, $\infty/\infty$ & \textbf{Indéterminées} \\
\hline
\end{tabular}
\end{center}

\textbf{4. Boîte à outils pour lever les indéterminées}
\begin{enumerate}
  \item \textbf{Polynôme/Fraction rationnelle en $\pm\infty$} : Factoriser par le terme de plus haut degré (ou garder termes dominants).
  \item \textbf{Fraction rationnelle en $x_0$ ($0/0$)} : Factoriser par $(x-x_0)$ (identités remarquables : $a^n-b^n$, $a^3\pm b^3$, etc.).
  \item \textbf{Racines ($\sqrt{A}\pm\sqrt{B}$)} : Quantité conjuguée $\frac{A-B}{\sqrt{A}\mp\sqrt{B}}$.
  \item \textbf{Composée} : Faire le changement de variable $y=f(x)$, calculer $\lim f$, puis $\lim g$. Vérifier continuité de $g$ si limite de $f$ finie.
\end{enumerate}

\textbf{5. Continuité \& TVI}
\begin{itemize}
  \item Polynômes, rationnelles (domaine), $\sqrt{x}$, $|x|$, $\exp$, $\ln$, $\sin$, $\cos$ : continues sur leur domaine.
  \item Somme, produit, quotient (dénominateur $\neq 0$), composée de continues = continue.
  \item TVI : $f$ continue sur $[a,b]$, $k\in[f(a),f(b)] \Rightarrow \exists c\in[a,b], f(c)=k$.
  \item Unicité si $f$ strictement monotone sur $[a,b]$.
\end{itemize}
\end{aretenir}

\begin{aretenir}[Checklist finale : Suis-je prêt pour le Bac ?]
Coche chaque case seulement si tu es capable de le faire \textbf{sans regarder le cours}.
\begin{itemize}
  \item[$\square$] Écrire la définition $\varepsilon-\eta$ (ou version française) de $\lim_{x\to x_0}f(x)=\ell$ et de la continuité en $x_0$.
  \item[$\square$] Calculer des limites latérales pour décider de l'existence d'une limite bilatérale.
  \item[$\square$] Réciter les limites de $x^n$, $1/x$, $\sqrt{x}$, $|x|$, $\exp$, $\ln$ en $0$, $\pm\infty$ sans hésiter.
  \item[$\square$] Appliquer les règles d'opérations (somme, produit, quotient) avec la règle des signes sur $\pm\infty$.
  \item[$\square$] Identifier instantanément les 4 formes indéterminées : $+\infty-\infty$, $0\times\infty$, $0/0$, $\infty/\infty$.
  \item[$\square$] Lever une indétermination en $\pm\infty$ par factorisation du terme dominant (méthode rapide et méthode rigoureuse).
  \item[$\square$] Lever un $0/0$ en $x_0$ par factorisation par $(x-x_0)$ (identités remarquables maîtrisées).
  \item[$\square$] Utiliser la quantité conjuguée pour $\sqrt{A}-\sqrt{B}$ (numérateur ou dénominateur).
  \item[$\square$] Calculer une limite de composée $g\circ f$ en 2 étapes (vérifier condition de continuité si limite intérieure finie).
  \item[$\square$] Vérifier les 3 conditions de continuité en un point (existence image, existence limite, égalité).
  \item[$\square$] Classer une discontinuité : évitable (trou), saut (1\textsuperscript{ère} espèce), infinie (2\textsuperscript{de} espèce / asymptote verticale).
  \item[$\square$] Justifier la continuité d'une fonction usuelle sur un intervalle de son domaine.
  \item[$\square$] Appliquer le TVI (existence) et le TVI + monotonie (unicité) pour résoudre $f(x)=k$.
  \item[$\square$] Lire sur un graphique : limite finie/infinie, asymptote H/V/O, continuité, types de discontinuité.
\end{itemize}
\end{aretenir}

\begin{savaistu}[Le saviez-vous ? Le pont de Wouri et la continuité]
Le pont de Wouri (Douala) supporte une charge maximale. La fonction \og contrainte dans l'acier \fg en fonction de la charge est continue... jusqu'à la rupture ! En mathématiques, une discontinuité brutale (saut) modéliserait une rupture soudaine. Les ingénieurs utilisent la continuité et les limites pour garantir que les structures ne présentent pas de \og sauts \fg de résistance. La dérivée (prochaine leçon) leur permettra de calculer la \emph{raideur} (pente de la courbe contrainte-déformation).
\end{savaistu}

\findecours

\end{document}
